% layout-v3-repair: keep adjective section opening together
% layout-v2-repairs: non-author v1 findings; masters unchanged
% Complete reviewed Unit 15; editorial master: chapters/ch15/student.md
% Markdown SHA-256: 73bbaac1b2d7de6c5a1c88fa3b546d226e7b2abb9dd9b0249483133e19882dda
% One-use initial assembly followed by complete content and page review.
\newcommand{\runningtitle}{Unit 15. Negation and quantity}
% Layout inherited from the reviewed previous course; set a running title per unit.
% This is a small hand-maintained TeX layout, not a conversion pipeline.
\RequirePackage{fix-cm}
\documentclass[11pt,twoside]{article}
\usepackage{fontspec}
\setmainfont{Cambria}[
  Renderer=HarfBuzz,
  BoldFont=cambriab.ttf,
  ItalicFont=cambriai.ttf,
  BoldItalicFont=cambriaz.ttf
]
\setmonofont{lmmono10-regular.otf}[BoldFont=lmmonolt10-bold.otf]
\usepackage[paperwidth=165mm,paperheight=241mm,inner=21mm,outer=19mm,
  top=22mm,bottom=20mm,headheight=14pt,headsep=7mm,footskip=10mm]{geometry}
\usepackage{xcolor,array,booktabs,longtable,tabularx}
\usepackage{enumitem,fancyhdr}
\usepackage[unicode,hidelinks]{hyperref}
\definecolor{ink}{HTML}{000000}
\definecolor{muted}{HTML}{4C5558}
\setlength{\parindent}{0pt}
\setlength{\parskip}{3.5pt}
\AtBeginDocument{\fontsize{11}{14.3}\selectfont}
\setlength{\tabcolsep}{4pt}
\renewcommand{\arraystretch}{1.10}
\setlist[itemize]{leftmargin=1.3em,label=\textbullet,itemsep=2pt,topsep=3pt,parsep=0pt}
\setlist[enumerate]{leftmargin=1.8em,itemsep=3pt,topsep=4pt,parsep=0pt}
\widowpenalty=10000
\clubpenalty=10000
\emergencystretch=1em
\raggedbottom
\pagestyle{fancy}
\fancyhf{}
\providecommand{\runningtitle}{Old English}
\fancyhead[LE,RO]{\small\itshape \runningtitle}
\fancyhead[LO,RE]{\small Old English}
\fancyfoot[LE,RO]{\small\thepage}
\renewcommand{\headrulewidth}{0.3pt}
\newcommand{\booktitle}[2]{
  {\fontsize{9}{11}\selectfont\color{muted}#1\par}
  \vspace{3pt}
  {\fontsize{20}{24}\selectfont\color{ink}#2\par}
  \vspace{6pt}
}
\newcommand{\source}[1]{\par{\footnotesize\color{muted}#1\par}}
\newcommand{\readingtext}[1]{
  \par\addvspace{5pt}
  \begingroup\leftskip=1.5em
    \fontsize{12}{16}\selectfont #1\par
  \endgroup\addvspace{5pt}
}
\makeatletter
\renewcommand{\section}{\@startsection{section}{1}{\z@}
  {-2.4ex \@plus -0.6ex \@minus -0.2ex}{0.8ex \@plus 0.2ex}
  {\normalfont\fontsize{12}{15}\selectfont\bfseries\color{ink}}}
\renewcommand{\subsection}{\@startsection{subsection}{2}{\z@}
  {-1.8ex \@plus -0.5ex \@minus -0.2ex}{0.5ex \@plus 0.2ex}
  {\normalfont\fontsize{11}{14}\selectfont\bfseries\color{ink}}}
\makeatother

\hyphenpenalty=10000
\exhyphenpenalty=10000
\emergencystretch=3em
\newcommand{\itempara}[2]{\par\addvspace{3pt}\noindent\hangindent=1.8em\hangafter=1\makebox[1.8em][l]{#1.}#2\par}
\newcommand{\reserve}[1]{\par\penalty-100\begingroup\dimen0=\pagegoal\advance\dimen0 by-\pagetotal\ifdim\dimen0<#1\newpage\fi\endgroup}
\begin{document}
\addfontfeatures{Ligatures=TeXOff}
\raggedright
\booktitle{UNIT 15}{Negation, quantity and comparison in discourse}

% mdparagraph:P001:begin
A negative clause can contain more than one negative word without becoming affirmative. Quantity words distinguish an individual from a group, while comparison can express equality, preference or a proportional change. This unit brings these distinctions into connected arguments and requests.
% mdparagraph:P001:end

% mdparagraph:P002:begin
Sections 1–5 give the vocabulary, complete defined systems and constructions needed by the practice. The eight translated examples and the lettered exercises are new teaching material. Cook IX.28–30 and X, XI.31 retain their complete original prompts. The three readings preserve Cook’s Old English and original notes; modern helps are identified separately.
% mdparagraph:P002:end

% mdparagraph:P003:begin
\textbf{Scope and notation.} Learn the forms for use in the meanings and constructions stated. A reference table supplies forms for recognition, without requiring their independent recall. Supplied expressions remain available during their exercises and readings. “Unassigned” means that no form or use in that position is required here; it does not assert that a historical form was impossible. In verbal tables, 1, 2 and 3 denote person. Preserve vowel length and the distinction between e and ę.
% mdparagraph:P003:end

\reserve{7\baselineskip}
% mdparagraph:P004:begin
\section*{1. Nouns and their meanings}
% mdparagraph:P004:end

% mdparagraph:P005:begin
The tables give nominative, genitive, dative and accusative in both numbers. Unless a number is marked as reference or unassigned, both numbers may be used. Endings that coincide still have distinct grammatical functions: the sentence determines the case. The lexical meaning, gender and stem belong to the whole word.\textsuperscript{1}
% mdparagraph:P005:end

\reserve{7\baselineskip}
% mdparagraph:P006:begin
\subsection*{Strong masculine nouns}
% mdparagraph:P006:end

% mdparagraph:P007:begin
The nominative and accusative singular have no ending. The other singular endings are genitive -es and dative -e. The plural endings are nominative/\allowbreak{}accusative -as, genitive -a and dative -um.
% mdparagraph:P007:end

% teaching: NEW-U15-VOC-HLAFORD
\reserve{11\baselineskip}
% mdparagraph:P008:begin
\textbf{hlāford} — masculine; lord; master. Stem \textbf{hlāford-}.
% mdparagraph:P008:end

% teaching: NEW-U15-NOM-HLAFORD
% table: NEW-U15-NOM-HLAFORD-DISPLAY
% mdtable:T01:begin
\par\begingroup
\fontsize{10.5}{13.2}\selectfont
\setlength{\LTleft}{0pt}\setlength{\LTright}{\fill}
\noindent
\begin{tabular}{@{}>{\raggedright\arraybackslash}p{39.46mm}>{\raggedright\arraybackslash}p{39.46mm}>{\raggedright\arraybackslash}p{39.46mm}@{}}
\toprule
\textbf{Case} & \textbf{Singular} & \textbf{Plural} \\
\midrule
Nominative & hlāford & hlāfordas \\
Genitive & hlāfordes & hlāforda \\
Dative & hlāforde & hlāfordum \\
Accusative & hlāford & hlāfordas \\
\bottomrule
\end{tabular}
\par\endgroup

% mdtable:T01:end
% teaching: NEW-U15-VOC-WINGEARD
\reserve{11\baselineskip}
% mdparagraph:P009:begin
\textbf{wīngeard} — masculine; vineyard. Stem \textbf{wīngeard-}.
% mdparagraph:P009:end

% teaching: NEW-U15-NOM-WINGEARD
% table: NEW-U15-NOM-WINGEARD-DISPLAY
% mdtable:T02:begin
\par\begingroup
\fontsize{10.5}{13.2}\selectfont
\setlength{\LTleft}{0pt}\setlength{\LTright}{\fill}
\noindent
\begin{tabular}{@{}>{\raggedright\arraybackslash}p{39.46mm}>{\raggedright\arraybackslash}p{39.46mm}>{\raggedright\arraybackslash}p{39.46mm}@{}}
\toprule
\textbf{Case} & \textbf{Singular} & \textbf{Plural} \\
\midrule
Nominative & wīngeard & wīngeardas \\
Genitive & wīngeardes & wīngearda \\
Dative & wīngearde & wīngeardum \\
Accusative & wīngeard & wīngeardas \\
\bottomrule
\end{tabular}
\par\endgroup

% mdtable:T02:end
% teaching: NEW-U15-VOC-WER
\reserve{11\baselineskip}
% mdparagraph:P010:begin
\textbf{wer} — masculine; man. Stem \textbf{wer-}.
% mdparagraph:P010:end

% teaching: NEW-U15-NOM-WER
% table: NEW-U15-NOM-WER-DISPLAY
% mdtable:T03:begin
\par\begingroup
\fontsize{10.5}{13.2}\selectfont
\setlength{\LTleft}{0pt}\setlength{\LTright}{\fill}
\noindent
\begin{tabular}{@{}>{\raggedright\arraybackslash}p{39.46mm}>{\raggedright\arraybackslash}p{39.46mm}>{\raggedright\arraybackslash}p{39.46mm}@{}}
\toprule
\textbf{Case} & \textbf{Singular} & \textbf{Plural} \\
\midrule
Nominative & wer & weras \\
Genitive & weres & wera \\
Dative & were & werum \\
Accusative & wer & weras \\
\bottomrule
\end{tabular}
\par\endgroup

% mdtable:T03:end
% teaching: NEW-U15-VOC-LEORNUNGCNIHT
\reserve{11\baselineskip}
% mdparagraph:P011:begin
\textbf{leornungcniht} — masculine; disciple. Stem \textbf{leornungcniht-}.
% mdparagraph:P011:end

% teaching: NEW-U15-NOM-LEORNUNGCNIHT
% table: NEW-U15-NOM-LEORNUNGCNIHT-DISPLAY
% mdtable:T04:begin
\par\begingroup
\fontsize{10.5}{13.2}\selectfont
\setlength{\LTleft}{0pt}\setlength{\LTright}{\fill}
\noindent
\begin{tabular}{@{}>{\raggedright\arraybackslash}p{29.59mm}>{\raggedright\arraybackslash}p{44.39mm}>{\raggedright\arraybackslash}p{44.39mm}@{}}
\toprule
\textbf{Case} & \textbf{Singular} & \textbf{Plural} \\
\midrule
Nominative & leornungcniht & leornungcnihtas \\
Genitive & leornungcnihtes & leornungcnihta \\
Dative & leornungcnihte & leornungcnihtum \\
Accusative & leornungcniht & leornungcnihtas \\
\bottomrule
\end{tabular}
\par\endgroup

% mdtable:T04:end
% teaching: NEW-U15-VOC-LAREOW
\reserve{11\baselineskip}
% mdparagraph:P012:begin
\textbf{lārēow} — masculine; master; teacher. Stem \textbf{lārēow-}.
% mdparagraph:P012:end

% teaching: NEW-U15-NOM-LAREOW
% table: NEW-U15-NOM-LAREOW-DISPLAY
% mdtable:T05:begin
\par\begingroup
\fontsize{10.5}{13.2}\selectfont
\setlength{\LTleft}{0pt}\setlength{\LTright}{\fill}
\noindent
\begin{tabular}{@{}>{\raggedright\arraybackslash}p{39.46mm}>{\raggedright\arraybackslash}p{39.46mm}>{\raggedright\arraybackslash}p{39.46mm}@{}}
\toprule
\textbf{Case} & \textbf{Singular} & \textbf{Plural} \\
\midrule
Nominative & lārēow & lārēowas \\
Genitive & lārēowes & lārēowa \\
Dative & lārēowe & lārēowum \\
Accusative & lārēow & lārēowas \\
\bottomrule
\end{tabular}
\par\endgroup

% mdtable:T05:end
% teaching: NEW-U15-VOC-HAM-NOUN
\reserve{11\baselineskip}
% mdparagraph:P013:begin
\textbf{hām} — masculine; home; dwelling. Stem \textbf{hām-}. This is the noun “home; dwelling”, distinct from the directional adverb hām.
% mdparagraph:P013:end

% teaching: NEW-U15-NOM-HAM-NOUN
% table: NEW-U15-NOM-HAM-NOUN-DISPLAY
% mdtable:T06:begin
\par\begingroup
\fontsize{10.5}{13.2}\selectfont
\setlength{\LTleft}{0pt}\setlength{\LTright}{\fill}
\noindent
\begin{tabular}{@{}>{\raggedright\arraybackslash}p{39.46mm}>{\raggedright\arraybackslash}p{39.46mm}>{\raggedright\arraybackslash}p{39.46mm}@{}}
\toprule
\textbf{Case} & \textbf{Singular} & \textbf{Plural} \\
\midrule
Nominative & hām & hāmas \\
Genitive & hāmes & hāma \\
Dative & hāme & hāmum \\
Accusative & hām & hāmas \\
\bottomrule
\end{tabular}
\par\endgroup

% mdtable:T06:end
% mdparagraph:P014:begin
Fōtsceamul retains the vowel of its final element -sceamul- before endings. The long vowel in the first element fōt- does not require the loss of that later vowel.
% mdparagraph:P014:end

% teaching: NEW-U15-VOC-FOTSCEAMUL
\reserve{11\baselineskip}
% mdparagraph:P015:begin
\textbf{fōtsceamul} — masculine; footstool. Stem \textbf{fōtsceamul-}.
% mdparagraph:P015:end

% teaching: NEW-U15-NOM-FOTSCEAMUL
% table: NEW-U15-NOM-FOTSCEAMUL-DISPLAY
% mdtable:T07:begin
\par\begingroup
\fontsize{10.5}{13.2}\selectfont
\setlength{\LTleft}{0pt}\setlength{\LTright}{\fill}
\noindent
\begin{tabular}{@{}>{\raggedright\arraybackslash}p{34.82mm}>{\raggedright\arraybackslash}p{41.78mm}>{\raggedright\arraybackslash}p{41.78mm}@{}}
\toprule
\textbf{Case} & \textbf{Singular} & \textbf{Plural} \\
\midrule
Nominative & fōtsceamul & fōtsceamulas \\
Genitive & fōtsceamules & fōtsceamula \\
Dative & fōtsceamule & fōtsceamulum \\
Accusative & fōtsceamul & fōtsceamulas \\
\bottomrule
\end{tabular}
\par\endgroup

% mdtable:T07:end
\reserve{7\baselineskip}
% mdparagraph:P016:begin
\subsection*{Strong neuter nouns}
% mdparagraph:P016:end

% mdparagraph:P017:begin
Heavy stems have no ending in nominative and accusative, singular or plural. Genitive and dative singular take -es and -e; genitive and dative plural take -a and -um. Bebod has a short final stem and takes nominative/\allowbreak{}accusative plural -u.
% mdparagraph:P017:end

% teaching: NEW-U15-VOC-HAER
\reserve{11\baselineskip}
% mdparagraph:P018:begin
\textbf{hǣr} — neuter; hair; a hair. Stem \textbf{hǣr-}.
% mdparagraph:P018:end

% teaching: NEW-U15-NOM-HAER
% table: NEW-U15-NOM-HAER-DISPLAY
% mdtable:T08:begin
\par\begingroup
\fontsize{10.5}{13.2}\selectfont
\setlength{\LTleft}{0pt}\setlength{\LTright}{\fill}
\noindent
\begin{tabular}{@{}>{\raggedright\arraybackslash}p{39.46mm}>{\raggedright\arraybackslash}p{39.46mm}>{\raggedright\arraybackslash}p{39.46mm}@{}}
\toprule
\textbf{Case} & \textbf{Singular} & \textbf{Plural} \\
\midrule
Nominative & hǣr & hǣr \\
Genitive & hǣres & hǣra \\
Dative & hǣre & hǣrum \\
Accusative & hǣr & hǣr \\
\bottomrule
\end{tabular}
\par\endgroup

% mdtable:T08:end
% teaching: NEW-U15-VOC-LEOTH
\reserve{11\baselineskip}
% mdparagraph:P019:begin
\textbf{lēoð} — neuter; poem; poetry; verse. Stem \textbf{lēoð-}.
% mdparagraph:P019:end

% teaching: NEW-U15-NOM-LEOTH
% table: NEW-U15-NOM-LEOTH-DISPLAY
% mdtable:T09:begin
\par\begingroup
\fontsize{10.5}{13.2}\selectfont
\setlength{\LTleft}{0pt}\setlength{\LTright}{\fill}
\noindent
\begin{tabular}{@{}>{\raggedright\arraybackslash}p{39.46mm}>{\raggedright\arraybackslash}p{39.46mm}>{\raggedright\arraybackslash}p{39.46mm}@{}}
\toprule
\textbf{Case} & \textbf{Singular} & \textbf{Plural} \\
\midrule
Nominative & lēoð & lēoð \\
Genitive & lēoðes & lēoða \\
Dative & lēoðe & lēoðum \\
Accusative & lēoð & lēoð \\
\bottomrule
\end{tabular}
\par\endgroup

% mdtable:T09:end
% teaching: NEW-U15-VOC-BEBOD
\reserve{11\baselineskip}
% mdparagraph:P020:begin
\textbf{bebod} — neuter; commandment. Stem \textbf{bebod-}.
% mdparagraph:P020:end

% teaching: NEW-U15-NOM-BEBOD
% table: NEW-U15-NOM-BEBOD-DISPLAY
% mdtable:T10:begin
\par\begingroup
\fontsize{10.5}{13.2}\selectfont
\setlength{\LTleft}{0pt}\setlength{\LTright}{\fill}
\noindent
\begin{tabular}{@{}>{\raggedright\arraybackslash}p{39.46mm}>{\raggedright\arraybackslash}p{39.46mm}>{\raggedright\arraybackslash}p{39.46mm}@{}}
\toprule
\textbf{Case} & \textbf{Singular} & \textbf{Plural} \\
\midrule
Nominative & bebod & bebodu \\
Genitive & bebodes & beboda \\
Dative & bebode & bebodum \\
Accusative & bebod & bebodu \\
\bottomrule
\end{tabular}
\par\endgroup

% mdtable:T10:end
% mdparagraph:P021:begin
Hēafod loses the unstressed o before an ending: the ending-bearing stem is hēafd-. Its selected nominative/\allowbreak{}accusative plural is hēafdu. Compare fōtsceamul above: medial vowels are not all deleted by one general rule.
% mdparagraph:P021:end

% teaching: NEW-U15-VOC-HEAFOD
\reserve{11\baselineskip}
% mdparagraph:P022:begin
\textbf{hēafod} — neuter; head. Stem \textbf{hēafd-}.
% mdparagraph:P022:end

% teaching: NEW-U15-NOM-HEAFOD
% table: NEW-U15-NOM-HEAFOD-DISPLAY
% mdtable:T11:begin
\par\begingroup
\fontsize{10.5}{13.2}\selectfont
\setlength{\LTleft}{0pt}\setlength{\LTright}{\fill}
\noindent
\begin{tabular}{@{}>{\raggedright\arraybackslash}p{39.46mm}>{\raggedright\arraybackslash}p{39.46mm}>{\raggedright\arraybackslash}p{39.46mm}@{}}
\toprule
\textbf{Case} & \textbf{Singular} & \textbf{Plural} \\
\midrule
Nominative & hēafod & hēafdu \\
Genitive & hēafdes & hēafda \\
Dative & hēafde & hēafdum \\
Accusative & hēafod & hēafdu \\
\bottomrule
\end{tabular}
\par\endgroup

% mdtable:T11:end
% teaching: NEW-U15-VOC-FLAESC
\reserve{12\baselineskip}
% mdparagraph:P023:begin
\textbf{flǣsc} — neuter; flesh. Stem \textbf{flǣsc-}. \textbf{Plural reference:} the plural forms show the morphology; use the singular for the assigned meaning “flesh”.
% mdparagraph:P023:end

% teaching: NEW-U15-NOM-FLAESC
% table: NEW-U15-NOM-FLAESC-DISPLAY
% mdtable:T12:begin
\par\begingroup
\fontsize{10.5}{13.2}\selectfont
\setlength{\LTleft}{0pt}\setlength{\LTright}{\fill}
\noindent
\begin{tabular}{@{}>{\raggedright\arraybackslash}p{39.46mm}>{\raggedright\arraybackslash}p{39.46mm}>{\raggedright\arraybackslash}p{39.46mm}@{}}
\toprule
\textbf{Case} & \textbf{Singular} & \textbf{Plural} \\
\midrule
Nominative & flǣsc & flǣsc \\
Genitive & flǣsces & flǣsca \\
Dative & flǣsce & flǣscum \\
Accusative & flǣsc & flǣsc \\
\bottomrule
\end{tabular}
\par\endgroup

% mdtable:T12:end
\reserve{7\baselineskip}
% mdparagraph:P024:begin
\subsection*{Strong feminine nouns}
% mdparagraph:P024:end

% mdparagraph:P025:begin
These long-stem nouns have no nominative singular ending. The other singular cases take -e. The plural has nominative/\allowbreak{}accusative and genitive -a, dative -um. Nouns in -nes double s before an ending. The -ung and -ing nouns use the same selected endings; oblique singular forms in -a are recognized alternatives.
% mdparagraph:P025:end

% teaching: NEW-U15-VOC-BRYD
\reserve{11\baselineskip}
% mdparagraph:P026:begin
\textbf{brȳd} — feminine; bride. Stem \textbf{brȳd-}.
% mdparagraph:P026:end

% teaching: NEW-U15-NOM-BRYD
% table: NEW-U15-NOM-BRYD-DISPLAY
% mdtable:T13:begin
\par\begingroup
\fontsize{10.5}{13.2}\selectfont
\setlength{\LTleft}{0pt}\setlength{\LTright}{\fill}
\noindent
\begin{tabular}{@{}>{\raggedright\arraybackslash}p{39.46mm}>{\raggedright\arraybackslash}p{39.46mm}>{\raggedright\arraybackslash}p{39.46mm}@{}}
\toprule
\textbf{Case} & \textbf{Singular} & \textbf{Plural} \\
\midrule
Nominative & brȳd & brȳda \\
Genitive & brȳde & brȳda \\
Dative & brȳde & brȳdum \\
Accusative & brȳde & brȳda \\
\bottomrule
\end{tabular}
\par\endgroup

% mdtable:T13:end
% teaching: NEW-U15-VOC-HEFIGNES
\reserve{11\baselineskip}
% mdparagraph:P027:begin
\textbf{hęfignes} — feminine; burden. Stem \textbf{hęfigness-}.
% mdparagraph:P027:end

% teaching: NEW-U15-NOM-HEFIGNES
% table: NEW-U15-NOM-HEFIGNES-DISPLAY
% mdtable:T14:begin
\par\begingroup
\fontsize{10.5}{13.2}\selectfont
\setlength{\LTleft}{0pt}\setlength{\LTright}{\fill}
\noindent
\begin{tabular}{@{}>{\raggedright\arraybackslash}p{38.19mm}>{\raggedright\arraybackslash}p{38.19mm}>{\raggedright\arraybackslash}p{42.0mm}@{}}
\toprule
\textbf{Case} & \textbf{Singular} & \textbf{Plural} \\
\midrule
Nominative & hęfignes & hęfignessa \\
Genitive & hęfignesse & hęfignessa \\
Dative & hęfignesse & hęfignessum \\
Accusative & hęfignesse & hęfignessa \\
\bottomrule
\end{tabular}
\par\endgroup

% mdtable:T14:end
% teaching: NEW-U15-VOC-AWENDEDNES
\reserve{11\baselineskip}
% mdparagraph:P028:begin
\textbf{āwęndednes} — feminine; translation. Stem \textbf{āwęndedness-}.
% mdparagraph:P028:end

% teaching: NEW-U15-NOM-AWENDEDNES
% table: NEW-U15-NOM-AWENDEDNES-DISPLAY
% mdtable:T15:begin
\par\begingroup
\fontsize{10.5}{13.2}\selectfont
\setlength{\LTleft}{0pt}\setlength{\LTright}{\fill}
\noindent
\begin{tabular}{@{}>{\raggedright\arraybackslash}p{33.82mm}>{\raggedright\arraybackslash}p{40.59mm}>{\raggedright\arraybackslash}p{43.97mm}@{}}
\toprule
\textbf{Case} & \textbf{Singular} & \textbf{Plural} \\
\midrule
Nominative & āwęndednes & āwęndednessa \\
Genitive & āwęndednesse & āwęndednessa \\
Dative & āwęndednesse & āwęndednessum \\
Accusative & āwęndednesse & āwęndednessa \\
\bottomrule
\end{tabular}
\par\endgroup

% mdtable:T15:end
% teaching: NEW-U15-VOC-COSTNUNG
\reserve{11\baselineskip}
% mdparagraph:P029:begin
\textbf{costnung} — feminine; temptation. Stem \textbf{costnung-}.
% mdparagraph:P029:end

% teaching: NEW-U15-NOM-COSTNUNG
% table: NEW-U15-NOM-COSTNUNG-DISPLAY
% mdtable:T16:begin
\par\begingroup
\fontsize{10.5}{13.2}\selectfont
\setlength{\LTleft}{0pt}\setlength{\LTright}{\fill}
\noindent
\begin{tabular}{@{}>{\raggedright\arraybackslash}p{39.46mm}>{\raggedright\arraybackslash}p{39.46mm}>{\raggedright\arraybackslash}p{39.46mm}@{}}
\toprule
\textbf{Case} & \textbf{Singular} & \textbf{Plural} \\
\midrule
Nominative & costnung & costnunga \\
Genitive & costnunge & costnunga \\
Dative & costnunge & costnungum \\
Accusative & costnunge & costnunga \\
\bottomrule
\end{tabular}
\par\endgroup

% mdtable:T16:end
% teaching: NEW-U15-VOC-BECRAEDING
\reserve{12\baselineskip}
% mdparagraph:P030:begin
\textbf{bēcrǣding} — feminine; reading; reading books. Stem \textbf{bēcrǣding-}. \textbf{Plural reference:} the meaning practised here is singular “reading”.
% mdparagraph:P030:end

% teaching: NEW-U15-NOM-BECRAEDING
% table: NEW-U15-NOM-BECRAEDING-DISPLAY
% mdtable:T17:begin
\par\begingroup
\fontsize{10.5}{13.2}\selectfont
\setlength{\LTleft}{0pt}\setlength{\LTright}{\fill}
\noindent
\begin{tabular}{@{}>{\raggedright\arraybackslash}p{38.19mm}>{\raggedright\arraybackslash}p{38.19mm}>{\raggedright\arraybackslash}p{42.0mm}@{}}
\toprule
\textbf{Case} & \textbf{Singular} & \textbf{Plural} \\
\midrule
Nominative & bēcrǣding & bēcrǣdinga \\
Genitive & bēcrǣdinge & bēcrǣdinga \\
Dative & bēcrǣdinge & bēcrǣdingum \\
Accusative & bēcrǣdinge & bēcrǣdinga \\
\bottomrule
\end{tabular}
\par\endgroup

% mdtable:T17:end
\reserve{7\baselineskip}
% mdparagraph:P031:begin
\subsection*{Weak masculine nouns}
% mdparagraph:P031:end

% mdparagraph:P032:begin
The nominative singular ends in -a. The other singular cases and nominative/\allowbreak{}accusative plural end in -an; genitive plural takes -ena and dative plural -um. In lēoðwyrhta the final element is the noun head, and the ending follows the whole stem.
% mdparagraph:P032:end

% teaching: NEW-U15-VOC-GEREFA
\reserve{11\baselineskip}
% mdparagraph:P033:begin
\textbf{gerēfa} — masculine; steward; reeve. Stem \textbf{gerēf-}.
% mdparagraph:P033:end

% teaching: NEW-U15-NOM-GEREFA
% table: NEW-U15-NOM-GEREFA-DISPLAY
% mdtable:T18:begin
\par\begingroup
\fontsize{10.5}{13.2}\selectfont
\setlength{\LTleft}{0pt}\setlength{\LTright}{\fill}
\noindent
\begin{tabular}{@{}>{\raggedright\arraybackslash}p{39.46mm}>{\raggedright\arraybackslash}p{39.46mm}>{\raggedright\arraybackslash}p{39.46mm}@{}}
\toprule
\textbf{Case} & \textbf{Singular} & \textbf{Plural} \\
\midrule
Nominative & gerēfa & gerēfan \\
Genitive & gerēfan & gerēfena \\
Dative & gerēfan & gerēfum \\
Accusative & gerēfan & gerēfan \\
\bottomrule
\end{tabular}
\par\endgroup

% mdtable:T18:end
% teaching: NEW-U15-VOC-LEOTHWYRHTA
\reserve{11\baselineskip}
% mdparagraph:P034:begin
\textbf{lēoðwyrhta} — masculine; poet. Stem \textbf{lēoðwyrht-}.
% mdparagraph:P034:end

% teaching: NEW-U15-NOM-LEOTHWYRHTA
% table: NEW-U15-NOM-LEOTHWYRHTA-DISPLAY
% mdtable:T19:begin
\par\begingroup
\fontsize{10.5}{13.2}\selectfont
\setlength{\LTleft}{0pt}\setlength{\LTright}{\fill}
\noindent
\begin{tabular}{@{}>{\raggedright\arraybackslash}p{35.87mm}>{\raggedright\arraybackslash}p{39.46mm}>{\raggedright\arraybackslash}p{43.05mm}@{}}
\toprule
\textbf{Case} & \textbf{Singular} & \textbf{Plural} \\
\midrule
Nominative & lēoðwyrhta & lēoðwyrhtan \\
Genitive & lēoðwyrhtan & lēoðwyrhtena \\
Dative & lēoðwyrhtan & lēoðwyrhtum \\
Accusative & lēoðwyrhtan & lēoðwyrhtan \\
\bottomrule
\end{tabular}
\par\endgroup

% mdtable:T19:end
% teaching: NEW-U15-VOC-WILLA
\reserve{11\baselineskip}
% mdparagraph:P035:begin
\textbf{willa} — masculine; will; disposition. Stem \textbf{will-}. \textbf{Plural reference:} the assigned use is singular “will; disposition”.
% mdparagraph:P035:end

% teaching: NEW-U15-NOM-WILLA
% table: NEW-U15-NOM-WILLA-DISPLAY
% mdtable:T20:begin
\par\begingroup
\fontsize{10.5}{13.2}\selectfont
\setlength{\LTleft}{0pt}\setlength{\LTright}{\fill}
\noindent
\begin{tabular}{@{}>{\raggedright\arraybackslash}p{39.46mm}>{\raggedright\arraybackslash}p{39.46mm}>{\raggedright\arraybackslash}p{39.46mm}@{}}
\toprule
\textbf{Case} & \textbf{Singular} & \textbf{Plural} \\
\midrule
Nominative & willa & willan \\
Genitive & willan & willena \\
Dative & willan & willum \\
Accusative & willan & willan \\
\bottomrule
\end{tabular}
\par\endgroup

% mdtable:T20:end
\reserve{7\baselineskip}
% mdparagraph:P036:begin
\subsection*{Root and kinship nouns}
% mdparagraph:P036:end

% mdparagraph:P037:begin
Learn the vowel changes with their case positions. Bōc has bēc in genitive and dative singular and nominative and accusative plural. Tōð has tēð in dative singular and nominative and accusative plural. Recognize genitive singular bōce, dative singular tōðe and nominative/\allowbreak{}accusative plural tōðas as alternatives to the selected forms below.
% mdparagraph:P037:end

% teaching: NEW-U15-VOC-BOC
\reserve{11\baselineskip}
% mdparagraph:P038:begin
\textbf{bōc} — feminine; book. Stem \textbf{bōc/\allowbreak{}bēc-}.
% mdparagraph:P038:end

% teaching: NEW-U15-NOM-BOC
% table: NEW-U15-NOM-BOC-DISPLAY
% mdtable:T21:begin
\par\begingroup
\fontsize{10.5}{13.2}\selectfont
\setlength{\LTleft}{0pt}\setlength{\LTright}{\fill}
\noindent
\begin{tabular}{@{}>{\raggedright\arraybackslash}p{39.46mm}>{\raggedright\arraybackslash}p{39.46mm}>{\raggedright\arraybackslash}p{39.46mm}@{}}
\toprule
\textbf{Case} & \textbf{Singular} & \textbf{Plural} \\
\midrule
Nominative & bōc & bēc \\
Genitive & bēc & bōca \\
Dative & bēc & bōcum \\
Accusative & bōc & bēc \\
\bottomrule
\end{tabular}
\par\endgroup

% mdtable:T21:end
% teaching: NEW-U15-VOC-TOTH
\reserve{11\baselineskip}
% mdparagraph:P039:begin
\textbf{tōð} — masculine; tooth. Stem \textbf{tōð/\allowbreak{}tēð-}.
% mdparagraph:P039:end

% teaching: NEW-U15-NOM-TOTH
% table: NEW-U15-NOM-TOTH-DISPLAY
% mdtable:T22:begin
\par\begingroup
\fontsize{10.5}{13.2}\selectfont
\setlength{\LTleft}{0pt}\setlength{\LTright}{\fill}
\noindent
\begin{tabular}{@{}>{\raggedright\arraybackslash}p{39.46mm}>{\raggedright\arraybackslash}p{39.46mm}>{\raggedright\arraybackslash}p{39.46mm}@{}}
\toprule
\textbf{Case} & \textbf{Singular} & \textbf{Plural} \\
\midrule
Nominative & tōð & tēð \\
Genitive & tōðes & tōða \\
Dative & tēð & tōðum \\
Accusative & tōð & tēð \\
\bottomrule
\end{tabular}
\par\endgroup

% mdtable:T22:end
% mdparagraph:P040:begin
Sweostor is unchanged throughout the singular and in nominative/\allowbreak{}accusative plural. Only genitive plural sweostra and dative plural sweostrum have different surfaces. Its pattern must be learned separately from the earlier brōðor and mōdor.
% mdparagraph:P040:end

% teaching: NEW-U15-VOC-SWEOSTOR
\reserve{11\baselineskip}
% mdparagraph:P041:begin
\textbf{sweostor} — feminine; sister. Stem \textbf{sweostr-}.
% mdparagraph:P041:end

% teaching: NEW-U15-NOM-SWEOSTOR
% table: NEW-U15-NOM-SWEOSTOR-DISPLAY
% mdtable:T23:begin
\par\begingroup
\fontsize{10.5}{13.2}\selectfont
\setlength{\LTleft}{0pt}\setlength{\LTright}{\fill}
\noindent
\begin{tabular}{@{}>{\raggedright\arraybackslash}p{39.46mm}>{\raggedright\arraybackslash}p{39.46mm}>{\raggedright\arraybackslash}p{39.46mm}@{}}
\toprule
\textbf{Case} & \textbf{Singular} & \textbf{Plural} \\
\midrule
Nominative & sweostor & sweostor \\
Genitive & sweostor & sweostra \\
Dative & sweostor & sweostrum \\
Accusative & sweostor & sweostor \\
\bottomrule
\end{tabular}
\par\endgroup

% mdtable:T23:end
\reserve{7\baselineskip}
% mdparagraph:P042:begin
\subsection*{Abstract nouns in -u}
% mdparagraph:P042:end

% mdparagraph:P043:begin
Hǣtu uses hǣte in the selected oblique singular cases. Hǣlu is unchanged throughout the selected singular. Oblique variants in -u, -o and -e occur, and nominative -o occurs beside -u. Learn the selected singular forms; no plural count meaning is assigned to either word.\textsuperscript{2}
% mdparagraph:P043:end

% teaching: NEW-U15-VOC-HAETU
\reserve{11\baselineskip}
% mdparagraph:P044:begin
\textbf{hǣtu} — feminine; heat. Stem \textbf{hǣt-}. \textbf{Plural unassigned.}
% mdparagraph:P044:end

% teaching: NEW-U15-NOM-HAETU
% table: NEW-U15-NOM-HAETU-DISPLAY
% mdtable:T24:begin
\par\begingroup
\fontsize{10.5}{13.2}\selectfont
\setlength{\LTleft}{0pt}\setlength{\LTright}{\fill}
\noindent
\begin{tabular}{@{}>{\raggedright\arraybackslash}p{39.46mm}>{\raggedright\arraybackslash}p{39.46mm}>{\raggedright\arraybackslash}p{39.46mm}@{}}
\toprule
\textbf{Case} & \textbf{Singular} & \textbf{Plural} \\
\midrule
Nominative & hǣtu & Unassigned \\
Genitive & hǣte & Unassigned \\
Dative & hǣte & Unassigned \\
Accusative & hǣte & Unassigned \\
\bottomrule
\end{tabular}
\par\endgroup

% mdtable:T24:end
% teaching: NEW-U15-VOC-HAELU
\reserve{11\baselineskip}
% mdparagraph:P045:begin
\textbf{hǣlu} — feminine; salvation. Stem \textbf{hǣl-}. \textbf{Plural unassigned.}
% mdparagraph:P045:end

% teaching: NEW-U15-NOM-HAELU
% table: NEW-U15-NOM-HAELU-DISPLAY
% mdtable:T25:begin
\par\begingroup
\fontsize{10.5}{13.2}\selectfont
\setlength{\LTleft}{0pt}\setlength{\LTright}{\fill}
\noindent
\begin{tabular}{@{}>{\raggedright\arraybackslash}p{39.46mm}>{\raggedright\arraybackslash}p{39.46mm}>{\raggedright\arraybackslash}p{39.46mm}@{}}
\toprule
\textbf{Case} & \textbf{Singular} & \textbf{Plural} \\
\midrule
Nominative & hǣlu & Unassigned \\
Genitive & hǣlu & Unassigned \\
Dative & hǣlu & Unassigned \\
Accusative & hǣlu & Unassigned \\
\bottomrule
\end{tabular}
\par\endgroup

% mdtable:T25:end
% mdparagraph:P046:begin
The complete class model strengu follows. It is \textbf{for recognition only in both numbers}. Its spelling is retained from Wright. Besides the selected forms, recognize nominative singular strengo, oblique singular strengu/\allowbreak{}strengo, and nominative/\allowbreak{}accusative plural strenge/\allowbreak{}strengu/\allowbreak{}strengo. The availability of a class pattern does not supply a count meaning for every abstract noun.
% mdparagraph:P046:end

% teaching: NEW-U15-VOC-STRENGU
\reserve{11\baselineskip}
% mdparagraph:P047:begin
\textbf{strengu} — feminine; strength. Stem \textbf{streng-}. \textbf{Singular and plural reference only.}
% mdparagraph:P047:end

% teaching: NEW-U15-NOM-STRENGU
% table: NEW-U15-NOM-STRENGU-DISPLAY
% mdtable:T26:begin
\par\begingroup
\fontsize{10.5}{13.2}\selectfont
\setlength{\LTleft}{0pt}\setlength{\LTright}{\fill}
\noindent
\begin{tabular}{@{}>{\raggedright\arraybackslash}p{39.46mm}>{\raggedright\arraybackslash}p{39.46mm}>{\raggedright\arraybackslash}p{39.46mm}@{}}
\toprule
\textbf{Case} & \textbf{Singular} & \textbf{Plural} \\
\midrule
Nominative & strengu & strenga \\
Genitive & strenge & strenga \\
Dative & strenge & strengum \\
Accusative & strenge & strenga \\
\bottomrule
\end{tabular}
\par\endgroup

% mdtable:T26:end
\reserve{7\baselineskip}
% mdparagraph:P048:begin
\subsection*{People names and an agent noun}
% mdparagraph:P048:end

% mdparagraph:P049:begin
Engle and Dęne name peoples in the plural. They take nominative/\allowbreak{}accusative -e, genitive -a and dative -um. No singular is assigned. Dęniga and Dęnia are recognized genitive plural alternatives.
% mdparagraph:P049:end

% teaching: NEW-U15-VOC-ENGLE
\reserve{11\baselineskip}
% mdparagraph:P050:begin
\textbf{Engle} — masculine; the English (people). Stem \textbf{Engl-}. \textbf{Singular unassigned.}
% mdparagraph:P050:end

% teaching: NEW-U15-NOM-ENGLE
% table: NEW-U15-NOM-ENGLE-DISPLAY
% mdtable:T27:begin
\par\begingroup
\fontsize{10.5}{13.2}\selectfont
\setlength{\LTleft}{0pt}\setlength{\LTright}{\fill}
\noindent
\begin{tabular}{@{}>{\raggedright\arraybackslash}p{39.46mm}>{\raggedright\arraybackslash}p{39.46mm}>{\raggedright\arraybackslash}p{39.46mm}@{}}
\toprule
\textbf{Case} & \textbf{Singular} & \textbf{Plural} \\
\midrule
Nominative & Unassigned & Engle \\
Genitive & Unassigned & Engla \\
Dative & Unassigned & Englum \\
Accusative & Unassigned & Engle \\
\bottomrule
\end{tabular}
\par\endgroup

% mdtable:T27:end
% teaching: NEW-U15-VOC-DENE
\reserve{11\baselineskip}
% mdparagraph:P051:begin
\textbf{Dęne} — masculine; the Danes. Stem \textbf{Dęn-}. \textbf{Singular unassigned.}
% mdparagraph:P051:end

% teaching: NEW-U15-NOM-DENE
% table: NEW-U15-NOM-DENE-DISPLAY
% mdtable:T28:begin
\par\begingroup
\fontsize{10.5}{13.2}\selectfont
\setlength{\LTleft}{0pt}\setlength{\LTright}{\fill}
\noindent
\begin{tabular}{@{}>{\raggedright\arraybackslash}p{39.46mm}>{\raggedright\arraybackslash}p{39.46mm}>{\raggedright\arraybackslash}p{39.46mm}@{}}
\toprule
\textbf{Case} & \textbf{Singular} & \textbf{Plural} \\
\midrule
Nominative & Unassigned & Dęne \\
Genitive & Unassigned & Dęna \\
Dative & Unassigned & Dęnum \\
Accusative & Unassigned & Dęne \\
\bottomrule
\end{tabular}
\par\endgroup

% mdtable:T28:end
% mdparagraph:P052:begin
The agent noun ālīesend has singular -es and -e as above. Its selected plural has unchanged nominative/\allowbreak{}accusative, genitive -ra and dative -um. Nominative/\allowbreak{}accusative plural forms in -e and -as are recognized alternatives. The plural is morphological reference; the assigned title “Redeemer” is singular.
% mdparagraph:P052:end

% teaching: NEW-U15-VOC-ALIESEND
\reserve{11\baselineskip}
% mdparagraph:P053:begin
\textbf{ālīesend} — masculine; Redeemer; one who delivers. Stem \textbf{ālīesend-}. \textbf{Plural reference.}
% mdparagraph:P053:end

% teaching: NEW-U15-NOM-ALIESEND
% table: NEW-U15-NOM-ALIESEND-DISPLAY
% mdtable:T29:begin
\par\begingroup
\fontsize{10.5}{13.2}\selectfont
\setlength{\LTleft}{0pt}\setlength{\LTright}{\fill}
\noindent
\begin{tabular}{@{}>{\raggedright\arraybackslash}p{39.46mm}>{\raggedright\arraybackslash}p{39.46mm}>{\raggedright\arraybackslash}p{39.46mm}@{}}
\toprule
\textbf{Case} & \textbf{Singular} & \textbf{Plural} \\
\midrule
Nominative & ālīesend & ālīesend \\
Genitive & ālīesendes & ālīesendra \\
Dative & ālīesende & ālīesendum \\
Accusative & ālīesend & ālīesend \\
\bottomrule
\end{tabular}
\par\endgroup

% mdtable:T29:end
\reserve{7\baselineskip}
% mdparagraph:P054:begin
\subsection*{Additional senses of familiar nouns}
% mdparagraph:P054:end

% mdparagraph:P055:begin
Keep the earlier declensions. These meanings are needed in the following exercises.
% mdparagraph:P055:end

% teaching: NEW-U15-RECORD-MED-ACC-PL
% teaching: NEW-U15-RECORD-MED-ACC-SG
% teaching: NEW-U15-RECORD-MED-DAT-PL
% teaching: NEW-U15-RECORD-MED-DAT-SG
% teaching: NEW-U15-RECORD-MED-GEN-PL
% teaching: NEW-U15-RECORD-MED-GEN-SG
% teaching: NEW-U15-RECORD-MED-NOM-PL
% teaching: NEW-U15-RECORD-MED-NOM-SG
% teaching: NEW-U15-VOC-MED-SENSE
% mdparagraph:P056:begin
\textbf{mēd} — hire; wages paid for work.
% mdparagraph:P056:end

% teaching: NEW-U15-RECORD-TID-ACC-PL
% teaching: NEW-U15-RECORD-TID-ACC-SG
% teaching: NEW-U15-RECORD-TID-DAT-PL
% teaching: NEW-U15-RECORD-TID-DAT-SG
% teaching: NEW-U15-RECORD-TID-GEN-PL
% teaching: NEW-U15-RECORD-TID-GEN-SG
% teaching: NEW-U15-RECORD-TID-NOM-PL
% teaching: NEW-U15-RECORD-TID-NOM-SG
% teaching: NEW-U15-VOC-TID-SENSE
% mdparagraph:P057:begin
\textbf{tīd} — hour.
% mdparagraph:P057:end

% teaching: NEW-U15-RECORD-STRAET-ACC-PL
% teaching: NEW-U15-RECORD-STRAET-ACC-SG
% teaching: NEW-U15-RECORD-STRAET-DAT-PL
% teaching: NEW-U15-RECORD-STRAET-DAT-SG
% teaching: NEW-U15-RECORD-STRAET-GEN-PL
% teaching: NEW-U15-RECORD-STRAET-GEN-SG
% teaching: NEW-U15-RECORD-STRAET-NOM-PL
% teaching: NEW-U15-RECORD-STRAET-NOM-SG
% teaching: NEW-U15-VOC-STRAET-SENSE
% mdparagraph:P058:begin
\textbf{strǣt} — public place; street as a gathering place.
% mdparagraph:P058:end

% teaching: NEW-U15-VOC-WORD
% mdparagraph:P059:begin
\textbf{word} — speech; utterance.
% mdparagraph:P059:end

\reserve{7\baselineskip}
% mdparagraph:P060:begin
\section*{2. Negation}
% mdparagraph:P060:end

\reserve{7\baselineskip}
% mdparagraph:P061:begin
\subsection*{Finite negation and contractions}
% mdparagraph:P061:end

% teaching: NEW-U15-GRAM-U15-NEGATIVE-COPULA
% mdparagraph:P062:begin
Place ne before the finite verb to deny a proposition. Some common verbs have contracted negative forms. The following complete table covers \textbf{the present indicative eom/\allowbreak{}eart/\allowbreak{}is/\allowbreak{}sindon and the preterite indicative wæs/\allowbreak{}wǣre/\allowbreak{}wǣron only}. It does not replace every form of bēon/\allowbreak{}wesan. Ne eom, ne eart and ne sindon remain two words; nis and the negative preterites are contracted. Separate ne with the ordinary positive form is also valid.\textsuperscript{3}
% mdparagraph:P062:end

% teaching: NEW-U15-GRAM-U15-NEGATIVE-COPULA
% teaching: NEW-U15-PAR-WESAN-NEG-IND
\reserve{13\baselineskip}
\reserve{11\baselineskip}
% mdparagraph:P063:begin
\textbf{Negative be-series: indicative}
% mdparagraph:P063:end

% table: NEW-U15-PAR-WESAN-NEG-IND-DISPLAY
% mdtable:T30:begin
% course-table-style: 2026-10
% course-panel:T30-P01:begin
\par\begingroup
\fontsize{10.5}{13.2}\selectfont
\setlength{\LTleft}{0pt}\setlength{\LTright}{\fill}
\noindent
\begin{tabular}{@{}>{\raggedright\arraybackslash}p{39.459mm}>{\raggedright\arraybackslash}p{39.459mm}>{\raggedright\arraybackslash}p{39.459mm}@{}}
\toprule
\textbf{Person and number} & \textbf{Present} & \textbf{Preterite} \\
\midrule
1 sg. & ne eom & næs \\
2 sg. & ne eart & nǣre \\
3 sg. & nis & næs \\
1–3 pl. & ne sindon & nǣron \\
\bottomrule
\end{tabular}
\par\endgroup
% course-panel:T30-P01:end
% mdtable:T30:end
% teaching: NEW-U15-VOC-NABBAN
% mdparagraph:P064:begin
\textbf{nabban} — have not.
% mdparagraph:P064:end

% teaching: NEW-U15-GRAM-U15-NABBAN-INDICATIVE
% mdparagraph:P065:begin
The complete negative possession series is now available for use. Distinguish first singular næbbe from second and third singular næfst, næfð; the present plural has bb again. A present form denies present possession; a preterite form denies past possession.
% mdparagraph:P065:end

% teaching: NEW-U15-GRAM-U15-NABBAN-INDICATIVE
% teaching: NEW-U15-PAR-NABBAN-IND
\reserve{13\baselineskip}
\reserve{11\baselineskip}
% mdparagraph:P066:begin
\textbf{Nabban: indicative}
% mdparagraph:P066:end

% table: NEW-U15-PAR-NABBAN-IND-DISPLAY
% mdtable:T31:begin
% course-table-style: 2026-10
% course-panel:T31-P01:begin
\par\begingroup
\fontsize{10.5}{13.2}\selectfont
\setlength{\LTleft}{0pt}\setlength{\LTright}{\fill}
\noindent
\begin{tabular}{@{}>{\raggedright\arraybackslash}p{39.459mm}>{\raggedright\arraybackslash}p{39.459mm}>{\raggedright\arraybackslash}p{39.459mm}@{}}
\toprule
\textbf{Person and number} & \textbf{Present} & \textbf{Preterite} \\
\midrule
1 sg. & næbbe & næfde \\
2 sg. & næfst & næfdest \\
3 sg. & næfð & næfde \\
1–3 pl. & nabbað & næfdon \\
\bottomrule
\end{tabular}
\par\endgroup
% course-panel:T31-P01:end
% mdtable:T31:end
% teaching: NEW-U15-VOC-WILLAN
% mdparagraph:P067:begin
\textbf{willan} — be unwilling; refuse (with ne).
% mdparagraph:P067:end

% teaching: NEW-U15-GRAM-U15-NEGATIVE-WILLAN
% mdparagraph:P068:begin
Negative willan expresses refusal or unwillingness and governs a plain infinitive. The present and preterite distinguish the time of that unwillingness; English “would not” may need a past refusal, not a future form. These are indicative forms. A supplied subjunctive nolde encountered earlier does not thereby acquire every use of this table.
% mdparagraph:P068:end

% teaching: NEW-U15-GRAM-U15-NEGATIVE-WILLAN
% teaching: NEW-U15-PAR-WILLAN-NEG-IND
\reserve{13\baselineskip}
\reserve{11\baselineskip}
% mdparagraph:P069:begin
\textbf{Negative willan: indicative}
% mdparagraph:P069:end

% table: NEW-U15-PAR-WILLAN-NEG-IND-DISPLAY
% mdtable:T32:begin
% course-table-style: 2026-10
% course-panel:T32-P01:begin
\par\begingroup
\fontsize{10.5}{13.2}\selectfont
\setlength{\LTleft}{0pt}\setlength{\LTright}{\fill}
\noindent
\begin{tabular}{@{}>{\raggedright\arraybackslash}p{39.459mm}>{\raggedright\arraybackslash}p{39.459mm}>{\raggedright\arraybackslash}p{39.459mm}@{}}
\toprule
\textbf{Person and number} & \textbf{Present} & \textbf{Preterite} \\
\midrule
1 sg. & nelle & nolde \\
2 sg. & nelt & noldest \\
3 sg. & nele & nolde \\
1–3 pl. & nellað & noldon \\
\bottomrule
\end{tabular}
\par\endgroup
% course-panel:T32-P01:end
% mdtable:T32:end
\reserve{7\baselineskip}
% mdparagraph:P070:begin
\subsection*{Nān and negative concord}
% mdparagraph:P070:end

% teaching: NEW-U15-VOC-NAN
% mdparagraph:P071:begin
\textbf{nān} — no; no one.
% mdparagraph:P071:end

% mdparagraph:P072:begin
Nān agrees with its noun or stands alone for “no one”. It has the strong long-stem adjective endings in the four cases below. Keep long ā; note masculine accusative singular nānne, feminine plural nāna and neuter nominative/\allowbreak{}accusative plural nān. Only this strong four-case system is required here.
% mdparagraph:P072:end

% teaching: NEW-U15-PAR-NAN-STRONG-NGDA
\reserve{11\baselineskip}
\reserve{16\baselineskip}
% mdparagraph:P073:begin
\textbf{Nān: strong declension}
% mdparagraph:P073:end

% table: NEW-U15-PAR-NAN-STRONG-NGDA-DISPLAY
% mdtable:T33:begin
% course-table-style: 2026-10
% course-panel:T33-P01:begin
\par\begingroup
\fontsize{10.5}{13.2}\selectfont
\setlength{\LTleft}{0pt}\setlength{\LTright}{\fill}
\noindent
\begin{tabular}{@{}>{\raggedright\arraybackslash}p{22.0mm}>{\raggedright\arraybackslash}p{31.188mm}>{\raggedright\arraybackslash}p{31.188mm}>{\raggedright\arraybackslash}p{31.188mm}@{}}
\toprule
\textbf{Case} & \textbf{Masculine singular} & \textbf{Neuter singular} & \textbf{Feminine singular} \\
\midrule
Nominative & nān & nān & nān \\
Genitive & nānes & nānes & nānre \\
Dative & nānum & nānum & nānre \\
Accusative & nānne & nān & nāne \\
\bottomrule
\end{tabular}
\par\endgroup
% course-panel:T33-P01:end
\par\addvspace{5pt}
% course-panel:T33-P02:begin
\par\begingroup
\fontsize{10.5}{13.2}\selectfont
\setlength{\LTleft}{0pt}\setlength{\LTright}{\fill}
\noindent
\begin{tabular}{@{}>{\raggedright\arraybackslash}p{22.0mm}>{\raggedright\arraybackslash}p{31.188mm}>{\raggedright\arraybackslash}p{31.188mm}>{\raggedright\arraybackslash}p{31.188mm}@{}}
\toprule
\textbf{Case} & \textbf{Masculine plural} & \textbf{Neuter plural} & \textbf{Feminine plural} \\
\midrule
Nominative & nāne & nān & nāna \\
Genitive & nānra & nānra & nānra \\
Dative & nānum & nānum & nānum \\
Accusative & nāne & nān & nāna \\
\bottomrule
\end{tabular}
\par\endgroup
% course-panel:T33-P02:end
% mdtable:T33:end

% teaching: NEW-U15-GRAM-U15-NEGATIVE-CONCORD
% mdparagraph:P074:begin
Nān, nāwiht and nǣfre can accompany finite ne or a contracted negative verb. They express \textbf{one denial}. This is negative concord: adding a second negative does not turn the clause into an affirmation. Determine what is denied from the roles of the negative expressions, not by counting them.
% mdparagraph:P074:end

% teaching: NEW-U15-VOC-NAWIHT
% mdparagraph:P075:begin
\textbf{nāwiht} — nothing. Use the selected accusative nāwiht as the thing denied: “nothing”. Other cases are not assigned by this form.
% mdparagraph:P075:end

% teaching: NEW-U15-VOC-NA
% mdparagraph:P076:begin
\textbf{nā} — not, negating a contrasted constituent.
% mdparagraph:P076:end

% teaching: NEW-U15-GRAM-U15-CONSTITUENT-NA
% mdparagraph:P077:begin
Nā can deny one contrasted participant rather than the whole finite assertion. In ic self and nā ōðer, it excludes “another”; the assertion about “I myself” can remain affirmative. Distinguish this constituent use from finite ne. For a negative command, put ne before the imperative; the familiar affirmative command form remains.
% mdparagraph:P077:end

\reserve{7\baselineskip}
% mdparagraph:P078:begin
\section*{3. Indefinite reference, groups and comparison}
% mdparagraph:P078:end

\reserve{7\baselineskip}
% mdparagraph:P079:begin
\subsection*{An unnamed person and each member of a group}
% mdparagraph:P079:end

% teaching: NEW-U15-VOC-AN
% mdparagraph:P080:begin
\textbf{ān} — a; a certain (introducing an unnamed singular individual). Use the earlier singular strong agreement forms.
% mdparagraph:P080:end

% teaching: NEW-U15-GRAM-U15-AN-INDEFINITE
% mdparagraph:P081:begin
Context distinguishes cardinal “one” from indefinite “a certain”. Both agree with the same singular noun. An indefinite use does not create a plural cardinal meaning.
% mdparagraph:P081:end

% teaching: NEW-U15-VOC-HWA
% mdparagraph:P082:begin
\textbf{hwā} — anyone. The earlier personal/\allowbreak{}neuter case forms are unchanged.
% mdparagraph:P082:end

% teaching: NEW-U15-GRAM-U15-HWA-INDEFINITE
% mdparagraph:P083:begin
Hwā/\allowbreak{}hwæt can ask “who?/\allowbreak{}what?” or refer indefinitely to “anyone/\allowbreak{}anything” in an appropriate condition or concession. The clause and its setting decide the function. English “any” cannot be replaced mechanically with hwā in every construction. Use the earlier ǣnig when an agreeing “any” is required.
% mdparagraph:P083:end

% teaching: NEW-U15-VOC-GEHWA
% mdparagraph:P084:begin
\textbf{gehwā /\allowbreak{} gehwæt} — each one; everyone /\allowbreak{} everything.
% mdparagraph:P084:end

% teaching: NEW-U15-GRAM-U15-GEHWA-DISTRIBUTIVE
% mdparagraph:P085:begin
Gehwā refers to each person; gehwæt to everything. Both are singular, even when they distribute over many members. The genitive can identify the group; the distributive itself takes the case required by its own role. There is one personal column, without separate masculine and feminine forms, and no plural. The selected dative is gehwǣm, and the personal accusative is gehwone.\textsuperscript{4}
% mdparagraph:P085:end

% teaching: NEW-U15-PAR-GEHWA-SINGULAR-NGDA
\reserve{11\baselineskip}
\reserve{9\baselineskip}
% mdparagraph:P086:begin
\textbf{Gehwā/\allowbreak{}gehwæt: singular}
% mdparagraph:P086:end

% table: NEW-U15-PAR-GEHWA-SINGULAR-NGDA-DISPLAY
% mdtable:T34:begin
% course-table-style: 2026-10
% course-panel:T34-P01:begin
\par\begingroup
\fontsize{10.5}{13.2}\selectfont
\setlength{\LTleft}{0pt}\setlength{\LTright}{\fill}
\noindent
\begin{tabular}{@{}>{\raggedright\arraybackslash}p{39.459mm}>{\raggedright\arraybackslash}p{39.459mm}>{\raggedright\arraybackslash}p{39.459mm}@{}}
\toprule
\textbf{Case} & \textbf{Neuter singular} & \textbf{Common singular} \\
\midrule
Nominative & gehwæt & gehwā \\
Genitive & gehwæs & gehwæs \\
Dative & gehwǣm & gehwǣm \\
Accusative & gehwæt & gehwone \\
\bottomrule
\end{tabular}
\par\endgroup
% course-panel:T34-P01:end
% mdtable:T34:end
% teaching: NEW-U15-VOC-MANN
% mdparagraph:P087:begin
\textbf{man} — one; people generally, impersonal subject with singular verb. Only nominative singular man is used in this impersonal function.
% mdparagraph:P087:end

% teaching: NEW-U15-GRAM-U15-MAN-IMPERSONAL
% mdparagraph:P088:begin
Impersonal man means “one; people in general” and takes a third-person singular verb. Do not give it the oblique or plural forms of the ordinary count noun mann. The count noun still refers to an individual man and retains its earlier declension.
% mdparagraph:P088:end

% teaching: NEW-U15-GRAM-U15-QUANTITY-PARTITIVES
% mdparagraph:P089:begin
Ǣlc distributes a singular share to each member; eall can gather the members as a plural whole. An explicit plural subject can retain a plural verb while ǣlc his pęning distributes one penny to each person. A genitive after a selection word identifies the whole group; it does not determine the selecting word’s case.
% mdparagraph:P089:end

% teaching: NEW-U15-VOC-IC
% mdparagraph:P090:begin
\textbf{uncer} — of us two, personal genitive after a selection word. Personal genitive uncer means “of us two”. Its form is fixed by the partitive relation; it is distinct from the earlier agreeing possessive uncer, “our (two people’s)”.
% mdparagraph:P090:end

\reserve{7\baselineskip}
% mdparagraph:P091:begin
\subsection*{Free choice and paired coordination}
% mdparagraph:P091:end

% teaching: NEW-U15-VOC-HWAETHER
% mdparagraph:P092:begin
\textbf{hwæðer} — whichever of two, in swā hwæðer swā.
% mdparagraph:P092:end

% teaching: NEW-U15-VOC-SWA
% mdparagraph:P093:begin
\textbf{swā} — free-choice framing: swā hwā/\allowbreak{}hwæt/\allowbreak{}hwæðer swā; the ... the, proportional relation.
% mdparagraph:P093:end

% teaching: NEW-U15-GRAM-U15-FREE-CHOICE
% mdparagraph:P094:begin
Swā hwā swā means “whoever”; swā hwæt swā means “whatever”. Swā hwæðer swā selects “whichever of two”. Give the pronoun the case needed within its clause. A pronoun in the main clause can resume the same referent in a different case, such as a subject within the choice clause and a dative recipient in the main clause. A possible or contingent event may take the subjunctive; free choice does not mechanically require one mood. Cook’s exceptional attracted dative in 31.12 is supplied with that item.
% mdparagraph:P094:end

% teaching: NEW-U15-VOC-AEGTHER
% mdparagraph:P095:begin
\textbf{ǣgðer} — both, in ǣgðer ge…ge.
% mdparagraph:P095:end

% teaching: NEW-U15-GRAM-U15-EITHER-BOTH-COORDINATION
% mdparagraph:P096:begin
Ǣgðer ge … ge means “both … and”. It includes both coordinated constituents. Here ǣgðer is an invariant coordinator; this expression does not require the separate inflectional family meaning “either”.
% mdparagraph:P096:end

\reserve{7\baselineskip}
% mdparagraph:P097:begin
\subsection*{Equality, preference and proportion}
% mdparagraph:P097:end

% teaching: NEW-U15-GRAM-U15-PROPORTIONAL-SWA
% mdparagraph:P098:begin
In equality, swā … swā accompanies a positive adjective: “as … as”. In a proportional comparison, each swā clause contains a comparative: “the more …, the more …”. Each clause retains its own subject, verb and adjective agreement. A proportional relation does not assert equal degree.
% mdparagraph:P098:end

% teaching: NEW-U15-VOC-SWITHE
% mdparagraph:P099:begin
\textbf{swīðe} — rather; more, comparing preference; most; very greatly. The positive swīðe “very; greatly” is familiar. Learn comparative swīðor for preference; superlative swīðost is for recognition.
% mdparagraph:P099:end

% teaching: NEW-U15-PAR-SWITHE-DEGREES
% teaching: NEW-U15-VOC-SWITHE
\reserve{10\baselineskip}
\reserve{8\baselineskip}
% mdparagraph:P100:begin
\textbf{Swīðe: degrees of the adverb}
% mdparagraph:P100:end

% table: NEW-U15-PAR-SWITHE-DEGREES-DISPLAY
% mdtable:T35:begin
\par\begingroup
\fontsize{10.5}{13.2}\selectfont
\setlength{\LTleft}{0pt}\setlength{\LTright}{\fill}
\noindent
\begin{tabular}{@{}>{\raggedright\arraybackslash}p{63.48mm}>{\raggedright\arraybackslash}p{57.71mm}@{}}
\toprule
\textbf{Position} & \textbf{Form} \\
\midrule
Positive & swīðe \\
Comparative & swīðor \\
Superlative & swīðost \\
\bottomrule
\end{tabular}
\par\endgroup

% mdtable:T35:end
% teaching: NEW-U15-GRAM-U15-RATHER-THAN
% mdparagraph:P101:begin
Swīðor … ðonne compares preference: “more/\allowbreak{}rather … than”. The compared expressions retain the cases required by their verb. With lufian, both the preferred object and the alternative are accusative; neither becomes nominative merely by following ðonne.
% mdparagraph:P101:end

% teaching: NEW-U15-VOC-MICEL
% mdparagraph:P102:begin
\textbf{micel} — much; an amount. Use the earlier neuter forms substantivally for an amount.
% mdparagraph:P102:end

% teaching: NEW-U15-VOC-MICEL-COMP
% mdparagraph:P103:begin
\textbf{māre} — more, a greater amount. The neuter substantive māre means “more; a greater amount”.
% mdparagraph:P103:end

% teaching: NEW-U15-RECORD-EALL-U15-DEGREE-ADVERB
% teaching: NEW-U15-VOC-EALL
% mdparagraph:P104:begin
\textbf{eall} — just; exactly, in eall swā. This adverbial use is invariant and distinct from agreeing “all”.
% mdparagraph:P104:end

% teaching: NEW-U15-GRAM-U15-EALL-EQUATIVE
% mdparagraph:P105:begin
Eall strengthens an equative in eall swā micel swā, “just as much as”. Micel denotes the amount. When the comparison leaves “give” understood, the second recipient remains dative, just as the first recipient does.
% mdparagraph:P105:end

% mdparagraph:P106:begin
\textbf{Reference limits.} Āwiht “anything” and hwæthwega “something” are supplied as exact objects in A10; they are not models for a new case system. The variants ōht and nōht, and the further family citations nǣnig, āhwæðer, nāhwæðer, ǣghwæðer, āhwā, ǣghwā and swā hwilc swā, require no new recall here. The earlier scope of other indefinite families remains unchanged. Instrumental degree expressions are reference material, and the possessive sīn is not introduced.
% mdparagraph:P106:end

\reserve{7\baselineskip}
% mdparagraph:P107:begin
\section*{4. Additional verbs}
% mdparagraph:P107:end

% mdparagraph:P108:begin
The following systems supply the forms required by the exercises. Each indicative table is complete for present and preterite, all three persons, singular and plural. Other moods and nonfinite forms are defined separately. A principal part identifies a stem; its citation alone does not require every possible use.\textsuperscript{5}
% mdparagraph:P108:end

\reserve{7\baselineskip}
% mdparagraph:P109:begin
\subsection*{Strong indicative systems}
% mdparagraph:P109:end

% teaching: NEW-U15-VOC-FORGIELDAN
% mdparagraph:P110:begin
\textbf{forgieldan} — pay; requite; recompense. Strong III; present forgield-, preterite singular forgeald-, preterite plural forguld-, past participle forgolden.
% mdparagraph:P110:end

% mdparagraph:P111:begin
The second- and third-singular present forms forgieltst and forgielt use the selected t cluster. The preterite second singular uses the plural stem plus -e; first and third singular have the bare singular stem. All present plurals end in -að and preterite plurals in -on. The prefix remains attached throughout.
% mdparagraph:P111:end

% teaching: NEW-U15-PAR-FORGIELDAN-IND
\reserve{13\baselineskip}
\reserve{11\baselineskip}
% mdparagraph:P112:begin
\textbf{Forgieldan: indicative}
% mdparagraph:P112:end

% table: NEW-U15-PAR-FORGIELDAN-IND-DISPLAY
% mdtable:T36:begin
% course-table-style: 2026-10
% course-panel:T36-P01:begin
\par\begingroup
\fontsize{10.5}{13.2}\selectfont
\setlength{\LTleft}{0pt}\setlength{\LTright}{\fill}
\noindent
\begin{tabular}{@{}>{\raggedright\arraybackslash}p{39.459mm}>{\raggedright\arraybackslash}p{39.459mm}>{\raggedright\arraybackslash}p{39.459mm}@{}}
\toprule
\textbf{Person and number} & \textbf{Present} & \textbf{Preterite} \\
\midrule
1 sg. & forgielde & forgeald \\
2 sg. & forgieltst & forgulde \\
3 sg. & forgielt & forgeald \\
1–3 pl. & forgieldað & forguldon \\
\bottomrule
\end{tabular}
\par\endgroup
% course-panel:T36-P01:end
% mdtable:T36:end
% teaching: NEW-U15-VOC-TOWEORPAN
% mdparagraph:P113:begin
\textbf{tōweorpan} — break; overthrow, here break a commandment. Strong III; present tōweorp-, preterite singular tōwearp-, preterite plural tōwurp-, past participle tōworpen.
% mdparagraph:P113:end

% mdparagraph:P114:begin
The second- and third-singular present stem is tōwierp-. The preterite follows the same person/\allowbreak{}number distribution as forgieldan.
% mdparagraph:P114:end

% teaching: NEW-U15-PAR-TOWEORPAN-IND
\reserve{13\baselineskip}
\reserve{11\baselineskip}
% mdparagraph:P115:begin
\textbf{Tōweorpan: indicative}
% mdparagraph:P115:end

% table: NEW-U15-PAR-TOWEORPAN-IND-DISPLAY
% mdtable:T37:begin
% course-table-style: 2026-10
% course-panel:T37-P01:begin
\par\begingroup
\fontsize{10.5}{13.2}\selectfont
\setlength{\LTleft}{0pt}\setlength{\LTright}{\fill}
\noindent
\begin{tabular}{@{}>{\raggedright\arraybackslash}p{39.459mm}>{\raggedright\arraybackslash}p{39.459mm}>{\raggedright\arraybackslash}p{39.459mm}@{}}
\toprule
\textbf{Person and number} & \textbf{Present} & \textbf{Preterite} \\
\midrule
1 sg. & tōweorpe & tōwearp \\
2 sg. & tōwierpst & tōwurpe \\
3 sg. & tōwierpð & tōwearp \\
1–3 pl. & tōweorpað & tōwurpon \\
\bottomrule
\end{tabular}
\par\endgroup
% course-panel:T37-P01:end
% mdtable:T37:end
% teaching: NEW-U15-VOC-ONGIETAN
% mdparagraph:P116:begin
\textbf{ongietan} — understand; perceive. Strong V; present ongiet-, preterite singular ongeat-, preterite plural ongēat-, past participle ongieten.
% mdparagraph:P116:end

% mdparagraph:P117:begin
Second singular ongietst and third singular ongiett retain the selected t-cluster spellings. Third singular ongiet is a recognized alternative. Preterite second singular ongēate takes the plural stem; first and third singular have ongeat. The principal-part infinitive and past participle are reference citations; the complete new system here is indicative.
% mdparagraph:P117:end

% teaching: NEW-U15-PAR-ONGIETAN-IND
\reserve{13\baselineskip}
\reserve{11\baselineskip}
% mdparagraph:P118:begin
\textbf{Ongietan: indicative}
% mdparagraph:P118:end

% table: NEW-U15-PAR-ONGIETAN-IND-DISPLAY
% mdtable:T38:begin
% course-table-style: 2026-10
% course-panel:T38-P01:begin
\par\begingroup
\fontsize{10.5}{13.2}\selectfont
\setlength{\LTleft}{0pt}\setlength{\LTright}{\fill}
\noindent
\begin{tabular}{@{}>{\raggedright\arraybackslash}p{39.459mm}>{\raggedright\arraybackslash}p{39.459mm}>{\raggedright\arraybackslash}p{39.459mm}@{}}
\toprule
\textbf{Person and number} & \textbf{Present} & \textbf{Preterite} \\
\midrule
1 sg. & ongiete & ongeat \\
2 sg. & ongietst & ongēate \\
3 sg. & ongiett & ongeat \\
1–3 pl. & ongietað & ongēaton \\
\bottomrule
\end{tabular}
\par\endgroup
% course-panel:T38-P01:end
% mdtable:T38:end
% teaching: NEW-U15-VOC-GEHATAN
% mdparagraph:P119:begin
\textbf{gehātan} — promise. Strong VII. This indicative is a review of the earlier system.
% mdparagraph:P119:end

% mdparagraph:P120:begin
Gehātan keeps its prefix and the familiar hātan stem changes: ā becomes ǣ in the second- and third-singular present, while the preterite has ē. This review does not add a new prefixed negative or subjunctive system.
% mdparagraph:P120:end

% teaching: NEW-U15-PAR-GEHATAN-IND
\reserve{13\baselineskip}
\reserve{11\baselineskip}
% mdparagraph:P121:begin
\textbf{Gehātan: indicative review}
% mdparagraph:P121:end

% table: NEW-U15-PAR-GEHATAN-IND-DISPLAY
% mdtable:T39:begin
% course-table-style: 2026-10
% course-panel:T39-P01:begin
\par\begingroup
\fontsize{10.5}{13.2}\selectfont
\setlength{\LTleft}{0pt}\setlength{\LTright}{\fill}
\noindent
\begin{tabular}{@{}>{\raggedright\arraybackslash}p{39.459mm}>{\raggedright\arraybackslash}p{39.459mm}>{\raggedright\arraybackslash}p{39.459mm}@{}}
\toprule
\textbf{Person and number} & \textbf{Present} & \textbf{Preterite} \\
\midrule
1 sg. & gehāte & gehēt \\
2 sg. & gehǣtst & gehēte \\
3 sg. & gehǣtt & gehēt \\
1–3 pl. & gehātað & gehēton \\
\bottomrule
\end{tabular}
\par\endgroup
% course-panel:T39-P01:end
% mdtable:T39:end
\reserve{7\baselineskip}
% mdparagraph:P122:begin
\subsection*{Weak indicative systems}
% mdparagraph:P122:end

% mdparagraph:P123:begin
The long-stem weak I verbs below use present -e, -st, -ð, -að and preterite -de, -dest, -de, -don. The repeated plural ending serves all three persons. Dēman also has the complete subjunctive shown separately.
% mdparagraph:P123:end

% teaching: NEW-U15-VOC-WENAN
% mdparagraph:P124:begin
\textbf{wēnan} — suppose; expect; think. Weak I.
% mdparagraph:P124:end

% teaching: NEW-U15-PAR-WENAN-IND
\reserve{13\baselineskip}
\reserve{11\baselineskip}
% mdparagraph:P125:begin
\textbf{Wēnan: indicative}
% mdparagraph:P125:end

% table: NEW-U15-PAR-WENAN-IND-DISPLAY
% mdtable:T40:begin
% course-table-style: 2026-10
% course-panel:T40-P01:begin
\par\begingroup
\fontsize{10.5}{13.2}\selectfont
\setlength{\LTleft}{0pt}\setlength{\LTright}{\fill}
\noindent
\begin{tabular}{@{}>{\raggedright\arraybackslash}p{39.459mm}>{\raggedright\arraybackslash}p{39.459mm}>{\raggedright\arraybackslash}p{39.459mm}@{}}
\toprule
\textbf{Person and number} & \textbf{Present} & \textbf{Preterite} \\
\midrule
1 sg. & wēne & wēnde \\
2 sg. & wēnst & wēndest \\
3 sg. & wēnð & wēnde \\
1–3 pl. & wēnað & wēndon \\
\bottomrule
\end{tabular}
\par\endgroup
% course-panel:T40-P01:end
% mdtable:T40:end
% teaching: NEW-U15-VOC-FORBAERNAN
% mdparagraph:P126:begin
\textbf{forbærnan} — parch; scorch. Weak I.
% mdparagraph:P126:end

% teaching: NEW-U15-PAR-FORBAERNAN-IND
\reserve{13\baselineskip}
\reserve{11\baselineskip}
% mdparagraph:P127:begin
\textbf{Forbærnan: indicative}
% mdparagraph:P127:end

% table: NEW-U15-PAR-FORBAERNAN-IND-DISPLAY
% mdtable:T41:begin
% course-table-style: 2026-10
% course-panel:T41-P01:begin
\par\begingroup
\fontsize{10.5}{13.2}\selectfont
\setlength{\LTleft}{0pt}\setlength{\LTright}{\fill}
\noindent
\begin{tabular}{@{}>{\raggedright\arraybackslash}p{39.459mm}>{\raggedright\arraybackslash}p{39.459mm}>{\raggedright\arraybackslash}p{39.459mm}@{}}
\toprule
\textbf{Person and number} & \textbf{Present} & \textbf{Preterite} \\
\midrule
1 sg. & forbærne & forbærnde \\
2 sg. & forbærnst & forbærndest \\
3 sg. & forbærnð & forbærnde \\
1–3 pl. & forbærnað & forbærndon \\
\bottomrule
\end{tabular}
\par\endgroup
% course-panel:T41-P01:end
% mdtable:T41:end
% teaching: NEW-U15-VOC-LAERAN
\par\begin{minipage}{\linewidth}\raggedright
% mdparagraph:P128:begin
\textbf{lǣran} — teach; instruct. Weak I.
% mdparagraph:P128:end

% teaching: NEW-U15-PAR-LAERAN-IND
% mdparagraph:P129:begin
\textbf{Lǣran: indicative}
% mdparagraph:P129:end

% table: NEW-U15-PAR-LAERAN-IND-DISPLAY
% mdtable:T42:begin
% course-table-style: 2026-10
% course-panel:T42-P01:begin
\par\begingroup
\fontsize{10.5}{13.2}\selectfont
\setlength{\LTleft}{0pt}\setlength{\LTright}{\fill}
\noindent
\begin{tabular}{@{}>{\raggedright\arraybackslash}p{39.459mm}>{\raggedright\arraybackslash}p{39.459mm}>{\raggedright\arraybackslash}p{39.459mm}@{}}
\toprule
\textbf{Person and number} & \textbf{Present} & \textbf{Preterite} \\
\midrule
1 sg. & lǣre & lǣrde \\
2 sg. & lǣrst & lǣrdest \\
3 sg. & lǣrð & lǣrde \\
1–3 pl. & lǣrað & lǣrdon \\
\bottomrule
\end{tabular}
\par\endgroup
% course-panel:T42-P01:end
% mdtable:T42:end
\end{minipage}\par\addvspace{4pt}
% teaching: NEW-U15-VOC-DEMAN
% mdparagraph:P130:begin
\textbf{dēman} — judge; condemn. Weak I.
% mdparagraph:P130:end

% teaching: NEW-U15-PAR-DEMAN-IND
\reserve{13\baselineskip}
\reserve{11\baselineskip}
% mdparagraph:P131:begin
\textbf{Dēman: indicative}
% mdparagraph:P131:end

% table: NEW-U15-PAR-DEMAN-IND-DISPLAY
% mdtable:T43:begin
% course-table-style: 2026-10
% course-panel:T43-P01:begin
\par\begingroup
\fontsize{10.5}{13.2}\selectfont
\setlength{\LTleft}{0pt}\setlength{\LTright}{\fill}
\noindent
\begin{tabular}{@{}>{\raggedright\arraybackslash}p{39.459mm}>{\raggedright\arraybackslash}p{39.459mm}>{\raggedright\arraybackslash}p{39.459mm}@{}}
\toprule
\textbf{Person and number} & \textbf{Present} & \textbf{Preterite} \\
\midrule
1 sg. & dēme & dēmde \\
2 sg. & dēmst & dēmdest \\
3 sg. & dēmð & dēmde \\
1–3 pl. & dēmað & dēmdon \\
\bottomrule
\end{tabular}
\par\endgroup
% course-panel:T43-P01:end
% mdtable:T43:end
% mdparagraph:P132:begin
The subjunctive singular has one form for all three persons; the plural has -n. Dēme can therefore be indicative first singular or present subjunctive singular; the clause and its function decide. Use the earlier distinctions among assertion, purpose and wish.
% mdparagraph:P132:end

% teaching: NEW-U15-PAR-DEMAN-SUBJ
\reserve{13\baselineskip}
\reserve{11\baselineskip}
% mdparagraph:P133:begin
\textbf{Dēman: subjunctive}
% mdparagraph:P133:end

% table: NEW-U15-PAR-DEMAN-SUBJ-DISPLAY
% mdtable:T44:begin
% course-table-style: 2026-10
% course-panel:T44-P01:begin
\par\begingroup
\fontsize{10.5}{13.2}\selectfont
\setlength{\LTleft}{0pt}\setlength{\LTright}{\fill}
\noindent
\begin{tabular}{@{}>{\raggedright\arraybackslash}p{39.459mm}>{\raggedright\arraybackslash}p{39.459mm}>{\raggedright\arraybackslash}p{39.459mm}@{}}
\toprule
\textbf{Person and number} & \textbf{Present} & \textbf{Preterite} \\
\midrule
1 sg. & dēme & dēmde \\
2 sg. & dēme & dēmde \\
3 sg. & dēme & dēmde \\
1–3 pl. & dēmen & dēmden \\
\bottomrule
\end{tabular}
\par\endgroup
% course-panel:T44-P01:end
% mdtable:T44:end
% teaching: NEW-U15-VOC-HYRIAN
% mdparagraph:P134:begin
\textbf{hȳrian} — hire. Weak II. Keep the long ȳ throughout.
% mdparagraph:P134:end

% mdparagraph:P135:begin
Hȳrian follows lufian: present singular hȳrie, hȳrast, hȳrað, plural hȳriað; the preterite uses hȳrod- and the ordinary weak endings. Distinguish third-singular -að from plural -iað.
% mdparagraph:P135:end

% teaching: NEW-U15-PAR-HYRIAN-IND
\reserve{13\baselineskip}
\reserve{11\baselineskip}
% mdparagraph:P136:begin
\textbf{Hȳrian: indicative}
% mdparagraph:P136:end

% table: NEW-U15-PAR-HYRIAN-IND-DISPLAY
% mdtable:T45:begin
% course-table-style: 2026-10
% course-panel:T45-P01:begin
\par\begingroup
\fontsize{10.5}{13.2}\selectfont
\setlength{\LTleft}{0pt}\setlength{\LTright}{\fill}
\noindent
\begin{tabular}{@{}>{\raggedright\arraybackslash}p{39.459mm}>{\raggedright\arraybackslash}p{39.459mm}>{\raggedright\arraybackslash}p{39.459mm}@{}}
\toprule
\textbf{Person and number} & \textbf{Present} & \textbf{Preterite} \\
\midrule
1 sg. & hȳrie & hȳrode \\
2 sg. & hȳrast & hȳrodest \\
3 sg. & hȳrað & hȳrode \\
1–3 pl. & hȳriað & hȳrodon \\
\bottomrule
\end{tabular}
\par\endgroup
% course-panel:T45-P01:end
% mdtable:T45:end
\reserve{7\baselineskip}
% mdparagraph:P137:begin
\subsection*{Commands}
% mdparagraph:P137:end

% mdparagraph:P138:begin
Each table gives the complete second-person singular/\allowbreak{}plural command pair. Forgieldan has the strong singular stem without an ending. The weak I singulars ālīes and ætīew also have no ending; their plurals take -að. Weak II cleopa has singular -a and plural cleopiað.
% mdparagraph:P138:end

% teaching: NEW-U15-PAR-FORGIELDAN-IMP
\reserve{9\baselineskip}
\reserve{7\baselineskip}
% mdparagraph:P139:begin
\textbf{Forgieldan: imperative}
% mdparagraph:P139:end

% table: NEW-U15-PAR-FORGIELDAN-IMP-DISPLAY
% mdtable:T46:begin
% course-table-style: 2026-10
% course-panel:T46-P01:begin
\par\begingroup
\fontsize{10.5}{13.2}\selectfont
\setlength{\LTleft}{0pt}\setlength{\LTright}{\fill}
\noindent
\begin{tabular}{@{}>{\raggedright\arraybackslash}p{63.342mm}>{\raggedright\arraybackslash}p{57.846mm}@{}}
\toprule
\textbf{Person and number} & \textbf{Imperative} \\
\midrule
2 sg. & forgield \\
2 pl. & forgieldað \\
\bottomrule
\end{tabular}
\par\endgroup
% course-panel:T46-P01:end
% mdtable:T46:end
% teaching: NEW-U15-VOC-CLEOPIAN
% mdparagraph:P140:begin
\textbf{cleopian} — call; summon. Weak II; only the command pair is newly required for use.
% mdparagraph:P140:end

% teaching: NEW-U15-PAR-CLEOPIAN-IMP
\reserve{9\baselineskip}
\reserve{7\baselineskip}
% mdparagraph:P141:begin
\textbf{Cleopian: imperative}
% mdparagraph:P141:end

% table: NEW-U15-PAR-CLEOPIAN-IMP-DISPLAY
% mdtable:T47:begin
% course-table-style: 2026-10
% course-panel:T47-P01:begin
\par\begingroup
\fontsize{10.5}{13.2}\selectfont
\setlength{\LTleft}{0pt}\setlength{\LTright}{\fill}
\noindent
\begin{tabular}{@{}>{\raggedright\arraybackslash}p{63.342mm}>{\raggedright\arraybackslash}p{57.846mm}@{}}
\toprule
\textbf{Person and number} & \textbf{Imperative} \\
\midrule
2 sg. & cleopa \\
2 pl. & cleopiað \\
\bottomrule
\end{tabular}
\par\endgroup
% course-panel:T47-P01:end
% mdtable:T47:end
% teaching: NEW-U15-VOC-ALIESAN
% mdparagraph:P142:begin
\textbf{ālīesan} — deliver; release. Weak I; only the command pair is newly required for use.
% mdparagraph:P142:end

% teaching: NEW-U15-PAR-ALIESAN-IMP
\reserve{9\baselineskip}
\reserve{7\baselineskip}
% mdparagraph:P143:begin
\textbf{Ālīesan: imperative}
% mdparagraph:P143:end

% table: NEW-U15-PAR-ALIESAN-IMP-DISPLAY
% mdtable:T48:begin
% course-table-style: 2026-10
% course-panel:T48-P01:begin
\par\begingroup
\fontsize{10.5}{13.2}\selectfont
\setlength{\LTleft}{0pt}\setlength{\LTright}{\fill}
\noindent
\begin{tabular}{@{}>{\raggedright\arraybackslash}p{63.342mm}>{\raggedright\arraybackslash}p{57.846mm}@{}}
\toprule
\textbf{Person and number} & \textbf{Imperative} \\
\midrule
2 sg. & ālīes \\
2 pl. & ālīesað \\
\bottomrule
\end{tabular}
\par\endgroup
% course-panel:T48-P01:end
% mdtable:T48:end
% mdparagraph:P144:begin
\textbf{ætīewan} — show; the earlier verb, now with its complete command pair. Retain short initial æ in ætīew.
% mdparagraph:P144:end

% teaching: NEW-U15-PAR-AETIEWAN-IMP
\reserve{9\baselineskip}
\reserve{7\baselineskip}
% mdparagraph:P145:begin
\textbf{Ætīewan: imperative}
% mdparagraph:P145:end

% table: NEW-U15-PAR-AETIEWAN-IMP-DISPLAY
% mdtable:T49:begin
% course-table-style: 2026-10
% course-panel:T49-P01:begin
\par\begingroup
\fontsize{10.5}{13.2}\selectfont
\setlength{\LTleft}{0pt}\setlength{\LTright}{\fill}
\noindent
\begin{tabular}{@{}>{\raggedright\arraybackslash}p{63.342mm}>{\raggedright\arraybackslash}p{57.846mm}@{}}
\toprule
\textbf{Person and number} & \textbf{Imperative} \\
\midrule
2 sg. & ætīew \\
2 pl. & ætīewað \\
\bottomrule
\end{tabular}
\par\endgroup
% course-panel:T49-P01:end
% mdtable:T49:end
\reserve{7\baselineskip}
% mdparagraph:P146:begin
\subsection*{Infinitives}
% mdparagraph:P146:end

% mdparagraph:P147:begin
These are complete pairs of plain and inflected infinitives. Use the plain form after the taught modal verbs. The inflected infinitive follows tō in the uses already learned; particular difficult frames are supplied locally. The pair does not introduce a finite system for the verb.
% mdparagraph:P147:end

% teaching: NEW-U15-PAR-FORGIELDAN-INF
\reserve{9\baselineskip}
\reserve{7\baselineskip}
% mdparagraph:P148:begin
\textbf{Forgieldan: infinitives}
% mdparagraph:P148:end

% table: NEW-U15-PAR-FORGIELDAN-INF-DISPLAY
% mdtable:T50:begin
\par\begingroup
\fontsize{10.5}{13.2}\selectfont
\setlength{\LTleft}{0pt}\setlength{\LTright}{\fill}
\noindent
\begin{tabular}{@{}>{\raggedright\arraybackslash}p{66.31mm}>{\raggedright\arraybackslash}p{54.88mm}@{}}
\toprule
\textbf{Position} & \textbf{Form} \\
\midrule
Plain infinitive & forgieldan \\
Inflected infinitive & forgieldanne \\
\bottomrule
\end{tabular}
\par\endgroup

% mdtable:T50:end
% mdparagraph:P149:begin
\textbf{wyrcean} — work; do. The selected pair retains ea before the inflected ending.
% mdparagraph:P149:end

% teaching: NEW-U15-PAR-WYRCEAN-INF
\reserve{9\baselineskip}
\reserve{7\baselineskip}
% mdparagraph:P150:begin
\textbf{Wyrcean: infinitives}
% mdparagraph:P150:end

% table: NEW-U15-PAR-WYRCEAN-INF-DISPLAY
% mdtable:T51:begin
\par\begingroup
\fontsize{10.5}{13.2}\selectfont
\setlength{\LTleft}{0pt}\setlength{\LTright}{\fill}
\noindent
\begin{tabular}{@{}>{\raggedright\arraybackslash}p{78.74mm}>{\raggedright\arraybackslash}p{42.45mm}@{}}
\toprule
\textbf{Position} & \textbf{Form} \\
\midrule
Plain infinitive & wyrcean \\
Inflected infinitive after tō & wyrceanne \\
\bottomrule
\end{tabular}
\par\endgroup

% mdtable:T51:end
% mdparagraph:P151:begin
\textbf{onfōn} — receive; keep the contracted stem in both infinitives.
% mdparagraph:P151:end

% teaching: NEW-U15-PAR-ONFON-INF
\reserve{9\baselineskip}
\reserve{7\baselineskip}
% mdparagraph:P152:begin
\textbf{Onfōn: infinitives}
% mdparagraph:P152:end

% table: NEW-U15-PAR-ONFON-INF-DISPLAY
% mdtable:T52:begin
\par\begingroup
\fontsize{10.5}{13.2}\selectfont
\setlength{\LTleft}{0pt}\setlength{\LTright}{\fill}
\noindent
\begin{tabular}{@{}>{\raggedright\arraybackslash}p{76.63mm}>{\raggedright\arraybackslash}p{44.55mm}@{}}
\toprule
\textbf{Position} & \textbf{Form} \\
\midrule
Infinitive (citation form) & onfōn \\
Inflected infinitive & onfōnne \\
\bottomrule
\end{tabular}
\par\endgroup

% mdtable:T52:end
\par\begin{minipage}{\linewidth}\raggedright
% mdparagraph:P153:begin
\textbf{forlǣtan} — forsake; leave. Keep the long ǣ in both infinitives.
% mdparagraph:P153:end

% teaching: NEW-U15-PAR-FORLAETAN-INF
% mdparagraph:P154:begin
\textbf{Forlǣtan: infinitives}
% mdparagraph:P154:end

% table: NEW-U15-PAR-FORLAETAN-INF-DISPLAY
% mdtable:T53:begin
\par\begingroup
\fontsize{10.5}{13.2}\selectfont
\setlength{\LTleft}{0pt}\setlength{\LTright}{\fill}
\noindent
\begin{tabular}{@{}>{\raggedright\arraybackslash}p{71.72mm}>{\raggedright\arraybackslash}p{49.46mm}@{}}
\toprule
\textbf{Position} & \textbf{Form} \\
\midrule
Plain infinitive & forlǣtan \\
Inflected infinitive & forlǣtanne \\
\bottomrule
\end{tabular}
\par\endgroup

% mdtable:T53:end
\end{minipage}\par\addvspace{4pt}
% teaching: NEW-U15-VOC-BLOSTMIAN
% mdparagraph:P155:begin
\textbf{blōstmian} — blossom; bloom. Weak II; only the infinitive pair is required for use here.
% mdparagraph:P155:end

% teaching: NEW-U15-PAR-BLOSTMIAN-INF
\reserve{9\baselineskip}
\reserve{7\baselineskip}
% mdparagraph:P156:begin
\textbf{Blōstmian: infinitives}
% mdparagraph:P156:end

% table: NEW-U15-PAR-BLOSTMIAN-INF-DISPLAY
% mdtable:T54:begin
\par\begingroup
\fontsize{10.5}{13.2}\selectfont
\setlength{\LTleft}{0pt}\setlength{\LTright}{\fill}
\noindent
\begin{tabular}{@{}>{\raggedright\arraybackslash}p{68.91mm}>{\raggedright\arraybackslash}p{52.28mm}@{}}
\toprule
\textbf{Position} & \textbf{Form} \\
\midrule
Plain infinitive & blōstmian \\
Inflected infinitive & blōstmianne \\
\bottomrule
\end{tabular}
\par\endgroup

% mdtable:T54:end
\reserve{7\baselineskip}
% mdparagraph:P157:begin
\subsection*{Participial bases}
% mdparagraph:P157:end

% mdparagraph:P158:begin
Each pair gives the present and past participial base, without adding a complete adjective declension. The present participle can express an accompanying activity in the earlier predicative construction; the past participle can enter the earlier passive. Keep full -ed in lǣred and hǣled. Standan is strong VI: its present base standende retains n, and its past base is standen.
% mdparagraph:P158:end

% teaching: NEW-U15-PAR-FORGIELDAN-PART
\reserve{9\baselineskip}
\reserve{7\baselineskip}
% mdparagraph:P159:begin
\textbf{Forgieldan: participial bases}
% mdparagraph:P159:end

% table: NEW-U15-PAR-FORGIELDAN-PART-DISPLAY
% mdtable:T55:begin
\par\begingroup
\fontsize{10.5}{13.2}\selectfont
\setlength{\LTleft}{0pt}\setlength{\LTright}{\fill}
\noindent
\begin{tabular}{@{}>{\raggedright\arraybackslash}p{74.28mm}>{\raggedright\arraybackslash}p{46.91mm}@{}}
\toprule
\textbf{Position} & \textbf{Form} \\
\midrule
Present participle (base form) & forgieldende \\
Past participle (base form) & forgolden \\
\bottomrule
\end{tabular}
\par\endgroup

% mdtable:T55:end
% teaching: NEW-U15-PAR-LAERAN-PART
\reserve{9\baselineskip}
\reserve{7\baselineskip}
% mdparagraph:P160:begin
\textbf{Lǣran: participial bases}
% mdparagraph:P160:end

% table: NEW-U15-PAR-LAERAN-PART-DISPLAY
% mdtable:T56:begin
\par\begingroup
\fontsize{10.5}{13.2}\selectfont
\setlength{\LTleft}{0pt}\setlength{\LTright}{\fill}
\noindent
\begin{tabular}{@{}>{\raggedright\arraybackslash}p{79.4mm}>{\raggedright\arraybackslash}p{41.79mm}@{}}
\toprule
\textbf{Position} & \textbf{Form} \\
\midrule
Present participle (base form) & lǣrende \\
Past participle (base form) & lǣred \\
\bottomrule
\end{tabular}
\par\endgroup

% mdtable:T56:end
% teaching: NEW-U15-VOC-HAELAN
% mdparagraph:P161:begin
\textbf{hǣlan} — save. Weak I; the participial bases are required, not a new finite system.
% mdparagraph:P161:end

% teaching: NEW-U15-PAR-HAELAN-PART
\reserve{9\baselineskip}
\reserve{7\baselineskip}
% mdparagraph:P162:begin
\textbf{Hǣlan: participial bases}
% mdparagraph:P162:end

% table: NEW-U15-PAR-HAELAN-PART-DISPLAY
% mdtable:T57:begin
\par\begingroup
\fontsize{10.5}{13.2}\selectfont
\setlength{\LTleft}{0pt}\setlength{\LTright}{\fill}
\noindent
\begin{tabular}{@{}>{\raggedright\arraybackslash}p{79.4mm}>{\raggedright\arraybackslash}p{41.79mm}@{}}
\toprule
\textbf{Position} & \textbf{Form} \\
\midrule
Present participle (base form) & hǣlende \\
Past participle (base form) & hǣled \\
\bottomrule
\end{tabular}
\par\endgroup

% mdtable:T57:end
% mdparagraph:P163:begin
\textbf{standan} — stand. \textbf{giefan} — give. Their meanings are familiar; the following complete base pairs supply the forms needed here.
% mdparagraph:P163:end

% teaching: NEW-U15-PAR-STANDAN-PART
\reserve{9\baselineskip}
\reserve{7\baselineskip}
% mdparagraph:P164:begin
\textbf{Standan: participial bases}
% mdparagraph:P164:end

% table: NEW-U15-PAR-STANDAN-PART-DISPLAY
% mdtable:T58:begin
\par\begingroup
\fontsize{10.5}{13.2}\selectfont
\setlength{\LTleft}{0pt}\setlength{\LTright}{\fill}
\noindent
\begin{tabular}{@{}>{\raggedright\arraybackslash}p{80.04mm}>{\raggedright\arraybackslash}p{41.15mm}@{}}
\toprule
\textbf{Position} & \textbf{Form} \\
\midrule
Present participle (base form) & standende \\
Past participle (citation form) & standen \\
\bottomrule
\end{tabular}
\par\endgroup

% mdtable:T58:end
% teaching: NEW-U15-PAR-GIEFAN-PART
\reserve{9\baselineskip}
\reserve{7\baselineskip}
% mdparagraph:P165:begin
\textbf{Giefan: participial bases}
% mdparagraph:P165:end

% table: NEW-U15-PAR-GIEFAN-PART-DISPLAY
% mdtable:T59:begin
\par\begingroup
\fontsize{10.5}{13.2}\selectfont
\setlength{\LTleft}{0pt}\setlength{\LTright}{\fill}
\noindent
\begin{tabular}{@{}>{\raggedright\arraybackslash}p{80.04mm}>{\raggedright\arraybackslash}p{41.15mm}@{}}
\toprule
\textbf{Position} & \textbf{Form} \\
\midrule
Present participle (base form) & giefende \\
Past participle (citation form) & giefen \\
\bottomrule
\end{tabular}
\par\endgroup

% mdtable:T59:end
\reserve{32\baselineskip}
% mdparagraph:P166:begin
\subsection*{Principal parts for reference}
% mdparagraph:P166:end

% mdparagraph:P167:begin
For strong verbs the order is infinitive, preterite singular, preterite plural, past participle. Weak verbs have infinitive, preterite singular and past participle. This list is available for reference. Only forms in a separately defined system above, or already learned earlier, require independent use; the remaining citation forms are supplied.
% mdparagraph:P167:end

% teaching: NEW-U15-PAR-BLOSTMIAN-INF
% teaching: NEW-U15-PAR-DEMAN-IND
% teaching: NEW-U15-PAR-FORBAERNAN-IND
% teaching: NEW-U15-PAR-FORGIELDAN-IND
% teaching: NEW-U15-PAR-FORGIELDAN-INF
% teaching: NEW-U15-PAR-FORGIELDAN-PART
% teaching: NEW-U15-PAR-FORLAETAN-INF
% teaching: NEW-U15-PAR-GEHATAN-IND
% teaching: NEW-U15-PAR-GIEFAN-PART
% teaching: NEW-U15-PAR-HAELAN-PART
% teaching: NEW-U15-PAR-HYRIAN-IND
% teaching: NEW-U15-PAR-LAERAN-IND
% teaching: NEW-U15-PAR-LAERAN-PART
% teaching: NEW-U15-PAR-ONFON-INF
% teaching: NEW-U15-PAR-ONGIETAN-IND
% teaching: NEW-U15-PAR-STANDAN-PART
% teaching: NEW-U15-PAR-TOWEORPAN-IND
% teaching: NEW-U15-PAR-WENAN-IND
% teaching: NEW-U15-PAR-WYRCEAN-INF
% teaching: NEW-U15-PRINCIPAL-PARTS
% mdtable:T60:begin
\par\begingroup
\fontsize{10.5}{13.2}\selectfont
\setlength{\LTleft}{0pt}\setlength{\LTright}{\fill}
\begin{longtable}{@{}>{\raggedright\arraybackslash}p{19.63mm}>{\raggedright\arraybackslash}p{19.63mm}>{\raggedright\arraybackslash}p{25.61mm}>{\raggedright\arraybackslash}p{23.85mm}>{\raggedright\arraybackslash}p{24.04mm}@{}}
\toprule
\textbf{Infinitive} & \textbf{Type} & \textbf{Preterite singular} & \textbf{Preterite plural} & \textbf{Past participle} \\
\midrule
\endfirsthead
\toprule
\textbf{Infinitive} & \textbf{Type} & \textbf{Preterite singular} & \textbf{Preterite plural} & \textbf{Past participle} \\
\midrule
\endhead
\midrule
\multicolumn{5}{r}{\footnotesize Continued on the following page.}\\
\endfoot
\bottomrule
\endlastfoot
dēman & weak & dēmde & — & dēmed \\
forgieldan & strong & forgeald & forguldon & forgolden \\
hȳrian & weak & hȳrode & — & hȳrod \\
wēnan & weak & wēnde & — & wēned \\
tōweorpan & strong & tōwearp & tōwurpon & tōworpen \\
forbærnan & weak & forbærnde & — & forbærned \\
lǣran & weak & lǣrde & — & lǣred \\
ongietan & strong & ongeat & ongēaton & ongieten \\
gehātan & strong & gehēt & gehēton & gehāten \\
cleopian & weak & cleopode & — & cleopod \\
hǣlan & weak & hǣlde & — & hǣled \\
ālīesan & weak & ālīesde & — & ālīesed \\
blōstmian & weak & blōstmode & — & blōstmod \\
ætīewan & weak & ætīewde & — & ætīewed \\
wyrcean & weak & worhte & — & worht \\
onfōn & strong & onfēng & onfēngon & onfangen \\
forlǣtan & strong & forlēt & forlēton & forlǣten \\
standan & strong & stōd & stōdon & standen \\
giefan & strong & geaf & gēafon & giefen \\
\end{longtable}
\par\endgroup

% mdtable:T60:end
\reserve{7\baselineskip}
\par\begin{minipage}{\linewidth}\raggedright
% mdparagraph:P168:begin
\section*{5. Adjectives, meanings and government}
% mdparagraph:P168:end

% mdparagraph:P169:begin
\subsection*{Complete adjective systems}
% mdparagraph:P169:end

% mdparagraph:P170:begin
Each adjective has complete strong and weak tables for the four cases, both numbers and all three genders. Use strong forms in the familiar indefinite and predicative constructions, weak forms after the definite determiner as previously learned. No instrumental system is added.\textsuperscript{6}
% mdparagraph:P170:end
\end{minipage}\par\addvspace{4pt}


% mdparagraph:P171:begin
Riht and gelīc follow the ordinary long-stem pattern: strong feminine nominative singular and neuter nominative/\allowbreak{}accusative plural have no ending. The selected strong feminine plural has -a, masculine plural -e. Weak genitive plural has -ra.
% mdparagraph:P171:end

% teaching: NEW-U15-VOC-RIHT
% mdparagraph:P172:begin
\textbf{riht} — right; just.
% mdparagraph:P172:end

% teaching: NEW-U15-PAR-RIHT-STRONG-NGDA
\reserve{11\baselineskip}
\reserve{16\baselineskip}
% mdparagraph:P173:begin
\textbf{Riht: strong}
% mdparagraph:P173:end

% table: NEW-U15-PAR-RIHT-STRONG-NGDA-DISPLAY
% mdtable:T61:begin
% course-table-style: 2026-10
% course-panel:T61-P01:begin
\par\begingroup
\fontsize{10.5}{13.2}\selectfont
\setlength{\LTleft}{0pt}\setlength{\LTright}{\fill}
\noindent
\begin{tabular}{@{}>{\raggedright\arraybackslash}p{22.0mm}>{\raggedright\arraybackslash}p{31.188mm}>{\raggedright\arraybackslash}p{31.188mm}>{\raggedright\arraybackslash}p{31.188mm}@{}}
\toprule
\textbf{Case} & \textbf{Masculine singular} & \textbf{Neuter singular} & \textbf{Feminine singular} \\
\midrule
Nominative & riht & riht & riht \\
Genitive & rihtes & rihtes & rihtre \\
Dative & rihtum & rihtum & rihtre \\
Accusative & rihtne & riht & rihte \\
\bottomrule
\end{tabular}
\par\endgroup
% course-panel:T61-P01:end
\par\addvspace{5pt}
% course-panel:T61-P02:begin
\par\begingroup
\fontsize{10.5}{13.2}\selectfont
\setlength{\LTleft}{0pt}\setlength{\LTright}{\fill}
\noindent
\begin{tabular}{@{}>{\raggedright\arraybackslash}p{22.0mm}>{\raggedright\arraybackslash}p{31.188mm}>{\raggedright\arraybackslash}p{31.188mm}>{\raggedright\arraybackslash}p{31.188mm}@{}}
\toprule
\textbf{Case} & \textbf{Masculine plural} & \textbf{Neuter plural} & \textbf{Feminine plural} \\
\midrule
Nominative & rihte & riht & rihta \\
Genitive & rihtra & rihtra & rihtra \\
Dative & rihtum & rihtum & rihtum \\
Accusative & rihte & riht & rihta \\
\bottomrule
\end{tabular}
\par\endgroup
% course-panel:T61-P02:end
% mdtable:T61:end

% teaching: NEW-U15-PAR-RIHT-WEAK-NGDA
\reserve{11\baselineskip}
\reserve{16\baselineskip}
% mdparagraph:P174:begin
\textbf{Riht: weak}
% mdparagraph:P174:end

% table: NEW-U15-PAR-RIHT-WEAK-NGDA-DISPLAY
% mdtable:T62:begin
% course-table-style: 2026-10
% course-panel:T62-P01:begin
\par\begingroup
\fontsize{10.5}{13.2}\selectfont
\setlength{\LTleft}{0pt}\setlength{\LTright}{\fill}
\noindent
\begin{tabular}{@{}>{\raggedright\arraybackslash}p{22.0mm}>{\raggedright\arraybackslash}p{31.188mm}>{\raggedright\arraybackslash}p{31.188mm}>{\raggedright\arraybackslash}p{31.188mm}@{}}
\toprule
\textbf{Case} & \textbf{Masculine singular} & \textbf{Neuter singular} & \textbf{Feminine singular} \\
\midrule
Nominative & rihta & rihte & rihte \\
Genitive & rihtan & rihtan & rihtan \\
Dative & rihtan & rihtan & rihtan \\
Accusative & rihtan & rihte & rihtan \\
\bottomrule
\end{tabular}
\par\endgroup
% course-panel:T62-P01:end
\par\addvspace{5pt}
% course-panel:T62-P02:begin
\par\begingroup
\fontsize{10.5}{13.2}\selectfont
\setlength{\LTleft}{0pt}\setlength{\LTright}{\fill}
\noindent
\begin{tabular}{@{}>{\raggedright\arraybackslash}p{22.0mm}>{\raggedright\arraybackslash}p{31.188mm}>{\raggedright\arraybackslash}p{31.188mm}>{\raggedright\arraybackslash}p{31.188mm}@{}}
\toprule
\textbf{Case} & \textbf{Masculine plural} & \textbf{Neuter plural} & \textbf{Feminine plural} \\
\midrule
Nominative & rihtan & rihtan & rihtan \\
Genitive & rihtra & rihtra & rihtra \\
Dative & rihtum & rihtum & rihtum \\
Accusative & rihtan & rihtan & rihtan \\
\bottomrule
\end{tabular}
\par\endgroup
% course-panel:T62-P02:end
% mdtable:T62:end

% teaching: NEW-U15-VOC-GELIC
\par\begin{minipage}{\linewidth}\raggedright
% mdparagraph:P175:begin
\textbf{gelīc} — equal; like.
% mdparagraph:P175:end

% teaching: NEW-U15-PAR-GELIC-STRONG-NGDA
% mdparagraph:P176:begin
\textbf{Gelīc: strong}
% mdparagraph:P176:end

% table: NEW-U15-PAR-GELIC-STRONG-NGDA-DISPLAY
% mdtable:T63:begin
% course-table-style: 2026-10
% course-panel:T63-P01:begin
\par\begingroup
\fontsize{10.5}{13.2}\selectfont
\setlength{\LTleft}{0pt}\setlength{\LTright}{\fill}
\noindent
\begin{tabular}{@{}>{\raggedright\arraybackslash}p{22.0mm}>{\raggedright\arraybackslash}p{31.188mm}>{\raggedright\arraybackslash}p{31.188mm}>{\raggedright\arraybackslash}p{31.188mm}@{}}
\toprule
\textbf{Case} & \textbf{Masculine singular} & \textbf{Neuter singular} & \textbf{Feminine singular} \\
\midrule
Nominative & gelīc & gelīc & gelīc \\
Genitive & gelīces & gelīces & gelīcre \\
Dative & gelīcum & gelīcum & gelīcre \\
Accusative & gelīcne & gelīc & gelīce \\
\bottomrule
\end{tabular}
\par\endgroup
% course-panel:T63-P01:end
\par\addvspace{5pt}
% course-panel:T63-P02:begin
\par\begingroup
\fontsize{10.5}{13.2}\selectfont
\setlength{\LTleft}{0pt}\setlength{\LTright}{\fill}
\noindent
\begin{tabular}{@{}>{\raggedright\arraybackslash}p{22.0mm}>{\raggedright\arraybackslash}p{31.188mm}>{\raggedright\arraybackslash}p{31.188mm}>{\raggedright\arraybackslash}p{31.188mm}@{}}
\toprule
\textbf{Case} & \textbf{Masculine plural} & \textbf{Neuter plural} & \textbf{Feminine plural} \\
\midrule
Nominative & gelīce & gelīc & gelīca \\
Genitive & gelīcra & gelīcra & gelīcra \\
Dative & gelīcum & gelīcum & gelīcum \\
Accusative & gelīce & gelīc & gelīca \\
\bottomrule
\end{tabular}
\par\endgroup
% course-panel:T63-P02:end
% mdtable:T63:end
\end{minipage}\par\addvspace{4pt}

% teaching: NEW-U15-PAR-GELIC-WEAK-NGDA
\reserve{11\baselineskip}
\reserve{16\baselineskip}
% mdparagraph:P177:begin
\textbf{Gelīc: weak}
% mdparagraph:P177:end

% table: NEW-U15-PAR-GELIC-WEAK-NGDA-DISPLAY
% mdtable:T64:begin
% course-table-style: 2026-10
% course-panel:T64-P01:begin
\par\begingroup
\fontsize{10.5}{13.2}\selectfont
\setlength{\LTleft}{0pt}\setlength{\LTright}{\fill}
\noindent
\begin{tabular}{@{}>{\raggedright\arraybackslash}p{22.0mm}>{\raggedright\arraybackslash}p{31.188mm}>{\raggedright\arraybackslash}p{31.188mm}>{\raggedright\arraybackslash}p{31.188mm}@{}}
\toprule
\textbf{Case} & \textbf{Masculine singular} & \textbf{Neuter singular} & \textbf{Feminine singular} \\
\midrule
Nominative & gelīca & gelīce & gelīce \\
Genitive & gelīcan & gelīcan & gelīcan \\
Dative & gelīcan & gelīcan & gelīcan \\
Accusative & gelīcan & gelīce & gelīcan \\
\bottomrule
\end{tabular}
\par\endgroup
% course-panel:T64-P01:end
\par\addvspace{5pt}
% course-panel:T64-P02:begin
\par\begingroup
\fontsize{10.5}{13.2}\selectfont
\setlength{\LTleft}{0pt}\setlength{\LTright}{\fill}
\noindent
\begin{tabular}{@{}>{\raggedright\arraybackslash}p{22.0mm}>{\raggedright\arraybackslash}p{31.188mm}>{\raggedright\arraybackslash}p{31.188mm}>{\raggedright\arraybackslash}p{31.188mm}@{}}
\toprule
\textbf{Case} & \textbf{Masculine plural} & \textbf{Neuter plural} & \textbf{Feminine plural} \\
\midrule
Nominative & gelīcan & gelīcan & gelīcan \\
Genitive & gelīcra & gelīcra & gelīcra \\
Dative & gelīcum & gelīcum & gelīcum \\
Accusative & gelīcan & gelīcan & gelīcan \\
\bottomrule
\end{tabular}
\par\endgroup
% course-panel:T64-P02:end
% mdtable:T64:end

% teaching: NEW-U15-VOC-MIHTIG
% mdparagraph:P178:begin
\textbf{mihtig} — mighty; powerful.
% mdparagraph:P178:end

% mdparagraph:P179:begin
The selected forms lose i before a vowel-initial ending: mihtges, mihtgum, mihtge, mihtgan. They retain i before a consonantal ending: mihtigne, mihtigre, mihtigra. Strong feminine nominative singular and neuter nominative/\allowbreak{}accusative plural are mihtig, parallel to the earlier hālig. Unsyncopated forms are recognized variants.
% mdparagraph:P179:end

% teaching: NEW-U15-PAR-MIHTIG-STRONG-NGDA
\reserve{11\baselineskip}
\reserve{16\baselineskip}
% mdparagraph:P180:begin
\textbf{Mihtig: strong}
% mdparagraph:P180:end

% table: NEW-U15-PAR-MIHTIG-STRONG-NGDA-DISPLAY
% mdtable:T65:begin
% course-table-style: 2026-10
% course-panel:T65-P01:begin
\par\begingroup
\fontsize{10.5}{13.2}\selectfont
\setlength{\LTleft}{0pt}\setlength{\LTright}{\fill}
\noindent
\begin{tabular}{@{}>{\raggedright\arraybackslash}p{22.0mm}>{\raggedright\arraybackslash}p{31.188mm}>{\raggedright\arraybackslash}p{31.188mm}>{\raggedright\arraybackslash}p{31.188mm}@{}}
\toprule
\textbf{Case} & \textbf{Masculine singular} & \textbf{Neuter singular} & \textbf{Feminine singular} \\
\midrule
Nominative & mihtig & mihtig & mihtig \\
Genitive & mihtges & mihtges & mihtigre \\
Dative & mihtgum & mihtgum & mihtigre \\
Accusative & mihtigne & mihtig & mihtge \\
\bottomrule
\end{tabular}
\par\endgroup
% course-panel:T65-P01:end
\par\addvspace{5pt}
% course-panel:T65-P02:begin
\par\begingroup
\fontsize{10.5}{13.2}\selectfont
\setlength{\LTleft}{0pt}\setlength{\LTright}{\fill}
\noindent
\begin{tabular}{@{}>{\raggedright\arraybackslash}p{22.0mm}>{\raggedright\arraybackslash}p{31.188mm}>{\raggedright\arraybackslash}p{31.188mm}>{\raggedright\arraybackslash}p{31.188mm}@{}}
\toprule
\textbf{Case} & \textbf{Masculine plural} & \textbf{Neuter plural} & \textbf{Feminine plural} \\
\midrule
Nominative & mihtge & mihtig & mihtga \\
Genitive & mihtigra & mihtigra & mihtigra \\
Dative & mihtgum & mihtgum & mihtgum \\
Accusative & mihtge & mihtig & mihtga \\
\bottomrule
\end{tabular}
\par\endgroup
% course-panel:T65-P02:end
% mdtable:T65:end

% teaching: NEW-U15-PAR-MIHTIG-WEAK-NGDA
\reserve{11\baselineskip}
\reserve{16\baselineskip}
% mdparagraph:P181:begin
\textbf{Mihtig: weak}
% mdparagraph:P181:end

% table: NEW-U15-PAR-MIHTIG-WEAK-NGDA-DISPLAY
% mdtable:T66:begin
% course-table-style: 2026-10
% course-panel:T66-P01:begin
\par\begingroup
\fontsize{10.5}{13.2}\selectfont
\setlength{\LTleft}{0pt}\setlength{\LTright}{\fill}
\noindent
\begin{tabular}{@{}>{\raggedright\arraybackslash}p{22.0mm}>{\raggedright\arraybackslash}p{31.188mm}>{\raggedright\arraybackslash}p{31.188mm}>{\raggedright\arraybackslash}p{31.188mm}@{}}
\toprule
\textbf{Case} & \textbf{Masculine singular} & \textbf{Neuter singular} & \textbf{Feminine singular} \\
\midrule
Nominative & mihtga & mihtge & mihtge \\
Genitive & mihtgan & mihtgan & mihtgan \\
Dative & mihtgan & mihtgan & mihtgan \\
Accusative & mihtgan & mihtge & mihtgan \\
\bottomrule
\end{tabular}
\par\endgroup
% course-panel:T66-P01:end
\par\addvspace{5pt}
% course-panel:T66-P02:begin
\par\begingroup
\fontsize{10.5}{13.2}\selectfont
\setlength{\LTleft}{0pt}\setlength{\LTright}{\fill}
\noindent
\begin{tabular}{@{}>{\raggedright\arraybackslash}p{22.0mm}>{\raggedright\arraybackslash}p{31.188mm}>{\raggedright\arraybackslash}p{31.188mm}>{\raggedright\arraybackslash}p{31.188mm}@{}}
\toprule
\textbf{Case} & \textbf{Masculine plural} & \textbf{Neuter plural} & \textbf{Feminine plural} \\
\midrule
Nominative & mihtgan & mihtgan & mihtgan \\
Genitive & mihtigra & mihtigra & mihtigra \\
Dative & mihtgum & mihtgum & mihtgum \\
Accusative & mihtgan & mihtgan & mihtgan \\
\bottomrule
\end{tabular}
\par\endgroup
% course-panel:T66-P02:end
% mdtable:T66:end

% teaching: NEW-U15-VOC-MAERE
% mdparagraph:P182:begin
\textbf{mǣre} — great; renowned.
% mdparagraph:P182:end

% mdparagraph:P183:begin
Mǣre keeps its final e where there is no added ending and drops it before another ending. Note strong feminine nominative singular and neuter nominative/\allowbreak{}accusative plural mǣru. Stem-final r with -re/\allowbreak{}-ra gives mǣrre/\allowbreak{}mǣrra; weak genitive plural is also mǣrra.
% mdparagraph:P183:end

% teaching: NEW-U15-PAR-MAERE-STRONG-NGDA
\reserve{11\baselineskip}
\reserve{16\baselineskip}
% mdparagraph:P184:begin
\textbf{Mǣre: strong}
% mdparagraph:P184:end

% table: NEW-U15-PAR-MAERE-STRONG-NGDA-DISPLAY
% mdtable:T67:begin
% course-table-style: 2026-10
% course-panel:T67-P01:begin
\par\begingroup
\fontsize{10.5}{13.2}\selectfont
\setlength{\LTleft}{0pt}\setlength{\LTright}{\fill}
\noindent
\begin{tabular}{@{}>{\raggedright\arraybackslash}p{22.0mm}>{\raggedright\arraybackslash}p{31.188mm}>{\raggedright\arraybackslash}p{31.188mm}>{\raggedright\arraybackslash}p{31.188mm}@{}}
\toprule
\textbf{Case} & \textbf{Masculine singular} & \textbf{Neuter singular} & \textbf{Feminine singular} \\
\midrule
Nominative & mǣre & mǣre & mǣru \\
Genitive & mǣres & mǣres & mǣrre \\
Dative & mǣrum & mǣrum & mǣrre \\
Accusative & mǣrne & mǣre & mǣre \\
\bottomrule
\end{tabular}
\par\endgroup
% course-panel:T67-P01:end
\par\addvspace{5pt}
% course-panel:T67-P02:begin
\par\begingroup
\fontsize{10.5}{13.2}\selectfont
\setlength{\LTleft}{0pt}\setlength{\LTright}{\fill}
\noindent
\begin{tabular}{@{}>{\raggedright\arraybackslash}p{22.0mm}>{\raggedright\arraybackslash}p{31.188mm}>{\raggedright\arraybackslash}p{31.188mm}>{\raggedright\arraybackslash}p{31.188mm}@{}}
\toprule
\textbf{Case} & \textbf{Masculine plural} & \textbf{Neuter plural} & \textbf{Feminine plural} \\
\midrule
Nominative & mǣre & mǣru & mǣra \\
Genitive & mǣrra & mǣrra & mǣrra \\
Dative & mǣrum & mǣrum & mǣrum \\
Accusative & mǣre & mǣru & mǣra \\
\bottomrule
\end{tabular}
\par\endgroup
% course-panel:T67-P02:end
% mdtable:T67:end

% teaching: NEW-U15-PAR-MAERE-WEAK-NGDA
\reserve{11\baselineskip}
\reserve{16\baselineskip}
% mdparagraph:P185:begin
\textbf{Mǣre: weak}
% mdparagraph:P185:end

% table: NEW-U15-PAR-MAERE-WEAK-NGDA-DISPLAY
% mdtable:T68:begin
% course-table-style: 2026-10
% course-panel:T68-P01:begin
\par\begingroup
\fontsize{10.5}{13.2}\selectfont
\setlength{\LTleft}{0pt}\setlength{\LTright}{\fill}
\noindent
\begin{tabular}{@{}>{\raggedright\arraybackslash}p{22.0mm}>{\raggedright\arraybackslash}p{31.188mm}>{\raggedright\arraybackslash}p{31.188mm}>{\raggedright\arraybackslash}p{31.188mm}@{}}
\toprule
\textbf{Case} & \textbf{Masculine singular} & \textbf{Neuter singular} & \textbf{Feminine singular} \\
\midrule
Nominative & mǣra & mǣre & mǣre \\
Genitive & mǣran & mǣran & mǣran \\
Dative & mǣran & mǣran & mǣran \\
Accusative & mǣran & mǣre & mǣran \\
\bottomrule
\end{tabular}
\par\endgroup
% course-panel:T68-P01:end
\par\addvspace{5pt}
% course-panel:T68-P02:begin
\par\begingroup
\fontsize{10.5}{13.2}\selectfont
\setlength{\LTleft}{0pt}\setlength{\LTright}{\fill}
\noindent
\begin{tabular}{@{}>{\raggedright\arraybackslash}p{22.0mm}>{\raggedright\arraybackslash}p{31.188mm}>{\raggedright\arraybackslash}p{31.188mm}>{\raggedright\arraybackslash}p{31.188mm}@{}}
\toprule
\textbf{Case} & \textbf{Masculine plural} & \textbf{Neuter plural} & \textbf{Feminine plural} \\
\midrule
Nominative & mǣran & mǣran & mǣran \\
Genitive & mǣrra & mǣrra & mǣrra \\
Dative & mǣrum & mǣrum & mǣrum \\
Accusative & mǣran & mǣran & mǣran \\
\bottomrule
\end{tabular}
\par\endgroup
% course-panel:T68-P02:end
% mdtable:T68:end

% teaching: NEW-U15-VOC-UNGETHANCFULL
% mdparagraph:P186:begin
\textbf{ungeðancfull} — unthankful.
% mdparagraph:P186:end

% mdparagraph:P187:begin
Ungeðancfull retains ll in the bare form and before vowel-initial endings, as in ungeðancfulle and ungeðancfullum. Before consonantal endings the selected forms have one l: ungeðancfulne, ungeðancfulre, ungeðancfulra.
% mdparagraph:P187:end

% teaching: NEW-U15-PAR-UNGETHANCFULL-STRONG-NGDA
\reserve{11\baselineskip}
\reserve{16\baselineskip}
% mdparagraph:P188:begin
\textbf{Ungeðancfull: strong}
% mdparagraph:P188:end

% table: NEW-U15-PAR-UNGETHANCFULL-STRONG-NGDA-DISPLAY
% mdtable:T69:begin
% course-table-style: 2026-10
% course-panel:T69-P01:begin
\par\begingroup
\fontsize{10.5}{13.2}\selectfont
\setlength{\LTleft}{0pt}\setlength{\LTright}{\fill}
\noindent
\begin{tabular}{@{}>{\raggedright\arraybackslash}p{22.0mm}>{\raggedright\arraybackslash}p{31.188mm}>{\raggedright\arraybackslash}p{31.188mm}>{\raggedright\arraybackslash}p{31.188mm}@{}}
\toprule
\textbf{Case} & \textbf{Masculine singular} & \textbf{Neuter singular} & \textbf{Feminine singular} \\
\midrule
Nominative & ungeðancfull & ungeðancfull & ungeðancfull \\
Genitive & ungeðancfulles & ungeðancfulles & ungeðancfulre \\
Dative & ungeðancfullum & ungeðancfullum & ungeðancfulre \\
Accusative & ungeðancfulne & ungeðancfull & ungeðancfulle \\
\bottomrule
\end{tabular}
\par\endgroup
% course-panel:T69-P01:end
\par\addvspace{5pt}
% course-panel:T69-P02:begin
\par\begingroup
\fontsize{10.5}{13.2}\selectfont
\setlength{\LTleft}{0pt}\setlength{\LTright}{\fill}
\noindent
\begin{tabular}{@{}>{\raggedright\arraybackslash}p{22.0mm}>{\raggedright\arraybackslash}p{31.188mm}>{\raggedright\arraybackslash}p{31.188mm}>{\raggedright\arraybackslash}p{31.188mm}@{}}
\toprule
\textbf{Case} & \textbf{Masculine plural} & \textbf{Neuter plural} & \textbf{Feminine plural} \\
\midrule
Nominative & ungeðancfulle & ungeðancfull & ungeðancfulla \\
Genitive & ungeðancfulra & ungeðancfulra & ungeðancfulra \\
Dative & ungeðancfullum & ungeðancfullum & ungeðancfullum \\
Accusative & ungeðancfulle & ungeðancfull & ungeðancfulla \\
\bottomrule
\end{tabular}
\par\endgroup
% course-panel:T69-P02:end
% mdtable:T69:end

% teaching: NEW-U15-PAR-UNGETHANCFULL-WEAK-NGDA
\reserve{11\baselineskip}
\reserve{16\baselineskip}
% mdparagraph:P189:begin
\textbf{Ungeðancfull: weak}
% mdparagraph:P189:end

% table: NEW-U15-PAR-UNGETHANCFULL-WEAK-NGDA-DISPLAY
% mdtable:T70:begin
% course-table-style: 2026-10
% course-panel:T70-P01:begin
\par\begingroup
\fontsize{10.5}{13.2}\selectfont
\setlength{\LTleft}{0pt}\setlength{\LTright}{\fill}
\noindent
\begin{tabular}{@{}>{\raggedright\arraybackslash}p{22.0mm}>{\raggedright\arraybackslash}p{31.188mm}>{\raggedright\arraybackslash}p{31.188mm}>{\raggedright\arraybackslash}p{31.188mm}@{}}
\toprule
\textbf{Case} & \textbf{Masculine singular} & \textbf{Neuter singular} & \textbf{Feminine singular} \\
\midrule
Nominative & ungeðancfulla & ungeðancfulle & ungeðancfulle \\
Genitive & ungeðancfullan & ungeðancfullan & ungeðancfullan \\
Dative & ungeðancfullan & ungeðancfullan & ungeðancfullan \\
Accusative & ungeðancfullan & ungeðancfulle & ungeðancfullan \\
\bottomrule
\end{tabular}
\par\endgroup
% course-panel:T70-P01:end
\par\addvspace{5pt}
% course-panel:T70-P02:begin
\par\begingroup
\fontsize{10.5}{13.2}\selectfont
\setlength{\LTleft}{0pt}\setlength{\LTright}{\fill}
\noindent
\begin{tabular}{@{}>{\raggedright\arraybackslash}p{22.0mm}>{\raggedright\arraybackslash}p{31.188mm}>{\raggedright\arraybackslash}p{31.188mm}>{\raggedright\arraybackslash}p{31.188mm}@{}}
\toprule
\textbf{Case} & \textbf{Masculine plural} & \textbf{Neuter plural} & \textbf{Feminine plural} \\
\midrule
Nominative & ungeðancfullan & ungeðancfullan & ungeðancfullan \\
Genitive & ungeðancfulra & ungeðancfulra & ungeðancfulra \\
Dative & ungeðancfullum & ungeðancfullum & ungeðancfullum \\
Accusative & ungeðancfullan & ungeðancfullan & ungeðancfullan \\
\bottomrule
\end{tabular}
\par\endgroup
% course-panel:T70-P02:end
% mdtable:T70:end

\reserve{7\baselineskip}
% mdparagraph:P190:begin
\subsection*{Additional vocabulary}
% mdparagraph:P190:end

% mdparagraph:P191:begin
Retain the earlier inflections of the familiar words below. Their stated senses are needed here; the adverbs do not agree with a noun.
% mdparagraph:P191:end

% teaching: NEW-U15-VOC-IDEL
% mdparagraph:P192:begin
\textbf{īdel} — idle; not working.
% mdparagraph:P192:end

% teaching: NEW-U15-VOC-LYTEL-SUPERL
% mdparagraph:P193:begin
\textbf{lǣst} — least in rank or importance.
% mdparagraph:P193:end

% teaching: NEW-U15-VOC-GOD-ADJECTIVE
% mdparagraph:P194:begin
\textbf{gōd} — kind; benevolent.
% mdparagraph:P194:end

% teaching: NEW-U15-VOC-YFEL
% mdparagraph:P195:begin
\textbf{yfel} — evil; evil people.
% mdparagraph:P195:end

% teaching: NEW-U15-VOC-SLEAN
% mdparagraph:P196:begin
\textbf{slēan} — kill.
% mdparagraph:P196:end

% teaching: NEW-U15-RECORD-GELICE-U14-SOURCE-1
% teaching: NEW-U15-VOC-GELICE
% mdparagraph:P197:begin
\textbf{gelīce} — likewise; similarly.
% mdparagraph:P197:end

% teaching: NEW-U15-RECORD-EORNOSTLICE-INVARIANT
% teaching: NEW-U15-VOC-EORNOSTLICE
% mdparagraph:P198:begin
\textbf{eornostlice} — accordingly; thus.
% mdparagraph:P198:end

% teaching: NEW-U15-RECORD-THAERTO-EACAN-INVARIANT
% teaching: NEW-U15-VOC-THAERTO-EACAN
% mdparagraph:P199:begin
\textbf{ðǣr-tō-ēacan} — in addition to that; besides.
% mdparagraph:P199:end

% teaching: NEW-U15-VOC-DON
% mdparagraph:P200:begin
\textbf{dōn} — make; render, with agreeing object predicate.
% mdparagraph:P200:end

% teaching: NEW-U15-RECORD-UT-INVARIANT
% teaching: NEW-U15-VOC-UT
% mdparagraph:P201:begin
\textbf{ūt} — out.
% mdparagraph:P201:end

% teaching: NEW-U15-RECORD-YMBE-INVARIANT
% teaching: NEW-U15-VOC-YMBE
% mdparagraph:P202:begin
\textbf{ymbe} — about, of approximate time, with accusative. Use accusative for the approximate time, including an hour.
% mdparagraph:P202:end

% teaching: NEW-U15-VOC-WITH
% mdparagraph:P203:begin
\textbf{wið} — in exchange or return for, with dative.
% mdparagraph:P203:end

% teaching: NEW-U15-VOC-FOR
% mdparagraph:P204:begin
\textbf{for} — for; in exchange for, with dative.
% mdparagraph:P204:end

% teaching: NEW-U15-VOC-OFER
% mdparagraph:P205:begin
\textbf{ofer} — above in rank, with accusative; toward, in gōd ofer persons, with accusative.
% mdparagraph:P205:end

% teaching: NEW-U15-VOC-THURH
% mdparagraph:P206:begin
\textbf{ðurh} — by, naming an oath standard, with accusative. The spelling ðurh is supplied beside the earlier þurh.
% mdparagraph:P206:end

\reserve{7\baselineskip}
% mdparagraph:P207:begin
\subsection*{Government and clause relations}
% mdparagraph:P207:end

% teaching: NEW-U15-GRAM-U15-FOR-EXCHANGE
% teaching: NEW-U15-GRAM-U15-FORGIELDAN-PAYMENT
% teaching: NEW-U15-GRAM-U15-ONFON-PAYMENT-ACC
% teaching: NEW-U15-GRAM-U15-WID-EXCHANGE
% mdparagraph:P208:begin
\textbf{Payment and exchange.} Forgieldan takes the person paid in the dative and the thing paid in the accusative. Onfōn can take an accusative material payment, including pęning and substantive māre. This payment construction stands alongside the earlier dative use of “receive”. Wið with the dative gives payment or return; for with the dative gives equivalent exchange, as in an eye-for-eye formula. The first noun in that quoted formula is nominative.
% mdparagraph:P208:end

% teaching: NEW-U15-GRAM-U15-GELIC-OBJECT-PREDICATE
% mdparagraph:P209:begin
\textbf{Making equal.} Dōn takes an accusative object and an agreeing predicate gelīc. The person used as the comparison standard is dative. Distinguish the object’s agreement from the comparison person’s case.
% mdparagraph:P209:end

% teaching: NEW-U15-GRAM-U15-WENAN-CONTENT
% mdparagraph:P210:begin
\textbf{Content after wēnan.} Wēnan can introduce a ðæt clause giving the thought or expectation. In a past report, sceolden, the preterite subjunctive plural of sculan, may express an expected outcome, “that they would receive”, rather than a duty.
% mdparagraph:P210:end

% teaching: NEW-U15-GRAM-U15-OFER-KINDNESS-RANK
% mdparagraph:P211:begin
\textbf{Kindness and rank.} Gōd ofer with the accusative means “kind toward” people. The bare strong accusative plurals ungeðancfulle and yfele can denote “unthankful people” and “evil people” in this use. Ofer with an accusative teacher or master can also mean “above” in rank. Neither use requires physical movement across something.
% mdparagraph:P211:end

% teaching: NEW-U15-GRAM-U15-SWERIAN-THURH
% mdparagraph:P212:begin
\textbf{An oath.} Swęrian with ðurh plus accusative names what one swears by. When a negative command is followed by several ne ðurh phrases, each phrase shares that command; the repeated ne does not create several positive assertions.
% mdparagraph:P212:end

% teaching: NEW-U15-GRAM-U15-ALIESAN-SEPARATION
% teaching: NEW-U15-GRAM-U15-LAERAN-OBJECTS
% mdparagraph:P213:begin
\textbf{Teaching and deliverance.} Lǣran takes an accusative person or the accusative matter taught. Ālīes takes the person delivered in the accusative; fram with the dative names what the person is delivered from. The passive teaching expression in Cook 31.3 and the exact substantive dative yfele, “evil”, in 31.4 remain supplied there.
% mdparagraph:P213:end

\reserve{7\baselineskip}
% mdparagraph:P214:begin
\subsection*{Translated examples — new teaching material}
% mdparagraph:P214:end

% mdparagraph:P215:begin
These examples distinguish meanings and constructions; they are not exercises.
% mdparagraph:P215:end

% teaching: NEW-U15-EXAMPLE-01
% teaching: NEW-U15-GRAM-U15-NEGATIVE-CONCORD
% teaching: NEW-U15-PAR-NAN-STRONG-NGDA
% teaching: NEW-U15-VOC-NAN
% mdparagraph:P216:begin
\textbf{1.} Nān cyning ne gǣð.
% mdparagraph:P216:end

% mdparagraph:P217:begin
No king goes.
% mdparagraph:P217:end

% mdparagraph:P218:begin
The negative determiner and finite ne express one denial.
% mdparagraph:P218:end

% teaching: NEW-U15-EXAMPLE-02
% teaching: NEW-U15-GRAM-U15-NABBAN-INDICATIVE
% teaching: NEW-U15-PAR-NABBAN-IND
% teaching: NEW-U15-VOC-NABBAN
% mdparagraph:P219:begin
\textbf{2.} Se mann ne hæfð męte; se mann næfð męte.
% mdparagraph:P219:end

% mdparagraph:P220:begin
The man does not have food; the man does not have food.
% mdparagraph:P220:end

% mdparagraph:P221:begin
The two clauses have the same negative meaning. The second uses the selected contracted finite form.
% mdparagraph:P221:end

% teaching: NEW-U15-EXAMPLE-03
% mdparagraph:P222:begin
\textbf{3.} Hwā hæfð ðæt hūs?
% mdparagraph:P222:end

% mdparagraph:P223:begin
Who has that house?
% mdparagraph:P223:end

% mdparagraph:P224:begin
Hwā asks for the person; this is an ordinary direct question.
% mdparagraph:P224:end

% teaching: NEW-U15-EXAMPLE-04
% teaching: NEW-U15-GRAM-U15-HWA-INDEFINITE
% teaching: NEW-U15-VOC-HWA
% mdparagraph:P225:begin
\textbf{4.} Gif hwā gā, hē gesiehð ðæt hūs.
% mdparagraph:P225:end

% mdparagraph:P226:begin
If anyone goes, he will see that house.
% mdparagraph:P226:end

% mdparagraph:P227:begin
Hwā is indefinite in the contingent condition. Gā is third-singular present subjunctive here; the following present states the prospective result.
% mdparagraph:P227:end

% teaching: NEW-U15-EXAMPLE-05
% mdparagraph:P228:begin
\textbf{5.} Ǣlc mann hæfð sunu.
% mdparagraph:P228:end

% mdparagraph:P229:begin
Each man has a son.
% mdparagraph:P229:end

% mdparagraph:P230:begin
Singular mann and hæfð express distribution to each individual.
% mdparagraph:P230:end

% teaching: NEW-U15-EXAMPLE-06
% teaching: NEW-U15-GRAM-U15-QUANTITY-PARTITIVES
% teaching: NEW-U15-NOM-BOC
% teaching: NEW-U15-VOC-BOC
% mdparagraph:P231:begin
\textbf{6.} Se cyning ānum ðissa manna bōc sęlð.
% mdparagraph:P231:end

% mdparagraph:P232:begin
The king gives a book to one of these men.
% mdparagraph:P232:end

% mdparagraph:P233:begin
The selected recipient is dative ānum; the group is genitive ðissa manna. The book is the accusative thing given.
% mdparagraph:P233:end

% teaching: NEW-U15-EXAMPLE-07
% teaching: NEW-U15-GRAM-U15-PROPORTIONAL-SWA
% teaching: NEW-U15-VOC-SWA
% mdparagraph:P234:begin
\textbf{7.} Swā se cyning ieldra is, swā hē wīsra is.
% mdparagraph:P234:end

% mdparagraph:P235:begin
The older the king is, the wiser he is.
% mdparagraph:P235:end

% mdparagraph:P236:begin
The paired clauses correlate degrees. Each comparative agrees with the subject of its own clause.
% mdparagraph:P236:end

% teaching: NEW-U15-EXAMPLE-08
% teaching: NEW-U15-NOM-BOC
% teaching: NEW-U15-VOC-BOC
% mdparagraph:P237:begin
\textbf{8.} Ne sęle ðǣm cyninge ðā bōc.
% mdparagraph:P237:end

% mdparagraph:P238:begin
Do not give the book to the king.
% mdparagraph:P238:end

% mdparagraph:P239:begin
Ne negates a singular command. The king remains the dative recipient and the book the accusative thing.
% mdparagraph:P239:end

\reserve{7\baselineskip}
% mdparagraph:P240:begin
\section*{6. Practice — new teaching material}
% mdparagraph:P240:end

% mdparagraph:P241:begin
The form grids identify each requested position even when several answers have the same spelling. Give English explanations for the analysis questions and Old English for composition. In composition, the directions specify the construction being practised. A different grammatical expression may convey the meaning without meeting that particular direction. Keep all supplied help visible.
% mdparagraph:P241:end

\reserve{7\baselineskip}
% mdparagraph:P242:begin
\subsection*{1. Negative propositions and complete indicative systems}
% mdparagraph:P242:end

\reserve{5\baselineskip}
\par\begin{minipage}{\linewidth}\raggedright
% mdparagraph:P243:begin
\textbf{F01}
% mdparagraph:P243:end

% item: NEW-U15-F01
% mdparagraph:P244:begin
Give the complete present and preterite indicative negative be-series for I, you singular, he, we, you plural and they. Use the selected contracted forms where the table gives them; retain separate ne where it does not. Explain why these are twelve logical positions even where forms coincide.
% mdparagraph:P244:end

% mdparagraph:P245:begin
\textbf{Response grid: wesan}
% mdparagraph:P245:end

% output: NEW-U15-F01-O001
% output: NEW-U15-F01-O002
% output: NEW-U15-F01-O003
% output: NEW-U15-F01-O004
% output: NEW-U15-F01-O005
% output: NEW-U15-F01-O006
% output: NEW-U15-F01-O007
% output: NEW-U15-F01-O008
% output: NEW-U15-F01-O009
% output: NEW-U15-F01-O010
% output: NEW-U15-F01-O011
% output: NEW-U15-F01-O012
% mdtable:T71:begin
% course-table-style: 2026-10
% course-panel:T71-P01:begin
\par\begingroup
\fontsize{10.5}{13.2}\selectfont
\setlength{\LTleft}{0pt}\setlength{\LTright}{\fill}
\noindent
\begin{tabular}{@{}>{\raggedright\arraybackslash}p{39.459mm}>{\raggedright\arraybackslash}p{39.459mm}>{\raggedright\arraybackslash}p{39.459mm}@{}}
\toprule
\textbf{Person and number} & \textbf{Present} & \textbf{Preterite} \\
\midrule
1 sg. & (1) \_\_\_\_ & (7) \_\_\_\_ \\
2 sg. & (2) \_\_\_\_ & (8) \_\_\_\_ \\
3 sg. & (3) \_\_\_\_ & (9) \_\_\_\_ \\
1 pl. & (4) \_\_\_\_ & (10) \_\_\_\_ \\
2 pl. & (5) \_\_\_\_ & (11) \_\_\_\_ \\
3 pl. & (6) \_\_\_\_ & (12) \_\_\_\_ \\
\bottomrule
\end{tabular}
\par\endgroup
% course-panel:T71-P01:end
% mdtable:T71:end
% output: NEW-U15-F01-O013
% mdparagraph:P246:begin
\textbf{Then explain:} (13) why twelve positions remain distinct; where ne stays separate.
% mdparagraph:P246:end

\end{minipage}\par\addvspace{4pt}

\par\begin{minipage}{\linewidth}\raggedright
% mdparagraph:P247:begin
\textbf{F02}
% mdparagraph:P247:end

% item: NEW-U15-F02
% mdparagraph:P248:begin
Give the complete present and preterite indicative of nabban for the six person/\allowbreak{}number positions. Interpret the first-person singular and third-person plural in each tense as denial of possession.
% mdparagraph:P248:end

% mdparagraph:P249:begin
\textbf{Response grid: nabban}
% mdparagraph:P249:end

% output: NEW-U15-F02-O001
% output: NEW-U15-F02-O002
% output: NEW-U15-F02-O003
% output: NEW-U15-F02-O004
% output: NEW-U15-F02-O005
% output: NEW-U15-F02-O006
% output: NEW-U15-F02-O007
% output: NEW-U15-F02-O008
% output: NEW-U15-F02-O009
% output: NEW-U15-F02-O010
% output: NEW-U15-F02-O011
% output: NEW-U15-F02-O012
% mdtable:T72:begin
% course-table-style: 2026-10
% course-panel:T72-P01:begin
\par\begingroup
\fontsize{10.5}{13.2}\selectfont
\setlength{\LTleft}{0pt}\setlength{\LTright}{\fill}
\noindent
\begin{tabular}{@{}>{\raggedright\arraybackslash}p{39.459mm}>{\raggedright\arraybackslash}p{39.459mm}>{\raggedright\arraybackslash}p{39.459mm}@{}}
\toprule
\textbf{Person and number} & \textbf{Present} & \textbf{Preterite} \\
\midrule
1 sg. & (1) \_\_\_\_ & (7) \_\_\_\_ \\
2 sg. & (2) \_\_\_\_ & (8) \_\_\_\_ \\
3 sg. & (3) \_\_\_\_ & (9) \_\_\_\_ \\
1 pl. & (4) \_\_\_\_ & (10) \_\_\_\_ \\
2 pl. & (5) \_\_\_\_ & (11) \_\_\_\_ \\
3 pl. & (6) \_\_\_\_ & (12) \_\_\_\_ \\
\bottomrule
\end{tabular}
\par\endgroup
% course-panel:T72-P01:end
% mdtable:T72:end
% output: NEW-U15-F02-O013
% output: NEW-U15-F02-O014
% output: NEW-U15-F02-O015
% output: NEW-U15-F02-O016
% mdparagraph:P250:begin
\textbf{Then explain:} (13) meaning of present first singular; (14) meaning of present third plural; (15) meaning of preterite first singular; (16) meaning of preterite third plural.
% mdparagraph:P250:end

\end{minipage}\par\addvspace{4pt}

\par\begin{minipage}{\linewidth}\raggedright
% mdparagraph:P251:begin
\textbf{F03}
% mdparagraph:P251:end

% item: NEW-U15-F03
% mdparagraph:P252:begin
Give the complete present and preterite indicative of the selected negative willan system. Interpret the first-person singular and second-person singular of each tense as refusal to go.
% mdparagraph:P252:end

% mdparagraph:P253:begin
\textbf{Response grid: willan}
% mdparagraph:P253:end

% output: NEW-U15-F03-O001
% output: NEW-U15-F03-O002
% output: NEW-U15-F03-O003
% output: NEW-U15-F03-O004
% output: NEW-U15-F03-O005
% output: NEW-U15-F03-O006
% output: NEW-U15-F03-O007
% output: NEW-U15-F03-O008
% output: NEW-U15-F03-O009
% output: NEW-U15-F03-O010
% output: NEW-U15-F03-O011
% output: NEW-U15-F03-O012
% mdtable:T73:begin
% course-table-style: 2026-10
% course-panel:T73-P01:begin
\par\begingroup
\fontsize{10.5}{13.2}\selectfont
\setlength{\LTleft}{0pt}\setlength{\LTright}{\fill}
\noindent
\begin{tabular}{@{}>{\raggedright\arraybackslash}p{39.459mm}>{\raggedright\arraybackslash}p{39.459mm}>{\raggedright\arraybackslash}p{39.459mm}@{}}
\toprule
\textbf{Person and number} & \textbf{Present} & \textbf{Preterite} \\
\midrule
1 sg. & (1) \_\_\_\_ & (7) \_\_\_\_ \\
2 sg. & (2) \_\_\_\_ & (8) \_\_\_\_ \\
3 sg. & (3) \_\_\_\_ & (9) \_\_\_\_ \\
1 pl. & (4) \_\_\_\_ & (10) \_\_\_\_ \\
2 pl. & (5) \_\_\_\_ & (11) \_\_\_\_ \\
3 pl. & (6) \_\_\_\_ & (12) \_\_\_\_ \\
\bottomrule
\end{tabular}
\par\endgroup
% course-panel:T73-P01:end
% mdtable:T73:end
% output: NEW-U15-F03-O013
% output: NEW-U15-F03-O014
% output: NEW-U15-F03-O015
% output: NEW-U15-F03-O016
% mdparagraph:P254:begin
\textbf{Then explain:} (13) meaning of present first singular; (14) meaning of present second singular; (15) meaning of preterite first singular; (16) meaning of preterite second singular.
% mdparagraph:P254:end

\end{minipage}\par\addvspace{4pt}

\par\begin{minipage}{\linewidth}\raggedright
% mdparagraph:P255:begin
\textbf{A01}
% mdparagraph:P255:end

% item: NEW-U15-A01
% mdparagraph:P256:begin
Interpret the three clauses. Identify every negative word and state the single proposition denied in each; do not turn the accumulated negatives into an affirmation.
% mdparagraph:P256:end

% mdparagraph:P257:begin
(1) Nān mann ne cymð.
% mdparagraph:P257:end

% mdparagraph:P258:begin
(2) Nǣfre ne cōm se cyning.
% mdparagraph:P258:end

% mdparagraph:P259:begin
(3) Nān mann nāwiht ne hæfð.
% mdparagraph:P259:end

% output: NEW-U15-A01-O001
% output: NEW-U15-A01-O002
% output: NEW-U15-A01-O003
% output: NEW-U15-A01-O004
% output: NEW-U15-A01-O005
% output: NEW-U15-A01-O006
% mdparagraph:P260:begin
\textbf{Responses, in order:} (1) clause 1: negative words; (2) clause 1: proposition denied; (3) clause 2: negative words; (4) clause 2: proposition denied; (5) clause 3: negative words; (6) clause 3: proposition denied.
% mdparagraph:P260:end

\end{minipage}\par\addvspace{4pt}

\par\begin{minipage}{\linewidth}\raggedright
% mdparagraph:P261:begin
\textbf{A02}
% mdparagraph:P261:end

% item: NEW-U15-A02
% mdparagraph:P262:begin
Compare the two clauses. Say whether the verb or one contrasted participant is negated, and identify the contrasting participant.
% mdparagraph:P262:end

% mdparagraph:P263:begin
(1) Ic self gā, and nā ōðer.
% mdparagraph:P263:end

% mdparagraph:P264:begin
(2) Ic ne gā.
% mdparagraph:P264:end

% output: NEW-U15-A02-O001
% output: NEW-U15-A02-O002
% output: NEW-U15-A02-O003
% mdparagraph:P265:begin
\textbf{Responses, in order:} (1) first clause: scope of negation; (2) first clause: contrasting participant; (3) second clause: scope of negation.
% mdparagraph:P265:end

\end{minipage}\par\addvspace{4pt}

\par\begin{minipage}{\linewidth}\raggedright
% mdparagraph:P266:begin
\textbf{C02}
% mdparagraph:P266:end

% item: NEW-U15-C02
% mdparagraph:P267:begin
Translate: I did not have the book, and I did not wish to go. Use the selected contracted negative preterites of have and will.
% mdparagraph:P267:end

% output: NEW-U15-C02-O001
\end{minipage}\par\addvspace{4pt}

\par\begin{minipage}{\linewidth}\raggedright
% mdparagraph:P268:begin
\textbf{C14}
% mdparagraph:P268:end

% item: NEW-U15-C14
% mdparagraph:P269:begin
Translate: He himself goes, and not another. The referent is a man. Use self and constituent nā.
% mdparagraph:P269:end

% output: NEW-U15-C14-O001
\end{minipage}\par\addvspace{4pt}

\reserve{7\baselineskip}
\par\begin{minipage}{\linewidth}\raggedright
% mdparagraph:P270:begin
\subsection*{2. Negative, indefinite and distributive reference}
% mdparagraph:P270:end


% mdparagraph:P271:begin
\textbf{F04}
% mdparagraph:P271:end

% item: NEW-U15-F04
% mdparagraph:P272:begin
Give all nominative, genitive, dative and accusative forms of nān, singular and plural, in masculine, feminine and neuter. Explain why a form such as nānum needs its noun or sentence to identify the particular case/\allowbreak{}number/\allowbreak{}gender.
% mdparagraph:P272:end

% mdparagraph:P273:begin
\textbf{Response grid: nān}
% mdparagraph:P273:end

% output: NEW-U15-F04-O001
% output: NEW-U15-F04-O002
% output: NEW-U15-F04-O003
% output: NEW-U15-F04-O004
% output: NEW-U15-F04-O005
% output: NEW-U15-F04-O006
% output: NEW-U15-F04-O007
% output: NEW-U15-F04-O008
% output: NEW-U15-F04-O009
% output: NEW-U15-F04-O010
% output: NEW-U15-F04-O011
% output: NEW-U15-F04-O012
% output: NEW-U15-F04-O013
% output: NEW-U15-F04-O014
% output: NEW-U15-F04-O015
% output: NEW-U15-F04-O016
% output: NEW-U15-F04-O017
% output: NEW-U15-F04-O018
% output: NEW-U15-F04-O019
% output: NEW-U15-F04-O020
% output: NEW-U15-F04-O021
% output: NEW-U15-F04-O022
% output: NEW-U15-F04-O023
% output: NEW-U15-F04-O024
% mdtable:T74:begin
% course-table-style: 2026-10
% course-panel:T74-P01:begin
\par\begingroup
\fontsize{10.5}{13.2}\selectfont
\setlength{\LTleft}{0pt}\setlength{\LTright}{\fill}
\noindent
\begin{tabular}{@{}>{\raggedright\arraybackslash}p{22.0mm}>{\raggedright\arraybackslash}p{31.188mm}>{\raggedright\arraybackslash}p{31.188mm}>{\raggedright\arraybackslash}p{31.188mm}@{}}
\toprule
\textbf{Case} & \textbf{Masculine singular} & \textbf{Neuter singular} & \textbf{Feminine singular} \\
\midrule
Nominative & (1) \_\_\_\_ & (17) \_\_\_\_ & (9) \_\_\_\_ \\
Genitive & (2) \_\_\_\_ & (18) \_\_\_\_ & (10) \_\_\_\_ \\
Dative & (3) \_\_\_\_ & (19) \_\_\_\_ & (11) \_\_\_\_ \\
Accusative & (4) \_\_\_\_ & (20) \_\_\_\_ & (12) \_\_\_\_ \\
\bottomrule
\end{tabular}
\par\endgroup
% course-panel:T74-P01:end
\par\addvspace{5pt}
% course-panel:T74-P02:begin
\par\begingroup
\fontsize{10.5}{13.2}\selectfont
\setlength{\LTleft}{0pt}\setlength{\LTright}{\fill}
\noindent
\begin{tabular}{@{}>{\raggedright\arraybackslash}p{22.0mm}>{\raggedright\arraybackslash}p{31.188mm}>{\raggedright\arraybackslash}p{31.188mm}>{\raggedright\arraybackslash}p{31.188mm}@{}}
\toprule
\textbf{Case} & \textbf{Masculine plural} & \textbf{Neuter plural} & \textbf{Feminine plural} \\
\midrule
Nominative & (5) \_\_\_\_ & (21) \_\_\_\_ & (13) \_\_\_\_ \\
Genitive & (6) \_\_\_\_ & (22) \_\_\_\_ & (14) \_\_\_\_ \\
Dative & (7) \_\_\_\_ & (23) \_\_\_\_ & (15) \_\_\_\_ \\
Accusative & (8) \_\_\_\_ & (24) \_\_\_\_ & (16) \_\_\_\_ \\
\bottomrule
\end{tabular}
\par\endgroup
% course-panel:T74-P02:end
% mdtable:T74:end

% output: NEW-U15-F04-O025
% mdparagraph:P274:begin
\textbf{Then explain:} (25) how the sentence identifies the form’s case, number and gender.
% mdparagraph:P274:end

\end{minipage}\par\addvspace{4pt}

\par\begin{minipage}{\linewidth}\raggedright
% mdparagraph:P275:begin
\textbf{F05}
% mdparagraph:P275:end

% item: NEW-U15-F05
% mdparagraph:P276:begin
Give all nominative, genitive, dative and accusative singular forms of gehwā/\allowbreak{}gehwæt, distinguishing personal and neuter reference. Explain why no plural column is required for this system.
% mdparagraph:P276:end

% mdparagraph:P277:begin
\textbf{Response grid: gehwā /\allowbreak{} gehwæt}
% mdparagraph:P277:end

% output: NEW-U15-F05-O001
% output: NEW-U15-F05-O002
% output: NEW-U15-F05-O003
% output: NEW-U15-F05-O004
% output: NEW-U15-F05-O005
% output: NEW-U15-F05-O006
% output: NEW-U15-F05-O007
% output: NEW-U15-F05-O008
% mdtable:T75:begin
% course-table-style: 2026-10
% course-panel:T75-P01:begin
\par\begingroup
\fontsize{10.5}{13.2}\selectfont
\setlength{\LTleft}{0pt}\setlength{\LTright}{\fill}
\noindent
\begin{tabular}{@{}>{\raggedright\arraybackslash}p{39.459mm}>{\raggedright\arraybackslash}p{39.459mm}>{\raggedright\arraybackslash}p{39.459mm}@{}}
\toprule
\textbf{Case} & \textbf{Neuter singular} & \textbf{Common singular} \\
\midrule
Nominative & (5) \_\_\_\_ & (1) \_\_\_\_ \\
Genitive & (6) \_\_\_\_ & (2) \_\_\_\_ \\
Dative & (7) \_\_\_\_ & (3) \_\_\_\_ \\
Accusative & (8) \_\_\_\_ & (4) \_\_\_\_ \\
\bottomrule
\end{tabular}
\par\endgroup
% course-panel:T75-P01:end
% mdtable:T75:end
% output: NEW-U15-F05-O009
% mdparagraph:P278:begin
\textbf{Then explain:} (9) reason no plural column is needed.
% mdparagraph:P278:end

\end{minipage}\par\addvspace{4pt}

\par\begin{minipage}{\linewidth}\raggedright
% mdparagraph:P279:begin
\textbf{A03}
% mdparagraph:P279:end

% item: NEW-U15-A03
% mdparagraph:P280:begin
For each clause identify whether hwā is interrogative, indefinite, or part of a free-choice expression; give the clause meaning and identify its subject.
% mdparagraph:P280:end

% mdparagraph:P281:begin
(1) Hwā cymð?
% mdparagraph:P281:end

% mdparagraph:P282:begin
(2) Gif hwā cume, ic him bōc sęlle.
% mdparagraph:P282:end

% mdparagraph:P283:begin
(3) Swā hwā swā cume, ic him bōc sęlle.
% mdparagraph:P283:end

% output: NEW-U15-A03-O001
% output: NEW-U15-A03-O002
% output: NEW-U15-A03-O003
% output: NEW-U15-A03-O004
% output: NEW-U15-A03-O005
% output: NEW-U15-A03-O006
% output: NEW-U15-A03-O007
% output: NEW-U15-A03-O008
% output: NEW-U15-A03-O009
% mdparagraph:P284:begin
\textbf{Responses, in order:} (1) clause 1: function of hwā; (2) clause 1: meaning; (3) clause 1: subject; (4) clause 2: function of hwā; (5) clause 2: meaning; (6) clause 2: subject; (7) clause 3: function of hwā; (8) clause 3: meaning; (9) clause 3: subject.
% mdparagraph:P284:end

\end{minipage}\par\addvspace{4pt}

\par\begin{minipage}{\linewidth}\raggedright
% mdparagraph:P285:begin
\textbf{A04}
% mdparagraph:P285:end

% item: NEW-U15-A04
% mdparagraph:P286:begin
Compare the two statements. Explain each versus all, identify the head of each subject, and account for the finite-verb number.
% mdparagraph:P286:end

% mdparagraph:P287:begin
(1) Ǣlc mann hæfð ānne pęning.
% mdparagraph:P287:end

% mdparagraph:P288:begin
(2) Ealle ðā męnn habbað męte.
% mdparagraph:P288:end

% output: NEW-U15-A04-O001
% output: NEW-U15-A04-O002
% output: NEW-U15-A04-O003
% output: NEW-U15-A04-O004
% mdparagraph:P289:begin
\textbf{Responses, in order:} (1) clause 1: quantifier meaning; (2) clause 1: subject head and verb agreement; (3) clause 2: quantifier meaning; (4) clause 2: subject head and verb agreement.
% mdparagraph:P289:end

\end{minipage}\par\addvspace{4pt}

\par\begin{minipage}{\linewidth}\raggedright
% mdparagraph:P290:begin
\textbf{A06}
% mdparagraph:P290:end

% item: NEW-U15-A06
% mdparagraph:P291:begin
Compare the two uses of man/\allowbreak{}mann. Identify the subject and finite-verb number, and explain why the impersonal use does not introduce a new plural pronoun.
% mdparagraph:P291:end

% mdparagraph:P292:begin
(1) Man sceal gān.
% mdparagraph:P292:end

% mdparagraph:P293:begin
(2) Se mann gǣð.
% mdparagraph:P293:end

% output: NEW-U15-A06-O001
% output: NEW-U15-A06-O002
% output: NEW-U15-A06-O003
% output: NEW-U15-A06-O004
% output: NEW-U15-A06-O005
% mdparagraph:P294:begin
\textbf{Responses, in order:} (1) clause 1: subject and reference; (2) clause 1: verb number; (3) clause 2: subject and reference; (4) clause 2: verb number; (5) reason this does not supply a new plural pronoun.
% mdparagraph:P294:end

\end{minipage}\par\addvspace{4pt}

\par\begin{minipage}{\linewidth}\raggedright
% mdparagraph:P295:begin
\textbf{A10}
% mdparagraph:P295:end

% item: NEW-U15-A10
% mdparagraph:P296:begin
Using the supplied meanings, interpret each clause and say whether it asks a question, denies possession, or asserts possession. Explain how nāwiht and ne work together in the middle clause.
% mdparagraph:P296:end

% mdparagraph:P297:begin
\textbf{Modern help — available throughout this item.}
% mdparagraph:P297:end

\end{minipage}\par\addvspace{4pt}

% teaching: NEW-U15-HELP-A10-INDEFINITES
% teaching: NEW-U15-VOC-AWIHT
% teaching: NEW-U15-VOC-HWAETHWEGA
% teaching: NEW-U15-VOC-NAWIHT
% support: NEW-U15-HELP-A10-INDEFINITES
% mdparagraph:P298:begin
āwiht anything; nāwiht nothing; hwæthwega something. Use these meanings to interpret the three clauses.
% mdparagraph:P298:end

% mdparagraph:P299:begin
(1) Hæfð hē āwiht?
% mdparagraph:P299:end

% mdparagraph:P300:begin
(2) Hē nāwiht ne hæfð.
% mdparagraph:P300:end

% mdparagraph:P301:begin
(3) Hē hæfð hwæthwega.
% mdparagraph:P301:end

% output: NEW-U15-A10-O001
% output: NEW-U15-A10-O002
% output: NEW-U15-A10-O003
% output: NEW-U15-A10-O004
% mdparagraph:P302:begin
\textbf{Responses, in order:} (1) clause 1: meaning and statement type; (2) clause 2: meaning and statement type; (3) clause 3: meaning and statement type; (4) how the two negatives work together.
% mdparagraph:P302:end

\reserve{5\baselineskip}
\par\begin{minipage}{\linewidth}\raggedright
% mdparagraph:P303:begin
\textbf{A11}
% mdparagraph:P303:end

% item: NEW-U15-A11
% mdparagraph:P304:begin
The same words Ān mann cōm occur in two settings: first, a story introduces an unnamed visitor; second, an explicit count contrasts one man with two. Give the sense of ān in each setting, and explain why the nominative agreement is the same.
% mdparagraph:P304:end

% mdparagraph:P305:begin
(1) Ān mann cōm.
% mdparagraph:P305:end

% mdparagraph:P306:begin
(2) Ān mann cōm.
% mdparagraph:P306:end

% output: NEW-U15-A11-O001
% output: NEW-U15-A11-O002
% output: NEW-U15-A11-O003
% mdparagraph:P307:begin
\textbf{Responses, in order:} (1) ān in the first setting; (2) ān in the second setting; (3) reason for the same agreement.
% mdparagraph:P307:end

\end{minipage}\par\addvspace{4pt}

\par\begin{minipage}{\linewidth}\raggedright
% mdparagraph:P308:begin
\textbf{C01}
% mdparagraph:P308:end

% item: NEW-U15-C01
% mdparagraph:P309:begin
Translate: No man has a book; each has food. Use nān and personal gehwā.
% mdparagraph:P309:end

% output: NEW-U15-C01-O001
\end{minipage}\par\addvspace{4pt}

\par\begin{minipage}{\linewidth}\raggedright
% mdparagraph:P310:begin
\textbf{C03}
% mdparagraph:P310:end

% item: NEW-U15-C03
% mdparagraph:P311:begin
Translate: Has the king any book? He has no book. Use ǣnig in the question and nān with finite negation in the reply.
% mdparagraph:P311:end

% output: NEW-U15-C03-O001
\end{minipage}\par\addvspace{4pt}

\par\begin{minipage}{\linewidth}\raggedright
% mdparagraph:P312:begin
\textbf{C04}
% mdparagraph:P312:end

% item: NEW-U15-C04
% mdparagraph:P313:begin
Translate: If anyone has a book, I shall give him a penny. Use indefinite hwā. Treat the condition as a possibility and use the present subjunctive in the if-clause.
% mdparagraph:P313:end

% output: NEW-U15-C04-O001
\end{minipage}\par\addvspace{4pt}

\par\begin{minipage}{\linewidth}\raggedright
% mdparagraph:P314:begin
\textbf{C07}
% mdparagraph:P314:end

% item: NEW-U15-C07
% mdparagraph:P315:begin
Translate: One can go. Use impersonal man with magan.
% mdparagraph:P315:end

% output: NEW-U15-C07-O001
\end{minipage}\par\addvspace{4pt}

\par\begin{minipage}{\linewidth}\raggedright
% mdparagraph:P316:begin
\textbf{C08}
% mdparagraph:P316:end

% item: NEW-U15-C08
% mdparagraph:P317:begin
Translate: Each son has one book; all the sons have food. Use ǣlc for each and eall for all.
% mdparagraph:P317:end

% output: NEW-U15-C08-O001
\end{minipage}\par\addvspace{4pt}

\par\begin{minipage}{\linewidth}\raggedright
% mdparagraph:P318:begin
\textbf{C15}
% mdparagraph:P318:end

% item: NEW-U15-C15
% mdparagraph:P319:begin
Translate: I give nothing to anyone. Use nāwiht for the thing and nān for the recipient. Use finite ne with the two specified negative expressions.
% mdparagraph:P319:end

% output: NEW-U15-C15-O001
\end{minipage}\par\addvspace{4pt}

\par\begin{minipage}{\linewidth}\raggedright
% mdparagraph:P320:begin
\textbf{C16}
% mdparagraph:P320:end

% item: NEW-U15-C16
% mdparagraph:P321:begin
Translate: No man’s son comes. Use nān within the dependent genitive phrase, with finite ne.
% mdparagraph:P321:end

% output: NEW-U15-C16-O001
\end{minipage}\par\addvspace{4pt}

\par\begin{minipage}{\linewidth}\raggedright
% mdparagraph:P322:begin
\textbf{C17}
% mdparagraph:P322:end

% item: NEW-U15-C17
% mdparagraph:P323:begin
Translate: The king has no books. Use contracted negative possession and nān.
% mdparagraph:P323:end

% output: NEW-U15-C17-O001
\end{minipage}\par\addvspace{4pt}

\par\begin{minipage}{\linewidth}\raggedright
% mdparagraph:P324:begin
\textbf{C18}
% mdparagraph:P324:end

% item: NEW-U15-C18
% mdparagraph:P325:begin
Translate: The king sees each person; each person’s son is good. Use personal gehwā for both each-person phrases.
% mdparagraph:P325:end

% output: NEW-U15-C18-O001
\end{minipage}\par\addvspace{4pt}

\par\begin{minipage}{\linewidth}\raggedright
% mdparagraph:P326:begin
\textbf{C19}
% mdparagraph:P326:end

% item: NEW-U15-C19
% mdparagraph:P327:begin
Translate: A certain man came, and he had a book. Use indefinite ān to introduce the unnamed man.
% mdparagraph:P327:end

% output: NEW-U15-C19-O001
\end{minipage}\par\addvspace{4pt}

\reserve{7\baselineskip}
% mdparagraph:P328:begin
\subsection*{3. Free choice, case, parts and groups}
% mdparagraph:P328:end

\reserve{5\baselineskip}
\par\begin{minipage}{\linewidth}\raggedright
% mdparagraph:P329:begin
\textbf{A05}
% mdparagraph:P329:end

% item: NEW-U15-A05
% mdparagraph:P330:begin
Explain the contribution of ǣgðer ge … ge and identify the two coordinated constituents. State whether the statement chooses just one of them.
% mdparagraph:P330:end

% mdparagraph:P331:begin
(1) Ic lufie ǣgðer ge ðone cyning ge his sunu.
% mdparagraph:P331:end

% output: NEW-U15-A05-O001
% output: NEW-U15-A05-O002
% output: NEW-U15-A05-O003
% output: NEW-U15-A05-O004
% mdparagraph:P332:begin
\textbf{Responses, in order:} (1) meaning of the construction; (2) first coordinated constituent; (3) second coordinated constituent; (4) whether just one is selected.
% mdparagraph:P332:end

\end{minipage}\par\addvspace{4pt}

\par\begin{minipage}{\linewidth}\raggedright
% mdparagraph:P333:begin
\textbf{A08}
% mdparagraph:P333:end

% item: NEW-U15-A08
% mdparagraph:P334:begin
In Swā hwā swā cume, ic him bōc sęlle, identify the case required for the subject that comes and the case required for the recipient of the book. Explain how the explicit pronoun in the main clause keeps the two roles separate.
% mdparagraph:P334:end

% mdparagraph:P335:begin
(1) Swā hwā swā cume, ic him bōc sęlle.
% mdparagraph:P335:end

% output: NEW-U15-A08-O001
% output: NEW-U15-A08-O002
% output: NEW-U15-A08-O003
% mdparagraph:P336:begin
\textbf{Responses, in order:} (1) subject of cume and its case; (2) recipient of the book and its case; (3) function of the explicit main-clause pronoun.
% mdparagraph:P336:end

\end{minipage}\par\addvspace{4pt}

\par\begin{minipage}{\linewidth}\raggedright
% mdparagraph:P337:begin
\textbf{A09}
% mdparagraph:P337:end

% item: NEW-U15-A09
% mdparagraph:P338:begin
Treat ānum uncer as the recipient of a gift and uncer hūs as a subject phrase. Explain the difference between the personal genitive uncer in the first phrase and the agreeing possessive uncer in the second. Identify the role and case of ānum and the number of the group selected by uncer.
% mdparagraph:P338:end

% mdparagraph:P339:begin
(1) ānum uncer
% mdparagraph:P339:end

% mdparagraph:P340:begin
(2) uncer hūs
% mdparagraph:P340:end

% output: NEW-U15-A09-O001
% output: NEW-U15-A09-O002
% output: NEW-U15-A09-O003
% output: NEW-U15-A09-O004
% mdparagraph:P341:begin
\textbf{Responses, in order:} (1) first uncer: function; (2) second uncer: function; (3) ānum: role and case; (4) number of the group.
% mdparagraph:P341:end

\end{minipage}\par\addvspace{4pt}

\par\begin{minipage}{\linewidth}\raggedright
% mdparagraph:P342:begin
\textbf{C05}
% mdparagraph:P342:end

% item: NEW-U15-C05
% mdparagraph:P343:begin
Translate: Whoever goes, I shall give him a penny. Use swā hwā swā and the present subjunctive for the open choice of a possible person.
% mdparagraph:P343:end

% output: NEW-U15-C05-O001
\end{minipage}\par\addvspace{4pt}

\par\begin{minipage}{\linewidth}\raggedright
% mdparagraph:P344:begin
\textbf{C06}
% mdparagraph:P344:end

% item: NEW-U15-C06
% mdparagraph:P345:begin
Translate: Whichever of these two sons comes first, I will give him the book. Use swā hwæðer swā and the present subjunctive for the possible first arrival. Keep the subject of comes nominative and express the recipient again with a dative pronoun in the main clause.
% mdparagraph:P345:end

% output: NEW-U15-C06-O001
\end{minipage}\par\addvspace{4pt}

\par\begin{minipage}{\linewidth}\raggedright
% mdparagraph:P346:begin
\textbf{C09}
% mdparagraph:P346:end

% item: NEW-U15-C09
% mdparagraph:P347:begin
Translate: I see both the king and his son. Use ǣgðer ge … ge.
% mdparagraph:P347:end

% output: NEW-U15-C09-O001
\end{minipage}\par\addvspace{4pt}

\reserve{7\baselineskip}
% mdparagraph:P348:begin
\subsection*{4. Equality, preference and proportion}
% mdparagraph:P348:end

\reserve{5\baselineskip}
\par\begin{minipage}{\linewidth}\raggedright
% mdparagraph:P349:begin
\textbf{A07}
% mdparagraph:P349:end

% item: NEW-U15-A07
% mdparagraph:P350:begin
Distinguish the meanings of the three comparisons. Identify the comparison standard or correlated clause in each.
% mdparagraph:P350:end

% mdparagraph:P351:begin
(1) Ðēos bōc is swā gōd swā sēo bōc.
% mdparagraph:P351:end

% mdparagraph:P352:begin
(2) Ic lufie lēoht swīðor ðonne ðīestru.
% mdparagraph:P352:end

% mdparagraph:P353:begin
(3) Swā ðēos bōc māre is, swā hēo bętere is.
% mdparagraph:P353:end

% output: NEW-U15-A07-O001
% output: NEW-U15-A07-O002
% output: NEW-U15-A07-O003
% output: NEW-U15-A07-O004
% output: NEW-U15-A07-O005
% output: NEW-U15-A07-O006
% mdparagraph:P354:begin
\textbf{Responses, in order:} (1) comparison 1: meaning; (2) comparison 1: standard or related clause; (3) comparison 2: meaning; (4) comparison 2: standard or related clause; (5) comparison 3: meaning; (6) comparison 3: standard or related clause.
% mdparagraph:P354:end

\end{minipage}\par\addvspace{4pt}

\par\begin{minipage}{\linewidth}\raggedright
% mdparagraph:P355:begin
\textbf{C10}
% mdparagraph:P355:end

% item: NEW-U15-C10
% mdparagraph:P356:begin
Translate: The larger that house is, the better it is. Use proportional swā … swā and a pronoun referring back to the house.
% mdparagraph:P356:end

% output: NEW-U15-C10-O001
\end{minipage}\par\addvspace{4pt}

\par\begin{minipage}{\linewidth}\raggedright
% mdparagraph:P357:begin
\textbf{C11}
% mdparagraph:P357:end

% item: NEW-U15-C11
% mdparagraph:P358:begin
Translate: This book is as good as that book, but the king’s book is better. Use equal-degree swā … swā in the first clause.
% mdparagraph:P358:end

% output: NEW-U15-C11-O001
\end{minipage}\par\addvspace{4pt}

\par\begin{minipage}{\linewidth}\raggedright
% mdparagraph:P359:begin
\textbf{C12}
% mdparagraph:P359:end

% item: NEW-U15-C12
% mdparagraph:P360:begin
Translate: He loves food more than wine. Use preferential swīðor.
% mdparagraph:P360:end

% output: NEW-U15-C12-O001
\end{minipage}\par\addvspace{4pt}

\reserve{7\baselineskip}
\par\begin{minipage}{\linewidth}\raggedright
% mdparagraph:P361:begin
\subsection*{5. Noun systems and commands}
% mdparagraph:P361:end


% mdparagraph:P362:begin
\textbf{F06}
% mdparagraph:P362:end

% item: NEW-U15-F06
% mdparagraph:P363:begin
Give all nominative, genitive, dative and accusative singular/\allowbreak{}plural forms of bōc and tōð. Identify the forms that change the stem vowel and distinguish a syncretic surface from a uniquely identifiable case.
% mdparagraph:P363:end

% mdparagraph:P364:begin
\textbf{Response grid: bōc}
% mdparagraph:P364:end

% output: NEW-U15-F06-O001
% output: NEW-U15-F06-O002
% output: NEW-U15-F06-O003
% output: NEW-U15-F06-O004
% output: NEW-U15-F06-O005
% output: NEW-U15-F06-O006
% output: NEW-U15-F06-O007
% output: NEW-U15-F06-O008
% mdtable:T76:begin
\par\begingroup
\fontsize{10.5}{13.2}\selectfont
\setlength{\LTleft}{0pt}\setlength{\LTright}{\fill}
\noindent
\begin{tabular}{@{}>{\raggedright\arraybackslash}p{39.46mm}>{\raggedright\arraybackslash}p{39.46mm}>{\raggedright\arraybackslash}p{39.46mm}@{}}
\toprule
\textbf{Case} & \textbf{Singular} & \textbf{Plural} \\
\midrule
Nominative & (1) \_\_\_\_ & (5) \_\_\_\_ \\
Genitive & (2) \_\_\_\_ & (6) \_\_\_\_ \\
Dative & (3) \_\_\_\_ & (7) \_\_\_\_ \\
Accusative & (4) \_\_\_\_ & (8) \_\_\_\_ \\
\bottomrule
\end{tabular}
\par\endgroup

% mdtable:T76:end
% mdparagraph:P365:begin
\textbf{Response grid: tōð}
% mdparagraph:P365:end

% output: NEW-U15-F06-O009
% output: NEW-U15-F06-O010
% output: NEW-U15-F06-O011
% output: NEW-U15-F06-O012
% output: NEW-U15-F06-O013
% output: NEW-U15-F06-O014
% output: NEW-U15-F06-O015
% output: NEW-U15-F06-O016
% mdtable:T77:begin
\par\begingroup
\fontsize{10.5}{13.2}\selectfont
\setlength{\LTleft}{0pt}\setlength{\LTright}{\fill}
\noindent
\begin{tabular}{@{}>{\raggedright\arraybackslash}p{39.46mm}>{\raggedright\arraybackslash}p{39.46mm}>{\raggedright\arraybackslash}p{39.46mm}@{}}
\toprule
\textbf{Case} & \textbf{Singular} & \textbf{Plural} \\
\midrule
Nominative & (9) \_\_\_\_ & (13) \_\_\_\_ \\
Genitive & (10) \_\_\_\_ & (14) \_\_\_\_ \\
Dative & (11) \_\_\_\_ & (15) \_\_\_\_ \\
Accusative & (12) \_\_\_\_ & (16) \_\_\_\_ \\
\bottomrule
\end{tabular}
\par\endgroup

% mdtable:T77:end
% output: NEW-U15-F06-O017
% output: NEW-U15-F06-O018
% mdparagraph:P366:begin
\textbf{Then explain:} (17) positions with vowel changes; (18) forms whose cases need context.
% mdparagraph:P366:end

\end{minipage}\par\addvspace{4pt}

\par\begin{minipage}{\linewidth}\raggedright
% mdparagraph:P367:begin
\textbf{F07}
% mdparagraph:P367:end

% item: NEW-U15-F07
% mdparagraph:P368:begin
Give the complete nominative, genitive, dative and accusative singular/\allowbreak{}plural of sweostor. Compare the number of distinct surfaces with bōc; explain why the frequent surface sweostor still requires a case analysis in a sentence.
% mdparagraph:P368:end

% mdparagraph:P369:begin
\textbf{Response grid: sweostor}
% mdparagraph:P369:end

% output: NEW-U15-F07-O001
% output: NEW-U15-F07-O002
% output: NEW-U15-F07-O003
% output: NEW-U15-F07-O004
% output: NEW-U15-F07-O005
% output: NEW-U15-F07-O006
% output: NEW-U15-F07-O007
% output: NEW-U15-F07-O008
% mdtable:T78:begin
\par\begingroup
\fontsize{10.5}{13.2}\selectfont
\setlength{\LTleft}{0pt}\setlength{\LTright}{\fill}
\noindent
\begin{tabular}{@{}>{\raggedright\arraybackslash}p{39.46mm}>{\raggedright\arraybackslash}p{39.46mm}>{\raggedright\arraybackslash}p{39.46mm}@{}}
\toprule
\textbf{Case} & \textbf{Singular} & \textbf{Plural} \\
\midrule
Nominative & (1) \_\_\_\_ & (5) \_\_\_\_ \\
Genitive & (2) \_\_\_\_ & (6) \_\_\_\_ \\
Dative & (3) \_\_\_\_ & (7) \_\_\_\_ \\
Accusative & (4) \_\_\_\_ & (8) \_\_\_\_ \\
\bottomrule
\end{tabular}
\par\endgroup

% mdtable:T78:end
% output: NEW-U15-F07-O009
% output: NEW-U15-F07-O010
% mdparagraph:P370:begin
\textbf{Then explain:} (9) contrast with bōc; (10) reason a sentence still requires case analysis.
% mdparagraph:P370:end

\end{minipage}\par\addvspace{4pt}

\par\begin{minipage}{\linewidth}\raggedright
% mdparagraph:P371:begin
\textbf{F08}
% mdparagraph:P371:end

% item: NEW-U15-F08
% mdparagraph:P372:begin
Give nominative, genitive, dative and accusative singular of hǣtu and hǣlu. Use the selected forms in the tables. Explain why the plural positions are labelled unassigned and why this neither licenses a count meaning nor claims that no historical plural ever existed.
% mdparagraph:P372:end

% mdparagraph:P373:begin
\textbf{Response grid: hǣtu}
% mdparagraph:P373:end

% output: NEW-U15-F08-O001
% output: NEW-U15-F08-O002
% output: NEW-U15-F08-O003
% output: NEW-U15-F08-O004
% mdtable:T79:begin
\par\begingroup
\fontsize{10.5}{13.2}\selectfont
\setlength{\LTleft}{0pt}\setlength{\LTright}{\fill}
\noindent
\begin{tabular}{@{}>{\raggedright\arraybackslash}p{60.59mm}>{\raggedright\arraybackslash}p{60.59mm}@{}}
\toprule
\textbf{Case} & \textbf{Singular} \\
\midrule
Nominative & (1) \_\_\_\_ \\
Genitive & (2) \_\_\_\_ \\
Dative & (3) \_\_\_\_ \\
Accusative & (4) \_\_\_\_ \\
\bottomrule
\end{tabular}
\par\endgroup

% mdtable:T79:end
% mdparagraph:P374:begin
\textbf{Response grid: hǣlu}
% mdparagraph:P374:end

% output: NEW-U15-F08-O005
% output: NEW-U15-F08-O006
% output: NEW-U15-F08-O007
% output: NEW-U15-F08-O008
% mdtable:T80:begin
\par\begingroup
\fontsize{10.5}{13.2}\selectfont
\setlength{\LTleft}{0pt}\setlength{\LTright}{\fill}
\noindent
\begin{tabular}{@{}>{\raggedright\arraybackslash}p{60.59mm}>{\raggedright\arraybackslash}p{60.59mm}@{}}
\toprule
\textbf{Case} & \textbf{Singular} \\
\midrule
Nominative & (5) \_\_\_\_ \\
Genitive & (6) \_\_\_\_ \\
Dative & (7) \_\_\_\_ \\
Accusative & (8) \_\_\_\_ \\
\bottomrule
\end{tabular}
\par\endgroup

% mdtable:T80:end
% output: NEW-U15-F08-O009
% output: NEW-U15-F08-O010
% mdparagraph:P375:begin
\textbf{Then explain:} (9) selected singular patterns; (10) meaning and limits of the unassigned plural cells.
% mdparagraph:P375:end

\end{minipage}\par\addvspace{4pt}

\par\begin{minipage}{\linewidth}\raggedright
% mdparagraph:P376:begin
\textbf{F09}
% mdparagraph:P376:end

% item: NEW-U15-F09
% mdparagraph:P377:begin
Give the complete nominative, genitive, dative and accusative singular/\allowbreak{}plural of gerēfa and lēoðwyrhta. Explain why the word-formation of lēoðwyrhta does not change where its weak-noun ending belongs.
% mdparagraph:P377:end

% mdparagraph:P378:begin
\textbf{Response grid: gerēfa}
% mdparagraph:P378:end

% output: NEW-U15-F09-O001
% output: NEW-U15-F09-O002
% output: NEW-U15-F09-O003
% output: NEW-U15-F09-O004
% output: NEW-U15-F09-O005
% output: NEW-U15-F09-O006
% output: NEW-U15-F09-O007
% output: NEW-U15-F09-O008
% mdtable:T81:begin
\par\begingroup
\fontsize{10.5}{13.2}\selectfont
\setlength{\LTleft}{0pt}\setlength{\LTright}{\fill}
\noindent
\begin{tabular}{@{}>{\raggedright\arraybackslash}p{39.46mm}>{\raggedright\arraybackslash}p{39.46mm}>{\raggedright\arraybackslash}p{39.46mm}@{}}
\toprule
\textbf{Case} & \textbf{Singular} & \textbf{Plural} \\
\midrule
Nominative & (1) \_\_\_\_ & (5) \_\_\_\_ \\
Genitive & (2) \_\_\_\_ & (6) \_\_\_\_ \\
Dative & (3) \_\_\_\_ & (7) \_\_\_\_ \\
Accusative & (4) \_\_\_\_ & (8) \_\_\_\_ \\
\bottomrule
\end{tabular}
\par\endgroup

% mdtable:T81:end
% mdparagraph:P379:begin
\textbf{Response grid: lēoðwyrhta}
% mdparagraph:P379:end

% output: NEW-U15-F09-O009
% output: NEW-U15-F09-O010
% output: NEW-U15-F09-O011
% output: NEW-U15-F09-O012
% output: NEW-U15-F09-O013
% output: NEW-U15-F09-O014
% output: NEW-U15-F09-O015
% output: NEW-U15-F09-O016
% mdtable:T82:begin
\par\begingroup
\fontsize{10.5}{13.2}\selectfont
\setlength{\LTleft}{0pt}\setlength{\LTright}{\fill}
\noindent
\begin{tabular}{@{}>{\raggedright\arraybackslash}p{39.46mm}>{\raggedright\arraybackslash}p{39.46mm}>{\raggedright\arraybackslash}p{39.46mm}@{}}
\toprule
\textbf{Case} & \textbf{Singular} & \textbf{Plural} \\
\midrule
Nominative & (9) \_\_\_\_ & (13) \_\_\_\_ \\
Genitive & (10) \_\_\_\_ & (14) \_\_\_\_ \\
Dative & (11) \_\_\_\_ & (15) \_\_\_\_ \\
Accusative & (12) \_\_\_\_ & (16) \_\_\_\_ \\
\bottomrule
\end{tabular}
\par\endgroup

% mdtable:T82:end
% output: NEW-U15-F09-O017
% mdparagraph:P380:begin
\textbf{Then explain:} (17) relation of the compound’s structure to its inflection.
% mdparagraph:P380:end

\end{minipage}\par\addvspace{4pt}

\par\begin{minipage}{\linewidth}\raggedright
% mdparagraph:P381:begin
\textbf{A12}
% mdparagraph:P381:end

% item: NEW-U15-A12
% mdparagraph:P382:begin
Use the complete displayed hēafod and fōtsceamul tables. Identify the ending-bearing stem of hēafod, give its selected nominative/\allowbreak{}accusative plural, and contrast the stem retained throughout fōtsceamul. Explain why a general instruction to delete every medial vowel would give wrong forms.
% mdparagraph:P382:end

% mdparagraph:P383:begin
(1) hēafod
% mdparagraph:P383:end

% mdparagraph:P384:begin
(2) fōtsceamul
% mdparagraph:P384:end

% output: NEW-U15-A12-O001
% output: NEW-U15-A12-O002
% output: NEW-U15-A12-O003
% output: NEW-U15-A12-O004
% mdparagraph:P385:begin
\textbf{Responses, in order:} (1) hēafod: stem before endings; (2) hēafod: nominative/\allowbreak{}accusative plural; (3) fōtsceamul: stem before endings; (4) why deleting every medial vowel would fail.
% mdparagraph:P385:end

\end{minipage}\par\addvspace{4pt}

\par\begin{minipage}{\linewidth}\raggedright
% mdparagraph:P386:begin
\textbf{C13}
% mdparagraph:P386:end

% item: NEW-U15-C13
% mdparagraph:P387:begin
Translate: Do not go; give the book to the poor man. Address one person with a negative command followed by an affirmative command.
% mdparagraph:P387:end

% output: NEW-U15-C13-O001
\end{minipage}\par\addvspace{4pt}

\reserve{7\baselineskip}
% mdparagraph:P388:begin
\section*{7. Cook’s exercises}
% mdparagraph:P388:end

% mdparagraph:P389:begin
The prompts, parentheses, references and translation controls below are Cook’s original material. Translate every item into Old English. Modern help is separate and remains available; it supplies unusual forms or constructions without requiring their unaided recall. Numbers in Cook’s parentheses refer to his grammar or earlier exercises.\textsuperscript{7}
% mdparagraph:P389:end

\reserve{7\baselineskip}
% mdparagraph:P390:begin
\subsection*{Cook IX.28}
% mdparagraph:P390:end

\reserve{5\baselineskip}
\par\begin{minipage}{\linewidth}\raggedright
% mdparagraph:P391:begin
\textbf{Cook IX.28.1}
% mdparagraph:P391:end

% item: EX-IX-28-01
% mdparagraph:P392:begin
Every one that (87. d) forsakes for my name his house, or his brethren, or sisters (52. 3), or father, or mother, or wife, or children, or lands, with a hundredfold (om. it shall) be requited to him (verb last), and he shall have besides eternal life.
% mdparagraph:P392:end

% mdparagraph:P393:begin
\textbf{Modern help — available throughout this item.}
% mdparagraph:P393:end

\end{minipage}\par\addvspace{4pt}

% teaching: U15-COOK-NOM-28-01
% support: U15-COOK-NOM-28-01
% mdparagraph:P394:begin
Keep every alternative, including sisters, and both promised results. Sweostor has its own table; retain the distinct earlier brōðor and mōdor patterns.
% mdparagraph:P394:end

% teaching: NEW-U11-LOCAL-HELPS
% teaching: NEW-U15-LOCAL-HELPS
% teaching: NEW-U15-RECORD-HUNDTEONTIG-U15-SOURCE-DAT-PL
% teaching: NEW-U15-RECORD-HUNDTEONTIG-U15-SOURCE-SENSE
% teaching: U15-H-28-01
% support: U15-H-28-01
% mdparagraph:P395:begin
Use him mid hundtēontigum forgolden bið, “he will be repaid a hundredfold”, with the finite verb last. The initial ǣlc has its own relative clause and is resumed by dative him, the recipient in this supplied passive construction. Mid hundtēontigum means “with a hundredfold”. For mīnum naman means “for the sake of my name”. Keep this whole construction available as help.
% mdparagraph:P395:end

% output: EX-IX-28-01-O001
\reserve{5\baselineskip}
\par\begin{minipage}{\linewidth}\raggedright
% mdparagraph:P396:begin
\textbf{Cook IX.28.2}
% mdparagraph:P396:end

% item: EX-IX-28-02
% mdparagraph:P397:begin
He gave (sęllan) to each a penny in return for his day’s work.
% mdparagraph:P397:end

% output: EX-IX-28-02-O001
\end{minipage}\par\addvspace{4pt}

\par\begin{minipage}{\linewidth}\raggedright
% mdparagraph:P398:begin
\textbf{Cook IX.28.3}
% mdparagraph:P398:end

% item: EX-IX-28-03
% mdparagraph:P399:begin
He sent them into (on, 172) his vineyard.
% mdparagraph:P399:end

% output: EX-IX-28-03-O001
\end{minipage}\par\addvspace{4pt}

\par\begin{minipage}{\linewidth}\raggedright
% mdparagraph:P400:begin
\textbf{Cook IX.28.4}
% mdparagraph:P400:end

% item: EX-IX-28-04
% mdparagraph:P401:begin
He saw others in (om. the) public place standing (169) idle.
% mdparagraph:P401:end

% mdparagraph:P402:begin
\textbf{Modern help — available throughout this item.}
% mdparagraph:P402:end

\end{minipage}\par\addvspace{4pt}

% teaching: U15-COOK-NOM-28-04
% support: U15-COOK-NOM-28-04
% mdparagraph:P403:begin
Strǣt here means “public place”, with the earlier feminine declension. Cook’s reference (169) selects the earlier perception construction: use accusative ōðre with the plain infinitive standan. In 28.7, by contrast, the printed direction explicitly requests a present participle.
% mdparagraph:P403:end

% output: EX-IX-28-04-O001
\reserve{5\baselineskip}
\par\begin{minipage}{\linewidth}\raggedright
% mdparagraph:P404:begin
\textbf{Cook IX.28.5}
% mdparagraph:P404:end

% item: EX-IX-28-05
% mdparagraph:P405:begin
Then said he, Go ye also into my vineyard, and I will give you what (ðæt) is right.
% mdparagraph:P405:end

% output: EX-IX-28-05-O001
\end{minipage}\par\addvspace{4pt}

\par\begin{minipage}{\linewidth}\raggedright
% mdparagraph:P406:begin
\textbf{Cook IX.28.6}
% mdparagraph:P406:end

% item: EX-IX-28-06
% mdparagraph:P407:begin
Once more he went out about (ymbe, 172) the sixth and ninth hour, and did similarly.
% mdparagraph:P407:end

% mdparagraph:P408:begin
\textbf{Modern help — available throughout this item.}
% mdparagraph:P408:end

\end{minipage}\par\addvspace{4pt}

% teaching: U15-COOK-NOM-28-06
% support: U15-COOK-NOM-28-06
% mdparagraph:P409:begin
Tīd means “hour”. Temporal ymbe takes accusative for the approximate time.
% mdparagraph:P409:end

% teaching: NEW-U15-LOCAL-HELPS
% teaching: U15-H-28-06
% support: U15-H-28-06
% mdparagraph:P410:begin
Use supplied nigoðan, the feminine accusative singular weak form of “ninth”, with tīd. Siextan, “sixth”, is already familiar.
% mdparagraph:P410:end

% output: EX-IX-28-06-O001
\reserve{5\baselineskip}
\par\begin{minipage}{\linewidth}\raggedright
% mdparagraph:P411:begin
\textbf{Cook IX.28.7}
% mdparagraph:P411:end

% item: EX-IX-28-07
% mdparagraph:P412:begin
Then about the eleventh hour he went out, and found others standing (pres. part.).
% mdparagraph:P412:end

% mdparagraph:P413:begin
\textbf{Modern help — available throughout this item.}
% mdparagraph:P413:end

\end{minipage}\par\addvspace{4pt}

% teaching: NEW-U15-LOCAL-HELPS
% teaching: U15-H-28-07
% support: U15-H-28-07
% mdparagraph:P414:begin
Use the supplied expression ðā endleftan tīd, “the eleventh hour”, with the feminine accusative singular weak ordinal.
% mdparagraph:P414:end

% output: EX-IX-28-07-O001
\reserve{5\baselineskip}
\par\begin{minipage}{\linewidth}\raggedright
% mdparagraph:P415:begin
\textbf{Cook IX.28.8}
% mdparagraph:P415:end

% item: EX-IX-28-08
% mdparagraph:P416:begin
Then said he, Why stand ye here all day idle?
% mdparagraph:P416:end

% output: EX-IX-28-08-O001
\end{minipage}\par\addvspace{4pt}

\par\begin{minipage}{\linewidth}\raggedright
% mdparagraph:P417:begin
\textbf{Cook IX.28.9}
% mdparagraph:P417:end

% item: EX-IX-28-09
% mdparagraph:P418:begin
Then said they, Because no man (om. hath) hired (hȳrian, 118) us.
% mdparagraph:P418:end

% output: EX-IX-28-09-O001
\end{minipage}\par\addvspace{4pt}

\par\begin{minipage}{\linewidth}\raggedright
% mdparagraph:P419:begin
\textbf{Cook IX.28.10}
% mdparagraph:P419:end

% item: EX-IX-28-10
% mdparagraph:P420:begin
Truly, when even was come [when it was evening become], then said the lord (hlāford) of the vineyard to his steward (gerēfa, 53), Call (cleopian) the workmen, and pay them (164. a) their hire (mēd).
% mdparagraph:P420:end

% mdparagraph:P421:begin
\textbf{Modern help — available throughout this item.}
% mdparagraph:P421:end



% teaching: U15-COOK-NOM-28-10
% support: U15-COOK-NOM-28-10
% mdparagraph:P422:begin
Ǣfen means “evening”; wyrhta, “workman”; gerēfa, “steward”; wīngeard, “vineyard”; mēd here, “hire; wages”.
% mdparagraph:P422:end
\end{minipage}\par\addvspace{4pt}

% teaching: NEW-U14-LOCAL-HELPS
% teaching: NEW-U15-LOCAL-HELPS
% teaching: U15-H-28-10
% support: U15-H-28-10
% mdparagraph:P423:begin
Use ðā hit ǣfen geworden wæs, “when it had become evening”. Geworden is the supplied past participle of geweorðan. With cweðan, tō his gerēfan names the person addressed: tō takes dative.
% mdparagraph:P423:end

% output: EX-IX-28-10-O001
\reserve{5\baselineskip}
\par\begin{minipage}{\linewidth}\raggedright
% mdparagraph:P424:begin
\textbf{Cook IX.28.11}
% mdparagraph:P424:end

% item: EX-IX-28-11
% mdparagraph:P425:begin
Accordingly, when those came that came about the eleventh hour (verb last), then received they each his penny.
% mdparagraph:P425:end

% mdparagraph:P426:begin
\textbf{Modern help — available throughout this item.}
% mdparagraph:P426:end

\end{minipage}\par\addvspace{4pt}

% teaching: NEW-U15-LOCAL-HELPS
% teaching: U15-H-28-07
% support: U15-H-28-07
% mdparagraph:P427:begin
Use the supplied expression ðā endleftan tīd, “the eleventh hour”, with the feminine accusative singular weak ordinal.
% mdparagraph:P427:end

% output: EX-IX-28-11-O001
\reserve{5\baselineskip}
\par\begin{minipage}{\linewidth}\raggedright
% mdparagraph:P428:begin
\textbf{Cook IX.28.12}
% mdparagraph:P428:end

% item: EX-IX-28-12
% mdparagraph:P429:begin
And those, that there (ðǣr) first came, supposed that they should (sculan) receive more (māre).
% mdparagraph:P429:end

% output: EX-IX-28-12-O001
\end{minipage}\par\addvspace{4pt}

\reserve{7\baselineskip}
% mdparagraph:P430:begin
\subsection*{Cook IX.29}
% mdparagraph:P430:end

\reserve{5\baselineskip}
\par\begin{minipage}{\linewidth}\raggedright
% mdparagraph:P431:begin
\textbf{Cook IX.29.1}
% mdparagraph:P431:end

% item: EX-IX-29-01
% mdparagraph:P432:begin
Thou (om. hast) made (dōn) them equal (gelīc) (om. unto) us, which (87. d) (om. have) borne burdens (51. 5) in this day’s heat.
% mdparagraph:P432:end

% mdparagraph:P433:begin
\textbf{Modern help — available throughout this item.}
% mdparagraph:P433:end

\end{minipage}\par\addvspace{4pt}

% teaching: U15-COOK-NOM-29-01
% support: U15-COOK-NOM-29-01
% mdparagraph:P434:begin
Hęfignessa is accusative plural “burdens”; hǣte is the selected dative singular “heat”.
% mdparagraph:P434:end

% teaching: NEW-U15-LOCAL-HELPS
% teaching: U15-H-29-01
% support: U15-H-29-01
% mdparagraph:P435:begin
On ðisses dæges hǣte means “in this day’s heat”: on with dative gives the circumstance; ðisses dæges is a dependent genitive.
% mdparagraph:P435:end

% output: EX-IX-29-01-O001
\reserve{5\baselineskip}
\par\begin{minipage}{\linewidth}\raggedright
% mdparagraph:P436:begin
\textbf{Cook IX.29.2}
% mdparagraph:P436:end

% item: EX-IX-29-02
% mdparagraph:P437:begin
Take that is thine, and go (om. thy way); I will give unto this (om. one) just (eall) as much as (om. unto) thee.
% mdparagraph:P437:end

% output: EX-IX-29-02-O001
\end{minipage}\par\addvspace{4pt}

\par\begin{minipage}{\linewidth}\raggedright
% mdparagraph:P438:begin
\textbf{Cook IX.29.3}
% mdparagraph:P438:end

% item: EX-IX-29-03
% mdparagraph:P439:begin
May I not do what (ðæt) I will?
% mdparagraph:P439:end

% output: EX-IX-29-03-O001
\end{minipage}\par\addvspace{4pt}

\par\begin{minipage}{\linewidth}\raggedright
% mdparagraph:P440:begin
\textbf{Cook IX.29.4}
% mdparagraph:P440:end

% item: EX-IX-29-04
% mdparagraph:P441:begin
Camest thou not to me to work (wyrcean, 200) for (wið) one penny?
% mdparagraph:P441:end

% output: EX-IX-29-04-O001
\end{minipage}\par\addvspace{4pt}

\par\begin{minipage}{\linewidth}\raggedright
% mdparagraph:P442:begin
\textbf{Cook IX.29.5}
% mdparagraph:P442:end

% item: EX-IX-29-05
% mdparagraph:P443:begin
He was (om. a) man (wer) and prophet, mighty in speech and in deed (weorc) before (beforan, 166) God and all (om. the) people (folc).
% mdparagraph:P443:end

% output: EX-IX-29-05-O001
\end{minipage}\par\addvspace{4pt}

\par\begin{minipage}{\linewidth}\raggedright
% mdparagraph:P444:begin
\textbf{Cook IX.29.6}
% mdparagraph:P444:end

% item: EX-IX-29-06
% mdparagraph:P445:begin
That is the judgment, that light (om. is) come into (on, 172) (om. the) world, and men loved darkness (plur.) rather (swīðor) than the light, because their deeds (weorc) were evil.
% mdparagraph:P445:end

% output: EX-IX-29-06-O001
\end{minipage}\par\addvspace{4pt}

\par\begin{minipage}{\linewidth}\raggedright
% mdparagraph:P446:begin
\textbf{Cook IX.29.7}
% mdparagraph:P446:end

% item: EX-IX-29-07
% mdparagraph:P447:begin
God sent not his son into (om. the) world that he (om. should) condemn the world, but that the world (om. may) be saved (hǣlan, 113) by (ðurh) him.
% mdparagraph:P447:end

% mdparagraph:P448:begin
\textbf{Modern help — available throughout this item.}
% mdparagraph:P448:end

\end{minipage}\par\addvspace{4pt}

% teaching: NEW-U15-LOCAL-HELPS
% teaching: U15-H-29-07
% support: U15-H-29-07
% mdparagraph:P449:begin
Use the supplied spelling ðurh beside the familiar þurh. Here ðurh hine means “through/\allowbreak{}by way of him”. Negation rejects the condemning purpose and contrasts it with saving; it does not deny that the Son was sent.
% mdparagraph:P449:end

% output: EX-IX-29-07-O001
\reserve{5\baselineskip}
\par\begin{minipage}{\linewidth}\raggedright
% mdparagraph:P450:begin
\textbf{Cook IX.29.8}
% mdparagraph:P450:end

% item: EX-IX-29-08
% mdparagraph:P451:begin
Your reward (om. shall) be great, and ye shall be (om. the) children of the Highest; because he is kind (gōd) unto (ofer, 172) (om. the) unthankful and evil.
% mdparagraph:P451:end

% mdparagraph:P452:begin
\textbf{Modern help — available throughout this item.}
% mdparagraph:P452:end

\end{minipage}\par\addvspace{4pt}

% teaching: NEW-U15-LOCAL-HELPS
% teaching: U15-H-29-08
% support: U15-H-29-08
% mdparagraph:P453:begin
Use ðæs Hēahstan, “of the Highest”, as a supplied masculine genitive singular divine title. No further superlative forms of hēah are required.
% mdparagraph:P453:end

% output: EX-IX-29-08-O001
\reserve{5\baselineskip}
\par\begin{minipage}{\linewidth}\raggedright
% mdparagraph:P454:begin
\textbf{Cook IX.29.9}
% mdparagraph:P454:end

% item: EX-IX-29-09
% mdparagraph:P455:begin
Give, and (om. it) shall be given you.
% mdparagraph:P455:end

% mdparagraph:P456:begin
\textbf{Modern help — available throughout this item.}
% mdparagraph:P456:end

\end{minipage}\par\addvspace{4pt}

% teaching: NEW-U15-LOCAL-HELPS
% teaching: U15-H-29-09
% support: U15-H-29-09
% mdparagraph:P457:begin
Use the supplied frame ēow bið giefen, “it will be given to you (plural).” Ēow is the dative recipient; no hit is expressed. For one addressee, use ðē bið giefen with a singular command. Keep the frame available while composing; its independent recall is not required.
% mdparagraph:P457:end

% output: EX-IX-29-09-O001
\reserve{5\baselineskip}
\par\begin{minipage}{\linewidth}\raggedright
% mdparagraph:P458:begin
\textbf{Cook IX.29.10}
% mdparagraph:P458:end

% item: EX-IX-29-10
% mdparagraph:P459:begin
The disciple (leornungcniht, 43) is not above (ofer, 172) his master (lārēow).
% mdparagraph:P459:end

% output: EX-IX-29-10-O001
\end{minipage}\par\addvspace{4pt}

\par\begin{minipage}{\linewidth}\raggedright
% mdparagraph:P460:begin
\textbf{Cook IX.29.11}
% mdparagraph:P460:end

% item: EX-IX-29-11
% mdparagraph:P461:begin
Show us thy mercy, and grant us thy salvation.
% mdparagraph:P461:end

% mdparagraph:P462:begin
\textbf{Modern help — available throughout this item.}
% mdparagraph:P462:end

\end{minipage}\par\addvspace{4pt}

% teaching: U15-COOK-NOM-29-11
% support: U15-COOK-NOM-29-11
% mdparagraph:P463:begin
Hǣlu, “salvation”, is unchanged through the selected singular. Unnan governs genitive of the thing granted, so use genitive ðīnre with it even though hǣlu shows no separate genitive ending. Hǣle is a recognized oblique alternative.
% mdparagraph:P463:end

% output: EX-IX-29-11-O001
\reserve{7\baselineskip}
% mdparagraph:P464:begin
\subsection*{Cook IX.30}
% mdparagraph:P464:end

\reserve{5\baselineskip}
\par\begin{minipage}{\linewidth}\raggedright
% mdparagraph:P465:begin
\textbf{Cook IX.30.1}
% mdparagraph:P465:end

% item: EX-IX-30-01
% mdparagraph:P466:begin
Swear (thou) not by (ðurh) heaven, because it is God’s throne; nor by earth, because it is his footstool (fōtsceamul, 43); nor by Jerusalem (Hierusalem), because it is (om. the) city of (om. the) great (mǣre) king.
% mdparagraph:P466:end

% mdparagraph:P467:begin
\textbf{Modern help — available throughout this item.}
% mdparagraph:P467:end

\end{minipage}\par\addvspace{4pt}

% teaching: U15-COOK-NOM-30-01
% support: U15-COOK-NOM-30-01
% mdparagraph:P468:begin
Setl means “seat; throne”. Fōtsceamul retains -sceamul- in its inflected stem.
% mdparagraph:P468:end

% teaching: NEW-U15-LOCAL-HELPS
% teaching: NEW-U15-VOC-THURH
% teaching: U15-H-30-01
% support: U15-H-30-01
% mdparagraph:P469:begin
Use the supplied spelling ðurh before each accusative oath standard. Hierusalem is the supplied place name; its other cases are not required.
% mdparagraph:P469:end

% output: EX-IX-30-01-O001
\reserve{5\baselineskip}
\par\begin{minipage}{\linewidth}\raggedright
% mdparagraph:P470:begin
\textbf{Cook IX.30.2}
% mdparagraph:P470:end

% item: EX-IX-30-02
% mdparagraph:P471:begin
Swear not by thy head, because thou canst not make (wyrcean) one hair (hǣr, 47) white or black.
% mdparagraph:P471:end

% mdparagraph:P472:begin
\textbf{Modern help — available throughout this item.}
% mdparagraph:P472:end

\end{minipage}\par\addvspace{4pt}

% teaching: U15-COOK-NOM-30-02
% support: U15-COOK-NOM-30-02
% mdparagraph:P473:begin
Hēafod uses hēafd- before endings. Hǣr is neuter and has no nominative/\allowbreak{}accusative plural ending. Fōtsceamul, by contrast, retains its medial vowel.
% mdparagraph:P473:end

% teaching: NEW-U15-LOCAL-HELPS
% teaching: NEW-U15-VOC-THURH
% teaching: U15-H-30-01
% support: U15-H-30-01
% mdparagraph:P474:begin
Use the supplied spelling ðurh before each accusative oath standard. Hierusalem is the supplied place name; its other cases are not required.
% mdparagraph:P474:end

% output: EX-IX-30-02-O001
\reserve{5\baselineskip}
\par\begin{minipage}{\linewidth}\raggedright
% mdparagraph:P475:begin
\textbf{Cook IX.30.3}
% mdparagraph:P475:end

% item: EX-IX-30-03
% mdparagraph:P476:begin
Be your (plur.) speech thus, It is, It is not.
% mdparagraph:P476:end

% output: EX-IX-30-03-O001
\end{minipage}\par\addvspace{4pt}

\par\begin{minipage}{\linewidth}\raggedright
% mdparagraph:P477:begin
\textbf{Cook IX.30.4}
% mdparagraph:P477:end

% item: EX-IX-30-04
% mdparagraph:P478:begin
Ye (om. have) heard that (om. it) was said, Eye for (for, 166) eye, and tooth (46) for tooth.
% mdparagraph:P478:end

% output: EX-IX-30-04-O001
\end{minipage}\par\addvspace{4pt}

\par\begin{minipage}{\linewidth}\raggedright
% mdparagraph:P479:begin
\textbf{Cook IX.30.5}
% mdparagraph:P479:end

% item: EX-IX-30-05
% mdparagraph:P480:begin
He that breaks (tōweorpan, 104) one of (of, 166) these least commandments (bebod, 47), and teacheth (lǣran) men so (verb last), he (sē) shall be named least in (om. the) kingdom (rīce) of heaven.
% mdparagraph:P480:end

% mdparagraph:P481:begin
\textbf{Modern help — available throughout this item.}
% mdparagraph:P481:end

\end{minipage}\par\addvspace{4pt}

% teaching: NEW-U15-LOCAL-HELPS
% teaching: NEW-U15-RECORD-OF-U15-SOURCE-PARTITIVE
% teaching: U15-H-30-05
% support: U15-H-30-05
% mdparagraph:P482:begin
The emphatic resumed sē is a supplied long-vowel nominative masculine singular. In ān of ðissum lǣstum bebodum, “one of these least commandments”, of takes dative to identify the group. The same long sē may be used in 30.6.
% mdparagraph:P482:end

% output: EX-IX-30-05-O001
\reserve{5\baselineskip}
\par\begin{minipage}{\linewidth}\raggedright
% mdparagraph:P483:begin
\textbf{Cook IX.30.6}
% mdparagraph:P483:end

% item: EX-IX-30-06
% mdparagraph:P484:begin
He that doeth and teacheth it, he shall be called great (micel).
% mdparagraph:P484:end

% mdparagraph:P485:begin
\textbf{Modern help — available throughout this item.}
% mdparagraph:P485:end

\end{minipage}\par\addvspace{4pt}

% teaching: NEW-U15-LOCAL-HELPS
% teaching: NEW-U15-RECORD-OF-U15-SOURCE-PARTITIVE
% teaching: U15-H-30-05
% support: U15-H-30-05
% mdparagraph:P486:begin
The emphatic resumed sē is a supplied long-vowel nominative masculine singular. In ān of ðissum lǣstum bebodum, “one of these least commandments”, of takes dative to identify the group. The same long sē may be used in 30.6.
% mdparagraph:P486:end

% output: EX-IX-30-06-O001
\reserve{5\baselineskip}
\par\begin{minipage}{\linewidth}\raggedright
% mdparagraph:P487:begin
\textbf{Cook IX.30.7}
% mdparagraph:P487:end

% item: EX-IX-30-07
% mdparagraph:P488:begin
Ye (om. have) heard that (om. it) was said in old times (tīd), Kill (thou) not.
% mdparagraph:P488:end

% mdparagraph:P489:begin
\textbf{Modern help — available throughout this item.}
% mdparagraph:P489:end

\end{minipage}\par\addvspace{4pt}

% teaching: NEW-U15-LOCAL-HELPS
% teaching: NEW-U15-RECORD-ON-U15-SOURCE-DURING
% teaching: U15-H-30-07
% support: U15-H-30-07
% mdparagraph:P490:begin
Use the supplied on ealdum tīdum, “in old times”, with temporal on and dative plural.
% mdparagraph:P490:end

% output: EX-IX-30-07-O001
\reserve{5\baselineskip}
\par\begin{minipage}{\linewidth}\raggedright
% mdparagraph:P491:begin
\textbf{Cook IX.30.8}
% mdparagraph:P491:end

% item: EX-IX-30-08
% mdparagraph:P492:begin
Is it your (sing.) earth?
% mdparagraph:P492:end

% mdparagraph:P493:begin
\textbf{Modern help — available throughout this item.}
% mdparagraph:P493:end

\end{minipage}\par\addvspace{4pt}

% teaching: U15-COOK-NOM-30-08
% support: U15-COOK-NOM-30-08
% mdparagraph:P494:begin
Eorðe, “earth”, is feminine. Its possessive must agree with the noun that it modifies.
% mdparagraph:P494:end

% output: EX-IX-30-08-O001
\reserve{5\baselineskip}
\par\begin{minipage}{\linewidth}\raggedright
% mdparagraph:P495:begin
\textbf{Cook IX.30.9}
% mdparagraph:P495:end

% item: EX-IX-30-09
% mdparagraph:P496:begin
He has from his labor what he can.
% mdparagraph:P496:end

% mdparagraph:P497:begin
\textbf{Modern help — available throughout this item.}
% mdparagraph:P497:end

\end{minipage}\par\addvspace{4pt}

% teaching: U15-COOK-NOM-30-09
% support: U15-COOK-NOM-30-09
% mdparagraph:P498:begin
Gewinne, “labour”, is a supplied noun form in the expression æt his gewinne.
% mdparagraph:P498:end

% teaching: NEW-U15-LOCAL-HELPS
% teaching: U15-H-30-09
% support: U15-H-30-09
% mdparagraph:P499:begin
Use æt his gewinne, “from his labour”. In the supported ðæt hē mæg, “what he can”, an action such as “obtain” is understood. Retain this supplied expression without deriving further forms of gewinn.
% mdparagraph:P499:end

% output: EX-IX-30-09-O001
\reserve{5\baselineskip}
\par\begin{minipage}{\linewidth}\raggedright
% mdparagraph:P500:begin
\textbf{Cook IX.30.10}
% mdparagraph:P500:end

% item: EX-IX-30-10
% mdparagraph:P501:begin
The sun parched his field.
% mdparagraph:P501:end

% output: EX-IX-30-10-O001
\end{minipage}\par\addvspace{4pt}

\par\begin{minipage}{\linewidth}\raggedright
% mdparagraph:P502:begin
\textbf{Cook IX.30.11}
% mdparagraph:P502:end

% item: EX-IX-30-11
% mdparagraph:P503:begin
The roads (weg) should blossom (blōstmian, 118), the roads should bloom, so fair (fæger) (om. a) bride (brȳd, 51. b) shall leave her home (hām).
% mdparagraph:P503:end

% mdparagraph:P504:begin
\textbf{Modern help — available throughout this item.}
% mdparagraph:P504:end

\end{minipage}\par\addvspace{4pt}

% teaching: U15-COOK-NOM-30-11
% support: U15-COOK-NOM-30-11
% mdparagraph:P505:begin
Brȳd is feminine; the selected accusative singular is brȳde. Hām is the noun “home”, here an accusative object, rather than the directional adverb. Use the given infinitive blōstmian for “blossom”.
% mdparagraph:P505:end

% output: EX-IX-30-11-O001
\reserve{5\baselineskip}
\par\begin{minipage}{\linewidth}\raggedright
% mdparagraph:P506:begin
\textbf{Cook IX.30.12}
% mdparagraph:P506:end

% item: EX-IX-30-12
% mdparagraph:P507:begin
If they can, let them live without me.
% mdparagraph:P507:end

% mdparagraph:P508:begin
\textbf{Modern help — available throughout this item.}
% mdparagraph:P508:end

\end{minipage}\par\addvspace{4pt}

% teaching: NEW-U15-LOCAL-HELPS
% teaching: U15-H-30-12
% support: U15-H-30-12
% mdparagraph:P509:begin
Use the supplied permission construction lǣt hīe libban, “let them live”: accusative hīe is the understood subject of libban. Gif hīe magon expresses ability. A third-person optative is also a valid way to express the request.
% mdparagraph:P509:end

% output: EX-IX-30-12-O001
\reserve{5\baselineskip}
\par\begin{minipage}{\linewidth}\raggedright
% mdparagraph:P510:begin
\textbf{Cook IX.30.13}
% mdparagraph:P510:end

% item: EX-IX-30-13
% mdparagraph:P511:begin
He never forsakes those that trust in (on) him (172).
% mdparagraph:P511:end

% output: EX-IX-30-13-O001
\end{minipage}\par\addvspace{4pt}

\reserve{7\baselineskip}
% mdparagraph:P512:begin
\subsection*{Cook X, XI.31}
% mdparagraph:P512:end

\reserve{5\baselineskip}
\par\begin{minipage}{\linewidth}\raggedright
% mdparagraph:P513:begin
\textbf{Cook X, XI.31.1}
% mdparagraph:P513:end

% item: EX-X-XI-31-01
% mdparagraph:P514:begin
I know truly that my Redeemer liveth.
% mdparagraph:P514:end

% mdparagraph:P515:begin
\textbf{Modern help — available throughout this item.}
% mdparagraph:P515:end

\end{minipage}\par\addvspace{4pt}

% teaching: U15-COOK-NOM-31-01
% support: U15-COOK-NOM-31-01
% mdparagraph:P516:begin
Use ālīesend as the singular title “Redeemer”. Its displayed plurals remain reference forms.
% mdparagraph:P516:end

% output: EX-X-XI-31-01-O001
\reserve{5\baselineskip}
\par\begin{minipage}{\linewidth}\raggedright
% mdparagraph:P517:begin
\textbf{Cook X, XI.31.2}
% mdparagraph:P517:end

% item: EX-X-XI-31-02
% mdparagraph:P518:begin
In my flesh I shall see God, I (om. my)self, and not (nā) (om. an)other.
% mdparagraph:P518:end

% output: EX-X-XI-31-02-O001
\end{minipage}\par\addvspace{4pt}

\par\begin{minipage}{\linewidth}\raggedright
% mdparagraph:P519:begin
\textbf{Cook X, XI.31.3}
% mdparagraph:P519:end

% item: EX-X-XI-31-03
% mdparagraph:P520:begin
He was taught (187) to do God’s will.
% mdparagraph:P520:end

% mdparagraph:P521:begin
\textbf{Modern help — available throughout this item.}
% mdparagraph:P521:end

\end{minipage}\par\addvspace{4pt}

% teaching: NEW-U15-LOCAL-HELPS
% teaching: U15-H-31-03
% support: U15-H-31-03
% mdparagraph:P522:begin
Use hē wæs lǣred … tō dōnne. The person taught is the passive subject and the understood doer of the infinitive. Godes willan is the infinitive’s accusative object, with dependent genitive Godes. This whole passive-and-infinitive frame is supplied.
% mdparagraph:P522:end

% output: EX-X-XI-31-03-O001
\reserve{5\baselineskip}
\par\begin{minipage}{\linewidth}\raggedright
% mdparagraph:P523:begin
\textbf{Cook X, XI.31.4}
% mdparagraph:P523:end

% item: EX-X-XI-31-04
% mdparagraph:P524:begin
Lead (thou) us not into temptation, but deliver us from evil.
% mdparagraph:P524:end

% mdparagraph:P525:begin
\textbf{Modern help — available throughout this item.}
% mdparagraph:P525:end

\end{minipage}\par\addvspace{4pt}

% teaching: NEW-U15-LOCAL-HELPS
% teaching: U15-H-31-04
% support: U15-H-31-04
% mdparagraph:P526:begin
Use fram yfele, “from evil”. Yfele is a supplied substantive noun in the dative singular, not the ordinary adjective dative yfelum.
% mdparagraph:P526:end

% output: EX-X-XI-31-04-O001
\reserve{5\baselineskip}
\par\begin{minipage}{\linewidth}\raggedright
% mdparagraph:P527:begin
\textbf{Cook X, XI.31.5}
% mdparagraph:P527:end

% item: EX-X-XI-31-05
% mdparagraph:P528:begin
Be thy name hallowed.
% mdparagraph:P528:end

% output: EX-X-XI-31-05-O001
\end{minipage}\par\addvspace{4pt}

\par\begin{minipage}{\linewidth}\raggedright
% mdparagraph:P529:begin
\textbf{Cook X, XI.31.6}
% mdparagraph:P529:end

% item: EX-X-XI-31-06
% mdparagraph:P530:begin
Thy will befall (geweorðan, verb first) in earth as (swā-swā) in heaven(s).
% mdparagraph:P530:end

% mdparagraph:P531:begin
\textbf{Modern help — available throughout this item.}
% mdparagraph:P531:end

\end{minipage}\par\addvspace{4pt}

% teaching: NEW-U09-SOURCE-HELPS
% teaching: NEW-U09-VOC-swa-swa
% teaching: NEW-U10-SOURCE-HELPS
% teaching: NEW-U12-REUSED-U12-SRC-HELP-15
% teaching: NEW-U12-SOURCE-HELPS
% teaching: NEW-U13-PAR-SWA-SWA
% teaching: NEW-U13-VOC-SWA-SWA
% teaching: NEW-U14-LOCAL-HELPS
% teaching: NEW-U15-LOCAL-HELPS
% teaching: U15-H-31-06
% support: U15-H-31-06
% mdparagraph:P532:begin
Begin with supplied Geweorðe, the third-singular present subjunctive of the prefixed verb. Swā-swā means “just as”; retain this supplied hyphenated expression.
% mdparagraph:P532:end

% output: EX-X-XI-31-06-O001
\reserve{5\baselineskip}
\par\begin{minipage}{\linewidth}\raggedright
% mdparagraph:P533:begin
\textbf{Cook X, XI.31.7}
% mdparagraph:P533:end

% item: EX-X-XI-31-07
% mdparagraph:P534:begin
No king of the English was so great as Alfred.
% mdparagraph:P534:end

% mdparagraph:P535:begin
\textbf{Modern help — available throughout this item.}
% mdparagraph:P535:end

\end{minipage}\par\addvspace{4pt}

% teaching: NEW-U15-LOCAL-HELPS
% teaching: NEW-U15-RECORD-AELFRED-U15-SOURCE-NOM-SG
% teaching: U15-H-31-NAMES
% support: U15-H-31-NAMES
% mdparagraph:P536:begin
The names Ælfred, Osburg, Æthelwulf and Asser are supplied; Ælfredes is the genitive for 31.14. Use supplied dative Lǣdene for “Latin”, following Cook’s glossary Lǣden. The Preface reading below retains its printed Lædene without a length mark.
% mdparagraph:P536:end

% output: EX-X-XI-31-07-O001
\reserve{5\baselineskip}
\par\begin{minipage}{\linewidth}\raggedright
% mdparagraph:P537:begin
\textbf{Cook X, XI.31.8}
% mdparagraph:P537:end

% item: EX-X-XI-31-08
% mdparagraph:P538:begin
He fought with the Danes (44. 2); he judged his people; he made translations from Latin; he was poet (lēoðwyrhta).
% mdparagraph:P538:end

% mdparagraph:P539:begin
\textbf{Modern help — available throughout this item.}
% mdparagraph:P539:end

\end{minipage}\par\addvspace{4pt}

% teaching: U15-COOK-NOM-31-08
% support: U15-COOK-NOM-31-08
% mdparagraph:P540:begin
Dęne and Engle name peoples in the plural. Lēoðwyrhta is a weak masculine noun, “poet”. Āwęndednes means “translation” and doubles s before endings. The language-name dative is supplied separately.
% mdparagraph:P540:end

% teaching: NEW-U15-LOCAL-HELPS
% teaching: NEW-U15-RECORD-AELFRED-U15-SOURCE-NOM-SG
% teaching: U15-H-31-NAMES
% support: U15-H-31-NAMES
% mdparagraph:P541:begin
The names Ælfred, Osburg, Æthelwulf and Asser are supplied; Ælfredes is the genitive for 31.14. Use supplied dative Lǣdene for “Latin”, following Cook’s glossary Lǣden. The Preface reading below retains its printed Lædene without a length mark.
% mdparagraph:P541:end

% output: EX-X-XI-31-08-O001
\reserve{5\baselineskip}
\par\begin{minipage}{\linewidth}\raggedright
% mdparagraph:P542:begin
\textbf{Cook X, XI.31.9}
% mdparagraph:P542:end

% item: EX-X-XI-31-09
% mdparagraph:P543:begin
His mother hight Osburg, and his father Æthelwulf.
% mdparagraph:P543:end

% mdparagraph:P544:begin
\textbf{Modern help — available throughout this item.}
% mdparagraph:P544:end

\end{minipage}\par\addvspace{4pt}

% teaching: NEW-U15-LOCAL-HELPS
% teaching: NEW-U15-RECORD-AELFRED-U15-SOURCE-NOM-SG
% teaching: U15-H-31-NAMES
% support: U15-H-31-NAMES
% mdparagraph:P545:begin
The names Ælfred, Osburg, Æthelwulf and Asser are supplied; Ælfredes is the genitive for 31.14. Use supplied dative Lǣdene for “Latin”, following Cook’s glossary Lǣden. The Preface reading below retains its printed Lædene without a length mark.
% mdparagraph:P545:end

% teaching: NEW-U15-LOCAL-HELPS
% teaching: U15-H-31-09
% support: U15-H-31-09
% mdparagraph:P546:begin
Supplied hātte means “is/\allowbreak{}was called”; context supplies past time here. It is a special passive naming form, not the active preterite hēt.
% mdparagraph:P546:end

% output: EX-X-XI-31-09-O001
\reserve{5\baselineskip}
\par\begin{minipage}{\linewidth}\raggedright
% mdparagraph:P547:begin
\textbf{Cook X, XI.31.10}
% mdparagraph:P547:end

% item: EX-X-XI-31-10
% mdparagraph:P548:begin
He lived fifty-two years.
% mdparagraph:P548:end

% mdparagraph:P549:begin
\textbf{Modern help — available throughout this item.}
% mdparagraph:P549:end

\end{minipage}\par\addvspace{4pt}

% teaching: NEW-U15-LOCAL-HELPS
% teaching: U15-H-31-10
% support: U15-H-31-10
% mdparagraph:P550:begin
Use the supplied number phrase twā and fīftig gēara, “fifty-two years”. Fīftig means “fifty”; gēara is genitive plural governed by this numeral expression. No further higher numerals are required.
% mdparagraph:P550:end

% output: EX-X-XI-31-10-O001
\reserve{5\baselineskip}
\par\begin{minipage}{\linewidth}\raggedright
% mdparagraph:P551:begin
\textbf{Cook X, XI.31.11}
% mdparagraph:P551:end

% item: EX-X-XI-31-11
% mdparagraph:P552:begin
He learned reading after he was twelve (155) years (154. c) old.
% mdparagraph:P552:end

% mdparagraph:P553:begin
\textbf{Modern help — available throughout this item.}
% mdparagraph:P553:end

\end{minipage}\par\addvspace{4pt}

% teaching: U15-COOK-NOM-31-11
% support: U15-COOK-NOM-31-11
% mdparagraph:P554:begin
Bēcrǣding means “reading”. Use its singular; the displayed plurals are for reference.
% mdparagraph:P554:end

% teaching: NEW-U15-LOCAL-HELPS
% teaching: U15-H-31-11
% support: U15-H-31-11
% mdparagraph:P555:begin
Use æfter ðǣm ðe … wæs for “after he was …”. Here ðe belongs to the supplied after-conjunction and has no separate subject or object role. In twęlf wintra eald, “twelve years old”, wintra is genitive plural with the numeral and adjective of age. Keep both expressions available as help.
% mdparagraph:P555:end

% output: EX-X-XI-31-11-O001
\reserve{5\baselineskip}
\par\begin{minipage}{\linewidth}\raggedright
% mdparagraph:P556:begin
\textbf{Cook X, XI.31.12}
% mdparagraph:P556:end

% item: EX-X-XI-31-12
% mdparagraph:P557:begin
His mother had a book of English poems, which she promised (om. to) whichever of her two sons (om. should) first learn it.
% mdparagraph:P557:end

% mdparagraph:P558:begin
\textbf{Modern help — available throughout this item.}
% mdparagraph:P558:end

\end{minipage}\par\addvspace{4pt}

% teaching: U15-COOK-NOM-31-12
% support: U15-COOK-NOM-31-12
% mdparagraph:P559:begin
Bōc is accusative singular “book”; bēc belongs to other cases in its table. Lēoða is genitive plural “poems”.
% mdparagraph:P559:end

% teaching: NEW-U15-LOCAL-HELPS
% teaching: U15-H-31-12
% support: U15-H-31-12
% mdparagraph:P560:begin
Use Engliscra lēoða, “of English poems”. The supplied recipient expression is swā hwæðrum hiere twēgra suna swā hīe ǣrest leornode. Hwæðrum is dative because the whole expression is the recipient of the promise; the same person is understood as the learner in the subordinate clause. Leornode is prospective preterite subjunctive, identical in spelling to the indicative; the learning is still to occur. This case attraction is supplied. In the independent free-choice exercises, keep the internal subject nominative and resume it with an explicit dative recipient.
% mdparagraph:P560:end

% output: EX-X-XI-31-12-O001
\reserve{5\baselineskip}
\par\begin{minipage}{\linewidth}\raggedright
% mdparagraph:P561:begin
\textbf{Cook X, XI.31.13}
% mdparagraph:P561:end

% item: EX-X-XI-31-13
% mdparagraph:P562:begin
Alfred asked her, Wilt thou in fact give that book to one of us two (uncer), (om. to) him who first understands (ongietan) it?
% mdparagraph:P562:end

% mdparagraph:P563:begin
\textbf{Modern help — available throughout this item.}
% mdparagraph:P563:end

\end{minipage}\par\addvspace{4pt}

% teaching: NEW-U15-LOCAL-HELPS
% teaching: NEW-U15-RECORD-AELFRED-U15-SOURCE-NOM-SG
% teaching: U15-H-31-NAMES
% support: U15-H-31-NAMES
% mdparagraph:P564:begin
The names Ælfred, Osburg, Æthelwulf and Asser are supplied; Ælfredes is the genitive for 31.14. Use supplied dative Lǣdene for “Latin”, following Cook’s glossary Lǣden. The Preface reading below retains its printed Lædene without a length mark.
% mdparagraph:P564:end

% output: EX-X-XI-31-13-O001
\reserve{5\baselineskip}
\par\begin{minipage}{\linewidth}\raggedright
% mdparagraph:P565:begin
\textbf{Cook X, XI.31.14}
% mdparagraph:P565:end

% item: EX-X-XI-31-14
% mdparagraph:P566:begin
Asser, Alfred’s bishop, wrote a book of (153) the king’s life.
% mdparagraph:P566:end

% mdparagraph:P567:begin
\textbf{Modern help — available throughout this item.}
% mdparagraph:P567:end

\end{minipage}\par\addvspace{4pt}

% teaching: NEW-U15-LOCAL-HELPS
% teaching: NEW-U15-RECORD-AELFRED-U15-SOURCE-NOM-SG
% teaching: U15-H-31-NAMES
% support: U15-H-31-NAMES
% mdparagraph:P568:begin
The names Ælfred, Osburg, Æthelwulf and Asser are supplied; Ælfredes is the genitive for 31.14. Use supplied dative Lǣdene for “Latin”, following Cook’s glossary Lǣden. The Preface reading below retains its printed Lædene without a length mark.
% mdparagraph:P568:end

% output: EX-X-XI-31-14-O001
\reserve{7\baselineskip}
% mdparagraph:P569:begin
\section*{8. Readings}
% mdparagraph:P569:end

% mdparagraph:P570:begin
The Old English passages and the numbered original notes are exact source material, with declared typographic normalization. Modern vocabulary and construction helps precede each passage and remain available for all five tasks on it. Original note numbers are prefixed D, P or R here only to distinguish the three sets. Cook’s grammatical references and historical questions are retained as apparatus; answer only the tasks explicitly assigned below.
% mdparagraph:P570:end

% teaching: NEW-U15-GRAM-U15-SOURCE-ARGUMENT
% mdparagraph:P571:begin
Attend to each passage’s movement: a rebuke becomes a promise, an account of translation becomes a request to readers, and an appeal becomes a series of commands and a concluding wish. Distinguish rhetorical questions, assertions, purposes and petitions.
% mdparagraph:P571:end

\reserve{7\baselineskip}
% mdparagraph:P572:begin
\subsection*{D. Duties of the rich toward the poor: the opening argument}
% mdparagraph:P572:end

% mdparagraph:P573:begin
Cook, First Book, printed p. 159 (PDF p. 175). This is the complete first argument, including its embedded gospel saying. The Savior’s outer speech continues on pp. 160–161; its opening quotation mark is intentionally left unmatched in this selection.
% mdparagraph:P573:end

% mdparagraph:P574:begin
\textbf{Modern vocabulary — supplied for this reading.} Forms are glossed in their local senses; the parses identify the forms encountered, not complete new paradigms.
% mdparagraph:P574:end

% teaching: NEW-U15-RECORD-AEGHWILC-U15-SOURCE-DAT-M-SG
% teaching: NEW-U15-RECORD-AEGHWILC-U15-SOURCE-SENSE
% teaching: U15-READ-LEX-AEGHWILC
% support: U15-READ-LEX-AEGHWILC
% mdparagraph:P575:begin
\textbf{ǣghwilcum} — each; every (strong dative masculine singular).
% mdparagraph:P575:end

% teaching: NEW-U15-RECORD-AFIERRAN-U15-SOURCE-PRES-1-SG
% teaching: NEW-U15-RECORD-AFIERRAN-U15-SOURCE-SENSE
% teaching: U15-READ-LEX-AFIERRAN
% support: U15-READ-LEX-AFIERRAN
% mdparagraph:P576:begin
\textbf{āfierre} — remove; take away (present indicative first singular, future reference).
% mdparagraph:P576:end

% teaching: NEW-U15-RECORD-CWETHAN-U15-SOURCE-GE-PAST-PART
% teaching: NEW-U15-RECORD-CWETHAN-U15-SOURCE-SENSE
% teaching: U15-READ-LEX-CWETHAN
% support: U15-READ-LEX-CWETHAN
% mdparagraph:P577:begin
\textbf{gecweden} — said (uninflected ge-prefixed past participle in passive).
% mdparagraph:P577:end

% teaching: NEW-U15-RECORD-DAED-U15-SOURCE-ACC-SG
% teaching: NEW-U15-RECORD-DAED-U15-SOURCE-SENSE
% teaching: U15-READ-LEX-DAED
% support: U15-READ-LEX-DAED
% mdparagraph:P578:begin
\textbf{dǣd} — deed (endingless accusative feminine singular).
% mdparagraph:P578:end

% teaching: NEW-U15-RECORD-DREAM-U15-SOURCE-ACC-SG
% teaching: NEW-U15-RECORD-DREAM-U15-SOURCE-SENSE
% teaching: U15-READ-LEX-DREAM
% support: U15-READ-LEX-DREAM
% mdparagraph:P579:begin
\textbf{drēam} — joy; bliss (accusative masculine singular).
% mdparagraph:P579:end

% teaching: NEW-U15-RECORD-FAESTHAFOL-U15-SOURCE-NOM-M-SG-STRONG
% teaching: NEW-U15-RECORD-FAESTHAFOL-U15-SOURCE-SENSE
% teaching: U15-READ-LEX-FAESTHAFOL
% support: U15-READ-LEX-FAESTHAFOL
% mdparagraph:P580:begin
\textbf{fæsthafol} — tenacious; withholding, with genitive of what is held (predicate nominative masculine singular).
% mdparagraph:P580:end

% teaching: NEW-U11-LOCAL-HELPS
% teaching: NEW-U15-RECORD-FOR-U15-SOURCE-WHY
% teaching: U15-READ-LEX-FOR
% support: U15-READ-LEX-FOR
% mdparagraph:P581:begin
\textbf{for} — for what reason; why, in for hwȳ/\allowbreak{}hwon. \textbf{for} — for; for the sake of, selected accusative/\allowbreak{}dative expressions supplied locally.
% mdparagraph:P581:end

% teaching: NEW-U15-RECORD-GEDON-U15-SOURCE-PRES-3-SG
% teaching: NEW-U15-RECORD-GEDON-U15-SOURCE-SENSE
% teaching: U15-READ-LEX-GEDON
% support: U15-READ-LEX-GEDON
% mdparagraph:P582:begin
\textbf{gedēð} — do; perform (present indicative third singular).
% mdparagraph:P582:end

% teaching: NEW-U15-RECORD-GESELLAN-U15-SOURCE-PRES-1-SG
% teaching: NEW-U15-RECORD-GESELLAN-U15-SOURCE-SENSE
% teaching: U15-READ-LEX-GESELLAN
% support: U15-READ-LEX-GESELLAN
% mdparagraph:P583:begin
\textbf{gesęlle} — give (present indicative first singular, future reference).
% mdparagraph:P583:end

% teaching: NEW-U15-RECORD-GETHENCEAN-U15-SOURCE-INF
% teaching: NEW-U15-RECORD-GETHENCEAN-U15-SOURCE-SENSE
% teaching: U15-READ-LEX-GETHENCEAN
% support: U15-READ-LEX-GETHENCEAN
% mdparagraph:P584:begin
\textbf{geðęncean} — remember; bear in mind (plain infinitive after noldest).
% mdparagraph:P584:end

% teaching: NEW-U15-RECORD-GIEFU-U15-SOURCE-GEN-SG
% teaching: NEW-U15-RECORD-GIEFU-U15-SOURCE-SENSE
% teaching: U15-READ-LEX-GIEFU
% support: U15-READ-LEX-GIEFU
% mdparagraph:P585:begin
\textbf{giefe} — gift (genitive feminine singular dependent on rēcelēasodest).
% mdparagraph:P585:end

% teaching: NEW-U15-RECORD-HUNDTEONTIG-U15-SOURCE-DAT-PL
% teaching: NEW-U15-RECORD-HUNDTEONTIG-U15-SOURCE-SENSE
% teaching: U15-READ-LEX-HUNDTEONTIG
% support: U15-READ-LEX-HUNDTEONTIG
% mdparagraph:P586:begin
\textbf{hundtēontigum} — a hundredfold; a hundred in return (dative plural after mid in fixed repayment expression).
% mdparagraph:P586:end

% teaching: NEW-U15-RECORD-HWA-U15-SOURCE-WHY
% teaching: U15-READ-LEX-HWA
% support: U15-READ-LEX-HWA
% mdparagraph:P587:begin
\textbf{hwon} — what (reason), instrumental hwon in fixed why expressions (instrumental singular neuter).
% mdparagraph:P587:end

% teaching: NEW-U15-RECORD-OF-U15-SOURCE-PARTITIVE
% teaching: U15-READ-LEX-OF
% support: U15-READ-LEX-OF
% mdparagraph:P588:begin
\textbf{of} — of; out of a set, selecting one of the least.
% mdparagraph:P588:end

% teaching: NEW-U15-RECORD-ON-U15-SOURCE-WORLDLY-LIFE
% teaching: U15-READ-LEX-ON
% support: U15-READ-LEX-ON
% mdparagraph:P589:begin
\textbf{on} — in worldly life, temporal/\allowbreak{}situational expression.
% mdparagraph:P589:end

% teaching: NEW-U15-RECORD-RECELEASIAN-U15-SOURCE-PRET-2-SG
% teaching: NEW-U15-RECORD-RECELEASIAN-U15-SOURCE-SENSE
% teaching: U15-READ-LEX-RECELEASIAN
% support: U15-READ-LEX-RECELEASIAN
% mdparagraph:P590:begin
\textbf{rēcelēasodest} — neglect; despise (preterite indicative second singular; genitive object).
% mdparagraph:P590:end

% teaching: NEW-U15-RECORD-SECGAN-U15-SOURCE-GE-PAST-PART
% teaching: NEW-U15-RECORD-SECGAN-U15-SOURCE-SENSE
% teaching: U15-READ-LEX-SECGAN
% support: U15-READ-LEX-SECGAN
% mdparagraph:P591:begin
\textbf{gesǣd} — said (uninflected ge-prefixed past participle, pleonastic with gecweden).
% mdparagraph:P591:end

% teaching: NEW-U15-RECORD-SELEN-U15-SOURCE-DAT-SG
% teaching: NEW-U15-RECORD-SELEN-U15-SOURCE-SENSE
% teaching: U15-READ-LEX-SELEN
% support: U15-READ-LEX-SELEN
% mdparagraph:P592:begin
\textbf{sęlene} — bounty; bestowal (dative feminine singular after fram).
% mdparagraph:P592:end

% teaching: NEW-U15-RECORD-SIMLE-U15-SOURCE-INVARIANT
% teaching: NEW-U15-RECORD-SIMLE-U15-SOURCE-SENSE
% teaching: U15-READ-LEX-SIMLE
% support: U15-READ-LEX-SIMLE
% mdparagraph:P593:begin
\textbf{simle} — always.
% mdparagraph:P593:end

% teaching: NEW-U15-RECORD-SWA-HWAET-SWA-U15-SOURCE-ACC-N-SG
% teaching: NEW-U15-RECORD-SWA-HWAET-SWA-U15-SOURCE-SENSE
% teaching: U15-READ-LEX-SWA-HWAET-SWA
% support: U15-READ-LEX-SWA-HWAET-SWA
% mdparagraph:P594:begin
\textbf{Swā-hwæt-swā} — whatsoever; whatever (free-choice relative pronoun, accusative neuter object).
% mdparagraph:P594:end

% teaching: NEW-U14-LOCAL-HELPS
% teaching: NEW-U15-RECORD-TO-U15-SOURCE-WHY
% teaching: U15-READ-LEX-TO
% support: U15-READ-LEX-TO
% mdparagraph:P595:begin
\textbf{tō} — to, the person addressed in speech. \textbf{tō} — for what purpose; why, in tō hwon.
% mdparagraph:P595:end

% teaching: NEW-U15-RECORD-WLANC-U15-SOURCE-DAT-M-SG-WEAK
% teaching: NEW-U15-RECORD-WLANC-U15-SOURCE-SENSE
% teaching: U15-READ-LEX-WLANC
% support: U15-READ-LEX-WLANC
% mdparagraph:P596:begin
\textbf{wlancan} — proud; wealthy, here the rich man addressed (substantival weak dative masculine singular).
% mdparagraph:P596:end

% teaching: NEW-U15-RECORD-WORULDLIF-U15-SOURCE-DAT-SG
% teaching: NEW-U15-RECORD-WORULDLIF-U15-SOURCE-SENSE
% teaching: U15-READ-LEX-WORULDLIF
% support: U15-READ-LEX-WORULDLIF
% mdparagraph:P597:begin
\textbf{woruldlīfe} — worldly life (dative neuter singular after on).
% mdparagraph:P597:end

\reserve{5\baselineskip}
% mdparagraph:P598:begin
\textbf{Modern construction help.}
% mdparagraph:P598:end

% teaching: NEW-U15-RECORD-U15-SOURCE-D-ORIENTATION
% teaching: U15-READ-D-ORIENTATION
% support: U15-READ-D-ORIENTATION
% mdparagraph:P599:begin
The Savior addresses a rich man. This selection contains the complete opening argument and its embedded gospel quotation. The Savior’s speech continues on the next page: the unmatched outer quotation mark is present in this bounded source selection.
% mdparagraph:P599:end

% teaching: NEW-U15-RECORD-U15-SOURCE-D-GENITIVE-CONTROLS
% teaching: U15-READ-D-GENITIVE-CONTROLS
% support: U15-READ-D-GENITIVE-CONTROLS
% mdparagraph:P600:begin
For hwȳ and tō hwon ask why. Fæsthafol governs the genitive mīnra gōda, “holding tightly to my goods”; rēcelēasodest takes genitive ðǣre giefe, “neglected the gift.” Relative ðe introduces the following clauses; long ðē is the personal pronoun “you/\allowbreak{}to you.”
% mdparagraph:P600:end

% teaching: NEW-U15-RECORD-U15-SOURCE-D-FUTURE-RESPONSE
% teaching: U15-READ-D-FUTURE-RESPONSE
% support: U15-READ-D-FUTURE-RESPONSE
% mdparagraph:P601:begin
Āfierre, bist, forgielde and gesęlle have present morphology with future reference in this promised response. Noldest is “you would not/\allowbreak{}refused”; geðęncean is its infinitive. Wille expresses the present speaker’s intention even within a remembered proposition: no mechanical backshift is required.
% mdparagraph:P601:end

% teaching: NEW-U15-RECORD-U15-SOURCE-D-GENERIC-AGENT
% teaching: U15-READ-D-GENERIC-AGENT
% support: U15-READ-D-GENERIC-AGENT
% mdparagraph:P602:begin
Ǣghwilcum męnn is the dative recipient; āne gōde dǣd is the accusative object. The subject of gedēð is generic mann, “one/\allowbreak{}a person.” Mid hundtēontigum means “with a hundredfold.” The singular him refers back to the person.
% mdparagraph:P602:end

% teaching: NEW-U15-RECORD-U15-SOURCE-D-EMBEDDED-QUOTATION
% teaching: U15-READ-D-EMBEDDED-QUOTATION
% support: U15-READ-D-EMBEDDED-QUOTATION
% mdparagraph:P603:begin
Gecweden and gesǣd are two past participles with is; Cook calls the pairing pleonastic. Swā-hwæt-swā means “whatever,” the object of sęllað. Ānum of mīnum ðǣm lǣstum means “to one of my least ones”; lǣstum is substantival and dative plural. The repeated gē/\allowbreak{}sęllað address the quoted audience.
% mdparagraph:P603:end

% teaching: NEW-U15-RECORD-U15-SOURCE-D-EXCHANGE
% teaching: U15-READ-D-EXCHANGE
% support: U15-READ-D-EXCHANGE
% mdparagraph:P604:begin
Wið ðǣm means “in return for that,” with dative singular. Drēam here means joy or bliss, not a dream. Ēcne agrees with accusative masculine drēam.
% mdparagraph:P604:end

\reserve{5\baselineskip}
% mdparagraph:P605:begin
\textbf{Cook’s text}
% mdparagraph:P605:end

% mdparagraph:P606:begin
IX. DUTIES OF THE RICH TOWARD THE POOR.
% mdparagraph:P606:end

% mdparagraph:P607:begin
(From the same Homily as the last.)
% mdparagraph:P607:end

% teaching: NEW-U15-READ-FB-IX-DUTIES-FIRST-ARGUMENT
\par\begin{minipage}{\linewidth}\raggedright
% passage: FB-IX-DUTIES-FIRST-ARGUMENT
% mdparagraph:P608:begin
\readingtext{Se Hǣlend cwæð tō ðǣm wlancan\textsuperscript{D1}: “For hwȳ wǣre ðū swā fæsthafol mīnra gōda, ðe ic ðē sealde? Tō hwon\textsuperscript{D2} rēcelēasodest ðū ðǣre giefe, ðe ic ðē geaf? Ic ðē nū āfierre\textsuperscript{D3} fram mīnre sęlene, ðe ic ðē forgeaf; ðonne bist\textsuperscript{D3} ðū wǣdla on woruldlīfe. For hwon\textsuperscript{D2} noldest\textsuperscript{D4} ðū geðęncean ðæt ic wille\textsuperscript{D5} forgieldan ǣghwilcum męnn āne gōde dǣd, ðe for mīnum naman mann gedēð? Mid hundtēontigum ic hit him forgielde,\textsuperscript{D3} swā hit is on mīnum godspelle gecweden and gesǣd,\textsuperscript{D6} ‘Swā-hwæt-swā\textsuperscript{D7} gē sęllað ānum of mīnum ðǣm lǣstum,\textsuperscript{D8} gē hit simle mē sęllað,\textsuperscript{D9} and ic ēow wið\textsuperscript{D10} ðǣm gesęlle\textsuperscript{D3} ēcne drēam\textsuperscript{D11} on heofonum.’}
% mdparagraph:P608:end

\end{minipage}\par
\reserve{5\baselineskip}
% mdparagraph:P609:begin
\textbf{Original notes — Cook.}
% mdparagraph:P609:end

% mdparagraph:P610:begin
\textsuperscript{D1}: From what OE. word is the Mod. Eng. rich derived (see Skeat, Prin., p. 61)? From what OE. meaning is the modern signification derived?
% mdparagraph:P610:end

% mdparagraph:P611:begin
\textsuperscript{D2}: Note Wulfstan’s use of tō hwon, for hwon, in the sense of why. See 88.
% mdparagraph:P611:end

% mdparagraph:P612:begin
\textsuperscript{D3}: See 188.
% mdparagraph:P612:end

% mdparagraph:P613:begin
\textsuperscript{D4}: See 139.
% mdparagraph:P613:end

% mdparagraph:P614:begin
\textsuperscript{D5}: Why should not the preterit be used here?
% mdparagraph:P614:end

% mdparagraph:P615:begin
\textsuperscript{D6}: Pleonastic.
% mdparagraph:P615:end

% mdparagraph:P616:begin
\textsuperscript{D7}: What portion of this is lost, and how is it replaced, in the Mod. Eng. whatsoever?
% mdparagraph:P616:end

% mdparagraph:P617:begin
\textsuperscript{D8}: See 66.
% mdparagraph:P617:end

% mdparagraph:P618:begin
\textsuperscript{D9}: Cf. the form of this sentence with that on p. 135, l. 14.
% mdparagraph:P618:end

% mdparagraph:P619:begin
\textsuperscript{D10}: = in return for. How is this to be reconciled with other senses of wið?
% mdparagraph:P619:end

% mdparagraph:P620:begin
\textsuperscript{D11}: Not dream, but joy, bliss.
% mdparagraph:P620:end

% mdparagraph:P621:begin
\textbf{Tasks D01–D05 — new teaching material.} Keep the modern vocabulary, construction helps and original notes available for every task.
% mdparagraph:P621:end

\reserve{5\baselineskip}
\par\begin{minipage}{\linewidth}\raggedright
% mdparagraph:P622:begin
\textbf{D01}
% mdparagraph:P622:end

% item: NEW-U15-D01
% mdparagraph:P623:begin
State the rich man’s fault, the consequence announced, and the promise made to a giver. Identify who speaks and to whom.
% mdparagraph:P623:end

% output: NEW-U15-D01-O001
% output: NEW-U15-D01-O002
% output: NEW-U15-D01-O003
% output: NEW-U15-D01-O004
% mdparagraph:P624:begin
\textbf{Responses, in order:} (1) fault; (2) consequence; (3) promise; (4) speaker and addressee.
% mdparagraph:P624:end

\end{minipage}\par\addvspace{4pt}

\par\begin{minipage}{\linewidth}\raggedright
% mdparagraph:P625:begin
\textbf{D02}
% mdparagraph:P625:end

% item: NEW-U15-D02
% mdparagraph:P626:begin
In the question beginning For hwon noldest ðū, explain the negative verb and the infinitive it governs. Explain why ic wille inside the following content clause can remain present rather than mechanically copying the question’s past tense.
% mdparagraph:P626:end

% output: NEW-U15-D02-O001
% output: NEW-U15-D02-O002
% mdparagraph:P627:begin
\textbf{Responses, in order:} (1) noldest and its infinitive; (2) tense of wille in this context.
% mdparagraph:P627:end

\end{minipage}\par\addvspace{4pt}

\par\begin{minipage}{\linewidth}\raggedright
% mdparagraph:P628:begin
\textbf{D03}
% mdparagraph:P628:end

% item: NEW-U15-D03
% mdparagraph:P629:begin
In the gospel saying beginning Swā-hwæt-swā, identify what is given, its immediate recipient, the recipient identified by the speaker, and the pronoun that resumes the thing given.
% mdparagraph:P629:end

% output: NEW-U15-D03-O001
% output: NEW-U15-D03-O002
% output: NEW-U15-D03-O003
% output: NEW-U15-D03-O004
% mdparagraph:P630:begin
\textbf{Responses, in order:} (1) thing given; (2) immediate recipient; (3) recipient identified by the speaker; (4) resuming pronoun.
% mdparagraph:P630:end

\end{minipage}\par\addvspace{4pt}

\par\begin{minipage}{\linewidth}\raggedright
% mdparagraph:P631:begin
\textbf{D04}
% mdparagraph:P631:end

% item: NEW-U15-D04
% mdparagraph:P632:begin
Identify the case relations in ānum of mīnum ðǣm lǣstum and wið ðǣm. Use the local helps; explain the difference between choosing a recipient from a group and receiving something in return.
% mdparagraph:P632:end

% output: NEW-U15-D04-O001
% output: NEW-U15-D04-O002
% mdparagraph:P633:begin
\textbf{Responses, in order:} (1) case relations in the first phrase; (2) case relation and meaning in wið ðǣm.
% mdparagraph:P633:end

\end{minipage}\par\addvspace{4pt}

\par\begin{minipage}{\linewidth}\raggedright
% mdparagraph:P634:begin
\textbf{D05}
% mdparagraph:P634:end

% item: NEW-U15-D05
% mdparagraph:P635:begin
Translate the complete assigned opening argument, from Se Hǣlend through heofonum. Preserve the change from rebuke to promise and the embedded gospel quotation.
% mdparagraph:P635:end

% output: NEW-U15-D05-O001
\end{minipage}\par\addvspace{4pt}

\reserve{7\baselineskip}
% mdparagraph:P636:begin
\subsection*{P. Alfred’s preface to Boethius}
% mdparagraph:P636:end

% mdparagraph:P637:begin
Cook, First Book, printed p. 162 (PDF p. 178). The complete preface.
% mdparagraph:P637:end

% mdparagraph:P638:begin
\textbf{Modern vocabulary — supplied for this reading.} Forms are glossed in their local senses; the parses identify the forms encountered, not complete new paradigms.
% mdparagraph:P638:end

% teaching: NEW-U15-LOCAL-HELPS
% teaching: NEW-U15-RECORD-AELFRED-U15-SOURCE-NOM-SG
% teaching: U15-READ-LEX-AELFRED
% support: U15-READ-LEX-AELFRED
% mdparagraph:P639:begin
\textbf{Ælfred} — Alfred (nominative masculine singular personal name).
% mdparagraph:P639:end

% teaching: NEW-U15-RECORD-AEMETTA-U15-SOURCE-DAT-SG
% teaching: NEW-U15-RECORD-AEMETTA-U15-SOURCE-SENSE
% teaching: U15-READ-LEX-AEMETTA
% support: U15-READ-LEX-AEMETTA
% mdparagraph:P640:begin
\textbf{ǣmettan} — leisure; available time (dative masculine singular after be).
% mdparagraph:P640:end

% teaching: NEW-U15-RECORD-ANDGIET-U15-SOURCE-ACC-SG
% teaching: NEW-U15-RECORD-ANDGIET-U15-SOURCE-DAT-SG
% teaching: NEW-U15-RECORD-ANDGIET-U15-SOURCE-GEN-SG
% teaching: NEW-U15-RECORD-ANDGIET-U15-SOURCE-SENSE
% teaching: NEW-U15-RECORD-ANDGIET-U15-SOURCE-UNDERSTANDING
% teaching: U15-READ-LEX-ANDGIET
% support: U15-READ-LEX-ANDGIET
% mdparagraph:P641:begin
\textbf{andgiet} — meaning; sense (accusative neuter singular object of sętte). \textbf{andgiete} — meaning; sense (dative neuter singular after of). \textbf{andgietes} — understanding (genitive neuter singular dependent on mǣðe).
% mdparagraph:P641:end

% teaching: NEW-U15-RECORD-ANDGIETFULLICE-U15-SOURCE-SENSE
% teaching: NEW-U15-RECORD-ANDGIETFULLICE-U15-SOURCE-SUPERL
% teaching: U15-READ-LEX-ANDGIETFULLICE
% support: U15-READ-LEX-ANDGIETFULLICE
% mdparagraph:P642:begin
\textbf{andgietfullicost} — most intelligibly (superlative adverb).
% mdparagraph:P642:end

% teaching: NEW-U06-SUPFORM-be-invariant-source
% teaching: NEW-U06-SUPSENSE-be-according-source
% teaching: NEW-U14-LOCAL-HELPS
% teaching: NEW-U15-RECORD-BE-U15-SOURCE-WORD-BY-WORD
% teaching: U15-READ-LEX-BE
% support: U15-READ-LEX-BE
% mdparagraph:P643:begin
\textbf{be} — word by word, literal correspondence in word be worde. \textbf{be} — according to.
% mdparagraph:P643:end

% teaching: NEW-U15-RECORD-BECUMAN-U15-SOURCE-PRET-3-PL
% teaching: NEW-U15-RECORD-BECUMAN-U15-SOURCE-SENSE
% teaching: U15-READ-LEX-BECUMAN
% support: U15-READ-LEX-BECUMAN
% mdparagraph:P644:begin
\textbf{becōmon} — come upon; befall (preterite indicative third plural).
% mdparagraph:P644:end

% teaching: NEW-U15-RECORD-BISGIAN-U15-SOURCE-PRET-3-PL
% teaching: NEW-U15-RECORD-BISGIAN-U15-SOURCE-SENSE
% teaching: U15-READ-LEX-BISGIAN
% support: U15-READ-LEX-BISGIAN
% mdparagraph:P645:begin
\textbf{bisgodon} — occupy; trouble (preterite indicative third plural; syncopated form).
% mdparagraph:P645:end

% teaching: NEW-U15-RECORD-BISGU-U15-SOURCE-NOM-PL
% teaching: NEW-U15-RECORD-BISGU-U15-SOURCE-SENSE
% teaching: U15-READ-LEX-BISGU
% support: U15-READ-LEX-BISGU
% mdparagraph:P646:begin
\textbf{bisga} — occupations; cares (nominative feminine plural subject).
% mdparagraph:P646:end

% teaching: NEW-U15-RECORD-BOCLAEDEN-U15-SOURCE-DAT-SG
% teaching: NEW-U15-RECORD-BOCLAEDEN-U15-SOURCE-SENSE
% teaching: U15-READ-LEX-BOCLAEDEN
% support: U15-READ-LEX-BOCLAEDEN
% mdparagraph:P647:begin
\textbf{Bōclædene} — book-Latin; Latin (dative neuter singular after of).
% mdparagraph:P647:end

% teaching: NEW-U15-RECORD-EARFOTHRIME-U15-SOURCE-NOM-F-PL
% teaching: NEW-U15-RECORD-EARFOTHRIME-U15-SOURCE-SENSE
% teaching: U15-READ-LEX-EARFOTHRIME
% support: U15-READ-LEX-EARFOTHRIME
% mdparagraph:P648:begin
\textbf{earfoðrīme} — hard to count; innumerable (predicate strong nominative feminine plural).
% mdparagraph:P648:end

% teaching: NEW-U15-RECORD-EFT-U15-SOURCE-INVARIANT
% teaching: NEW-U15-RECORD-EFT-U15-SOURCE-SENSE
% teaching: U15-READ-LEX-EFT
% support: U15-READ-LEX-EFT
% mdparagraph:P649:begin
\textbf{eft} — again; afterwards.
% mdparagraph:P649:end

% teaching: NEW-U15-LOCAL-HELPS
% teaching: NEW-U15-RECORD-ENGLISC-U15-SOURCE-DAT-N-SG
% teaching: U15-READ-LEX-ENGLISC
% support: U15-READ-LEX-ENGLISC
% mdparagraph:P650:begin
\textbf{Engliscum} — English (strong dative neuter singular).
% mdparagraph:P650:end

% teaching: NEW-U15-RECORD-ENGLISC-LANGUAGE-U15-SOURCE-ACC-SG
% teaching: NEW-U15-RECORD-ENGLISC-LANGUAGE-U15-SOURCE-SENSE
% teaching: U15-READ-LEX-ENGLISC-LANGUAGE
% support: U15-READ-LEX-ENGLISC-LANGUAGE
% mdparagraph:P651:begin
\textbf{Englisc} — English, the language (accusative language name after on, goal of translation).
% mdparagraph:P651:end

% teaching: NEW-U11-LOCAL-HELPS
% teaching: NEW-U15-RECORD-FOR-U15-SOURCE-CAUSE
% teaching: U15-READ-LEX-FOR
% support: U15-READ-LEX-FOR
% mdparagraph:P652:begin
\textbf{for} — because of, with dative of hindering circumstance. \textbf{for} — for; for the sake of, selected accusative/\allowbreak{}dative expressions supplied locally.
% mdparagraph:P652:end

% teaching: NEW-U15-RECORD-FOR-THAM-THE-U15-SOURCE-INVARIANT
% teaching: NEW-U15-RECORD-FOR-THAM-THE-U15-SOURCE-SENSE
% teaching: U15-READ-LEX-FOR-THAM-THE
% support: U15-READ-LEX-FOR-THAM-THE
% mdparagraph:P653:begin
\textbf{for-ðām-ðe} — because (causal conjunction).
% mdparagraph:P653:end

% teaching: NEW-U05-SUPFORM-gedon-past-participle
% teaching: NEW-U15-RECORD-GEDON-U15-SOURCE-ARRANGEMENT
% teaching: U15-READ-LEX-GEDON
% support: U15-READ-LEX-GEDON
% mdparagraph:P654:begin
\textbf{gedōn} — made; arranged, describing the present form of the translation (past participle).
% mdparagraph:P654:end

% teaching: NEW-U15-RECORD-GELEORNIAN-U15-SOURCE-PAST-PART
% teaching: NEW-U15-RECORD-GELEORNIAN-U15-SOURCE-SENSE
% teaching: U15-READ-LEX-GELEORNIAN
% support: U15-READ-LEX-GELEORNIAN
% mdparagraph:P655:begin
\textbf{geleornod} — learn; study (past participle in hæfde geleornod).
% mdparagraph:P655:end

% teaching: NEW-U15-RECORD-GERECCAN-U15-SOURCE-INF
% teaching: NEW-U15-RECORD-GERECCAN-U15-SOURCE-SENSE
% teaching: U15-READ-LEX-GERECCAN
% support: U15-READ-LEX-GERECCAN
% mdparagraph:P656:begin
\textbf{geręccan} — explain; render (plain infinitive after meahte).
% mdparagraph:P656:end

% teaching: NEW-U14-PAR-GEWENDAN
% teaching: NEW-U15-RECORD-GEWENDAN-U15-SOURCE-TRANSLATE
% teaching: U15-READ-LEX-GEWENDAN
% support: U15-READ-LEX-GEWENDAN
% mdparagraph:P657:begin
\textbf{gewęnd} — translate; turn from one language into another (past participle).
% mdparagraph:P657:end

% teaching: NEW-U15-LOCAL-HELPS
% teaching: NEW-U15-RECORD-LAEDEN-U15-SOURCE-DAT-SG
% teaching: U15-READ-LEX-LAEDEN
% support: U15-READ-LEX-LAEDEN
% mdparagraph:P658:begin
\textbf{Lædene} — Latin, the language (dative neuter singular language name after of).
% mdparagraph:P658:end

% teaching: NEW-U15-RECORD-LYSTAN-U15-SOURCE-PRES-SUBJ-3-SG
% teaching: NEW-U15-RECORD-LYSTAN-U15-SOURCE-SENSE
% teaching: U15-READ-LEX-LYSTAN
% support: U15-READ-LEX-LYSTAN
% mdparagraph:P659:begin
\textbf{lyste} — please; be a desire to (present subjunctive third singular impersonal; understood experiencer accusative from ǣlcne).
% mdparagraph:P659:end

% teaching: NEW-U15-RECORD-MAETH-U15-SOURCE-DAT-SG
% teaching: NEW-U15-RECORD-MAETH-U15-SOURCE-SENSE
% teaching: U15-READ-LEX-MAETH
% support: U15-READ-LEX-MAETH
% mdparagraph:P660:begin
\textbf{mǣðe} — measure; capacity (dative feminine singular after be).
% mdparagraph:P660:end

% teaching: NEW-U15-RECORD-MANIGFEALD-U15-SOURCE-DAT-F-PL-STRONG
% teaching: NEW-U15-RECORD-MANIGFEALD-U15-SOURCE-SENSE
% teaching: U15-READ-LEX-MANIGFEALD
% support: U15-READ-LEX-MANIGFEALD
% mdparagraph:P661:begin
\textbf{manigfealdum} — manifold; many (strong dative feminine plural).
% mdparagraph:P661:end

% teaching: NEW-U15-RECORD-MISLIC-U15-SOURCE-DAT-F-PL-STRONG
% teaching: NEW-U15-RECORD-MISLIC-U15-SOURCE-SENSE
% teaching: U15-READ-LEX-MISLIC
% support: U15-READ-LEX-MISLIC
% mdparagraph:P662:begin
\textbf{mislicum} — various (strong dative feminine plural).
% mdparagraph:P662:end

% teaching: NEW-U13-SOURCE-HELPS
% teaching: U15-READ-LEX-MOD
% support: U15-READ-LEX-MOD
% mdparagraph:P663:begin
\textbf{mōde} — mind; attention (dative neuter singular).
% mdparagraph:P663:end

% teaching: NEW-U15-RECORD-OF-U15-SOURCE-CORRESPONDENCE
% teaching: U15-READ-LEX-OF
% support: U15-READ-LEX-OF
% mdparagraph:P664:begin
\textbf{of} — from; corresponding to, sense for sense.
% mdparagraph:P664:end

% teaching: NEW-U15-RECORD-ON-U15-SOURCE-DURING
% teaching: U15-READ-LEX-ON
% support: U15-READ-LEX-ON
% mdparagraph:P665:begin
\textbf{on} — during, with dative time.
% mdparagraph:P665:end

% teaching: NEW-U15-RECORD-ONGIETAN-U15-SOURCE-PRES-SUBJ-3-SG
% teaching: NEW-U15-VOC-ONGIETAN
% teaching: U15-READ-LEX-ONGIETAN
% support: U15-READ-LEX-ONGIETAN
% mdparagraph:P666:begin
\textbf{ongiete} — understand; perceive (present subjunctive third singular in condition).
% mdparagraph:P666:end

% teaching: NEW-U09-VOC-raedan
% teaching: U15-READ-LEX-RAEDAN
% support: U15-READ-LEX-RAEDAN
% mdparagraph:P667:begin
\textbf{rǣdan} — read (plain infinitive).
% mdparagraph:P667:end

% teaching: NEW-U15-RECORD-RIHTLICE-U15-SOURCE-COMPARATIVE
% teaching: NEW-U15-RECORD-RIHTLICE-U15-SOURCE-SENSE
% teaching: U15-READ-LEX-RIHTLICE
% support: U15-READ-LEX-RIHTLICE
% mdparagraph:P668:begin
\textbf{rihtlicor} — more correctly (comparative adverb).
% mdparagraph:P668:end

% teaching: NEW-U15-RECORD-SPELL-U15-SOURCE-DAT-SG
% teaching: NEW-U15-RECORD-SPELL-U15-SOURCE-SENSE
% teaching: U15-READ-LEX-SPELL
% support: U15-READ-LEX-SPELL
% mdparagraph:P669:begin
\textbf{spelle} — account; discourse (dative neuter singular after tō).
% mdparagraph:P669:end

% teaching: NEW-U09-SOURCE-HELPS
% teaching: NEW-U09-VOC-swa-swa
% teaching: NEW-U10-SOURCE-HELPS
% teaching: NEW-U12-REUSED-U12-SRC-HELP-15
% teaching: NEW-U12-SOURCE-HELPS
% teaching: NEW-U13-PAR-SWA-SWA
% teaching: NEW-U13-VOC-SWA-SWA
% teaching: NEW-U14-LOCAL-HELPS
% teaching: U15-READ-LEX-SWA-SWA
% support: U15-READ-LEX-SWA-SWA
% mdparagraph:P670:begin
\textbf{swā-swā} — just as; like.
% mdparagraph:P670:end

% teaching: NEW-U15-RECORD-SWEOTOLE-U15-SOURCE-SENSE
% teaching: NEW-U15-RECORD-SWEOTOLE-U15-SOURCE-SUPERL
% teaching: U15-READ-LEX-SWEOTOLE
% support: U15-READ-LEX-SWEOTOLE
% mdparagraph:P671:begin
\textbf{sweotolost} — most clearly (superlative adverb).
% mdparagraph:P671:end

% teaching: NEW-U15-RECORD-THAET-THAET-U15-SOURCE-ACC-N-SG
% teaching: NEW-U15-RECORD-THAET-THAET-U15-SOURCE-SENSE
% teaching: U15-READ-LEX-THAET-THAET
% support: U15-READ-LEX-THAET-THAET
% mdparagraph:P672:begin
\textbf{ðæt-ðæt} — that which; what (demonstrative accusative neuter + relative accusative neuter).
% mdparagraph:P672:end

% teaching: NEW-U15-RECORD-THEAH-U15-SOURCE-NEVERTHELESS
% teaching: U15-READ-LEX-THEAH
% support: U15-READ-LEX-THEAH
% mdparagraph:P673:begin
\textbf{ðēah} — nevertheless; yet (independent adverb).
% mdparagraph:P673:end

% teaching: NEW-U15-RECORD-TO-U15-SOURCE-TRANSLATION-RESULT
% teaching: U15-READ-LEX-TO
% support: U15-READ-LEX-TO
% mdparagraph:P674:begin
\textbf{tō} — into; to, result in a new linguistic/\allowbreak{}literary form.
% mdparagraph:P674:end

% teaching: NEW-U15-RECORD-UNDERFON-U15-SOURCE-PAST-PART
% teaching: NEW-U15-RECORD-UNDERFON-U15-SOURCE-SENSE
% teaching: U15-READ-LEX-UNDERFON
% support: U15-READ-LEX-UNDERFON
% mdparagraph:P675:begin
\textbf{underfangen} — receive; take charge of (past participle in anterior construction hæfde underfangen).
% mdparagraph:P675:end

% teaching: NEW-U15-RECORD-WEALHSTOD-U15-SOURCE-NOM-SG
% teaching: NEW-U15-RECORD-WEALHSTOD-U15-SOURCE-SENSE
% teaching: U15-READ-LEX-WEALHSTOD
% support: U15-READ-LEX-WEALHSTOD
% mdparagraph:P676:begin
\textbf{wealhstōd} — interpreter; translator (predicate nominative masculine singular).
% mdparagraph:P676:end

% teaching: NEW-U15-RECORD-WENDAN-U15-SOURCE-PRET-3-SG
% teaching: NEW-U15-RECORD-WENDAN-U15-SOURCE-SENSE
% teaching: U15-READ-LEX-WENDAN
% support: U15-READ-LEX-WENDAN
% mdparagraph:P677:begin
\textbf{węnde} — translate (preterite indicative third singular).
% mdparagraph:P677:end

% teaching: NEW-U15-RECORD-WESAN-U15-SOURCE-SENSE
% teaching: NEW-U15-RECORD-WESAN-U15-SOURCE-SIND-PRES-3-PL
% teaching: U15-READ-LEX-WESAN
% support: U15-READ-LEX-WESAN
% mdparagraph:P678:begin
\textbf{sind} — are (present indicative third plural).
% mdparagraph:P678:end

% teaching: NEW-U15-RECORD-WITAN-BLAME-U15-SOURCE-PRES-SUBJ-3-SG
% teaching: NEW-U15-RECORD-WITAN-BLAME-U15-SOURCE-SENSE
% teaching: U15-READ-LEX-WITAN-BLAME
% support: U15-READ-LEX-WITAN-BLAME
% mdparagraph:P679:begin
\textbf{wīte} — blame; censure (present subjunctive third singular, governed request coordinated with gebidde).
% mdparagraph:P679:end

% teaching: NEW-U15-RECORD-WORULDBISGU-U15-SOURCE-DAT-PL
% teaching: NEW-U15-RECORD-WORULDBISGU-U15-SOURCE-SENSE
% teaching: U15-READ-LEX-WORULDBISGU
% support: U15-READ-LEX-WORULDBISGU
% mdparagraph:P680:begin
\textbf{woruldbisgum} — worldly occupations; cares (dative feminine plural after for).
% mdparagraph:P680:end

\reserve{5\baselineskip}
% mdparagraph:P681:begin
\textbf{Modern construction help.}
% mdparagraph:P681:end

% teaching: NEW-U15-RECORD-U15-SOURCE-P-TRANSLATION
% teaching: U15-READ-P-TRANSLATION
% support: U15-READ-P-TRANSLATION
% mdparagraph:P682:begin
Ælfred is subject. The feminine hīe/\allowbreak{}hēo refers to bōc. The source-language phrases use of plus dative; on Englisc gives the translation’s goal. Word be worde contrasts literal verbal correspondence with andgiet of andgiete, sense for sense. Gedōn describes the work’s resulting arrangement.
% mdparagraph:P682:end

% teaching: NEW-U15-RECORD-U15-SOURCE-P-HINDRANCES
% teaching: U15-READ-P-HINDRANCES
% support: U15-READ-P-HINDRANCES
% mdparagraph:P683:begin
The two superlative adverbs mean “most clearly and most intelligibly”. For with dative introduces the worldly cares that constrained how clearly Alfred could explain it. Ǣgðer ge … ge means “both … and”; on mōde/\allowbreak{}on līchaman identify mind and body. Sind means “are”. The relative after Ðā bisga refers to the cares, which are the subject of becōmon.
% mdparagraph:P683:end

% teaching: NEW-U15-RECORD-U15-SOURCE-P-ANTERIOR-ACTIONS
% teaching: U15-READ-P-ANTERIOR-ACTIONS
% support: U15-READ-P-ANTERIOR-ACTIONS
% mdparagraph:P684:begin
Hæfde underfangen = “had received/\allowbreak{}taken charge of.” The repeated hæfde governs both geleornod and gewęnd: “had studied ... and translated ...”. Ðēah here is “nevertheless”; ðā ... ðā frames “when ... then.” Tō Engliscum spelle and tō lēoðe mark resulting literary form. Hīe/\allowbreak{}hēo still refers to the book.
% mdparagraph:P684:end

% teaching: NEW-U15-RECORD-U15-SOURCE-P-APPEAL
% teaching: U15-READ-P-APPEAL
% support: U15-READ-P-APPEAL
% mdparagraph:P685:begin
Supply hē (Alfred) with bitt and hālsað, as Cook’s note 4 directs. Ǣlcne ðāra ðe ... means “each of those who ...”. Lyste is impersonal: “whom it may please to read this book”; its understood accusative experiencer is the prospective reader. The requested ðæt clause contains subjunctive gebidde, “that he pray for him.”
% mdparagraph:P685:end

% teaching: NEW-U15-RECORD-U15-SOURCE-P-READER-AUTHOR-REFERENCE
% teaching: U15-READ-P-READER-AUTHOR-REFERENCE
% support: U15-READ-P-READER-AUTHOR-REFERENCE
% mdparagraph:P686:begin
Ne wīte continues the request: “and not blame him.” In gif hē hit rihtlicor ongiete ðonne hē meahte, the first hē is the reader and the second is Alfred (Cook’s notes 6–7). Ongiete is subjunctive; wīte is the verb “blame,” not the noun “punishment.”
% mdparagraph:P686:end

% teaching: NEW-U15-RECORD-U15-SOURCE-P-CAPACITY
% teaching: U15-READ-P-CAPACITY
% support: U15-READ-P-CAPACITY
% mdparagraph:P687:begin
Be ... mǣðe and be ... ǣmettan mean according to the measure of one’s understanding and one’s available leisure. Sceal governs sprecan and dōn. Both ðæt ... and ðæt-ðæt ... function as “what/\allowbreak{}that which.” Preserve the source’s hyphenated spelling in the latter.
% mdparagraph:P687:end

\reserve{5\baselineskip}
% mdparagraph:P688:begin
\textbf{Cook’s text}
% mdparagraph:P688:end

% mdparagraph:P689:begin
X. ALFRED’S PREFACE TO BOETHIUS.
% mdparagraph:P689:end

% mdparagraph:P690:begin
(Prefixed to his translation.)
% mdparagraph:P690:end

% teaching: NEW-U15-READ-FB-X-ALFRED-PREFACE
\par\begin{minipage}{\linewidth}\raggedright
% passage: FB-X-ALFRED-PREFACE
% mdparagraph:P691:begin
\readingtext{Ælfred cyning wæs wealhstōd\textsuperscript{P1} ðisse bēc, and hīe of Bōclædene\textsuperscript{P2} on Englisc węnde, swā hēo nū is gedōn. Hwīlum hē sętte word be worde, hwīlum andgiet of andgiete, swā-swā hē hit ðā sweotolost and andgietfullicost geręccan meahte for ðǣm mislicum and manigfealdum woruldbisgum ðe hine oft ǣgðer ge on mōde ge on līchaman bisgodon. Ðā bisga\textsuperscript{P3} ūs sind swīðe earfoðrīme ðe on his dagum on ðā rīcu becōmon ðe hē underfangen hæfde, and ðēah, ðā hē ðās bōc hæfde geleornod, and of Lædene tō Engliscum spelle gewęnd, ðā geworhte hē hīe eft tō lēoðe, swā-swā hēo nū gedōn is. And nū bitt\textsuperscript{P4} and for Godes naman hālsað ǣlcne ðāra ðe ðās bōc rǣdan lyste,\textsuperscript{P5} ðæt hē for hine gebidde, and him ne wīte gif hē\textsuperscript{P6} hit rihtlicor ongiete ðonne hē\textsuperscript{P7} meahte; for-ðām-ðe ǣlc mann sceal be his andgietes mǣðe, and be his ǣmettan, sprecan ðæt hē spricð, and dōn ðæt-ðæt hē dēð.}
% mdparagraph:P691:end

\end{minipage}\par
\reserve{5\baselineskip}
% mdparagraph:P692:begin
\textbf{Original notes — Cook.}
% mdparagraph:P692:end

% mdparagraph:P693:begin
\textsuperscript{P1}: Wealh- signifies foreign (see walnut), and sometimes servant, orig. Celtic, Celt (cf. Wales, Welsh, Cornwall), from Volcæ, the name of a Celtic tribe (Cæsar, Gallic War, Bk. VII.).
% mdparagraph:P693:end

% mdparagraph:P694:begin
\textsuperscript{P2}: Perhaps originally in contrast to the Latin spoken in Britain.
% mdparagraph:P694:end

% mdparagraph:P695:begin
\textsuperscript{P3}: See 51. a.
% mdparagraph:P695:end

% mdparagraph:P696:begin
\textsuperscript{P4}: Supply hē.
% mdparagraph:P696:end

% mdparagraph:P697:begin
\textsuperscript{P5}: See 190.
% mdparagraph:P697:end

% mdparagraph:P698:begin
\textsuperscript{P6}: The reader.
% mdparagraph:P698:end

% mdparagraph:P699:begin
\textsuperscript{P7}: Alfred.
% mdparagraph:P699:end

% mdparagraph:P700:begin
\textbf{Tasks P01–P05 — new teaching material.} Keep the modern vocabulary, construction helps and original notes available for every task.
% mdparagraph:P700:end

\reserve{5\baselineskip}
\par\begin{minipage}{\linewidth}\raggedright
% mdparagraph:P701:begin
\textbf{P01}
% mdparagraph:P701:end

% item: NEW-U15-P01
% mdparagraph:P702:begin
State the two translation methods described, the obstacles Alfred faced, the later change he made to the translation, and his request to readers.
% mdparagraph:P702:end

% output: NEW-U15-P01-O001
% output: NEW-U15-P01-O002
% output: NEW-U15-P01-O003
% output: NEW-U15-P01-O004
% output: NEW-U15-P01-O005
% mdparagraph:P703:begin
\textbf{Responses, in order:} (1) first translation method; (2) second translation method; (3) obstacles; (4) later change; (5) request to readers.
% mdparagraph:P703:end

\end{minipage}\par\addvspace{4pt}

\par\begin{minipage}{\linewidth}\raggedright
% mdparagraph:P704:begin
\textbf{P02}
% mdparagraph:P704:end

% item: NEW-U15-P02
% mdparagraph:P705:begin
In the appeal from ðæt hē for hine gebidde through ðonne hē meahte, identify the referent of each of these five tokens in order: hē, hine, him, hē after gif, hē after ðonne. Explain the resulting contrast between the two comparison subjects.
% mdparagraph:P705:end

% output: NEW-U15-P02-O001
% output: NEW-U15-P02-O002
% output: NEW-U15-P02-O003
% output: NEW-U15-P02-O004
% output: NEW-U15-P02-O005
% output: NEW-U15-P02-O006
% mdparagraph:P706:begin
\textbf{Responses, in order:} (1) hē before gebidde; (2) hine after for; (3) him before ne wīte; (4) hē after gif; (5) hē after ðonne; (6) contrast between the comparison subjects.
% mdparagraph:P706:end

\end{minipage}\par\addvspace{4pt}

\par\begin{minipage}{\linewidth}\raggedright
% mdparagraph:P707:begin
\textbf{P03}
% mdparagraph:P707:end

% item: NEW-U15-P03
% mdparagraph:P708:begin
Explain ǣgðer ge on mōde ge on līchaman. Then trace the subject and finite verb of the sentence beginning Ðā bisga and identify the two relative clauses that follow.
% mdparagraph:P708:end

% output: NEW-U15-P03-O001
% output: NEW-U15-P03-O002
% output: NEW-U15-P03-O003
% output: NEW-U15-P03-O004
% mdparagraph:P709:begin
\textbf{Responses, in order:} (1) meaning and coordination in the first expression; (2) main subject and finite verb; (3) first relative clause; (4) second relative clause.
% mdparagraph:P709:end

\end{minipage}\par\addvspace{4pt}

\par\begin{minipage}{\linewidth}\raggedright
% mdparagraph:P710:begin
\textbf{P04}
% mdparagraph:P710:end

% item: NEW-U15-P04
% mdparagraph:P711:begin
Explain what is being compared by rihtlicor … ðonne and what the second clause leaves understood. State the function of ne in and him ne wīte and the relationship of that clause to the preceding request.
% mdparagraph:P711:end

% output: NEW-U15-P04-O001
% output: NEW-U15-P04-O002
% mdparagraph:P712:begin
\textbf{Responses, in order:} (1) comparison and understood words; (2) function of ne and relation of its clause to the request.
% mdparagraph:P712:end

\end{minipage}\par\addvspace{4pt}

\par\begin{minipage}{\linewidth}\raggedright
% mdparagraph:P713:begin
\textbf{P05}
% mdparagraph:P713:end

% item: NEW-U15-P05
% mdparagraph:P714:begin
Translate Alfred’s complete preface. Preserve the distinction between his account of the translation and his request for a charitable reading.
% mdparagraph:P714:end

% output: NEW-U15-P05-O001
\end{minipage}\par\addvspace{4pt}

\reserve{7\baselineskip}
% mdparagraph:P715:begin
\subsection*{R. A prayer of King Alfred}
% mdparagraph:P715:end

% mdparagraph:P716:begin
Cook, First Book, printed p. 163 (PDF p. 179). The complete prayer.
% mdparagraph:P716:end

% mdparagraph:P717:begin
\textbf{Modern vocabulary — supplied for this reading.} Forms are glossed in their local senses; the parses identify the forms encountered, not complete new paradigms.
% mdparagraph:P717:end

% teaching: NEW-U15-RECORD-A-EVER-U15-SOURCE-INVARIANT
% teaching: NEW-U15-RECORD-A-EVER-U15-SOURCE-SENSE
% teaching: U15-READ-LEX-A-EVER
% support: U15-READ-LEX-A-EVER
% mdparagraph:P718:begin
\textbf{ā} — ever; always (adverb repeated for intensification).
% mdparagraph:P718:end

% teaching: NEW-U15-RECORD-AEGHWILC-U15-SOURCE-DAT-M-SG
% teaching: NEW-U15-RECORD-AEGHWILC-U15-SOURCE-SENSE
% teaching: U15-READ-LEX-AEGHWILC
% support: U15-READ-LEX-AEGHWILC
% mdparagraph:P719:begin
\textbf{ǣghwilcum} — each; every (strong dative masculine singular).
% mdparagraph:P719:end

% teaching: NEW-U15-RECORD-AELMIHTIG-U15-SOURCE-SENSE
% teaching: NEW-U15-RECORD-AELMIHTIG-U15-SOURCE-VOC-M-SG-WEAK
% teaching: U15-READ-LEX-AELMIHTIG
% support: U15-READ-LEX-AELMIHTIG
% mdparagraph:P720:begin
\textbf{ælmihtiga} — almighty (weak masculine singular in direct address).
% mdparagraph:P720:end

% teaching: NEW-U15-RECORD-AFIERRAN-U15-SOURCE-IMP-2-SG
% teaching: NEW-U15-RECORD-AFIERRAN-U15-SOURCE-SENSE
% teaching: U15-READ-LEX-AFIERRAN
% support: U15-READ-LEX-AFIERRAN
% mdparagraph:P721:begin
\textbf{āfierr} — remove; take away (imperative second singular).
% mdparagraph:P721:end

% teaching: NEW-U15-RECORD-AMEN-U15-SOURCE-INVARIANT
% teaching: NEW-U15-RECORD-AMEN-U15-SOURCE-SENSE
% teaching: U15-READ-LEX-AMEN
% support: U15-READ-LEX-AMEN
% mdparagraph:P722:begin
\textbf{Amen} — Amen (liturgical closing word).
% mdparagraph:P722:end

% teaching: NEW-U15-RECORD-AWYRCEAN-U15-SOURCE-PRET-1-SG
% teaching: NEW-U15-RECORD-AWYRCEAN-U15-SOURCE-SENSE
% teaching: U15-READ-LEX-AWYRCEAN
% support: U15-READ-LEX-AWYRCEAN
% mdparagraph:P723:begin
\textbf{āworhte} — do; perform, here deserve by one's conduct (preterite indicative first singular with tō ðē).
% mdparagraph:P723:end

% teaching: NEW-U15-RECORD-CLAENE-U15-SOURCE-DAT-M-SG-STRONG
% teaching: NEW-U15-RECORD-CLAENE-U15-SOURCE-DAT-N-SG-STRONG
% teaching: NEW-U15-RECORD-CLAENE-U15-SOURCE-SENSE
% teaching: U15-READ-LEX-CLAENE
% support: U15-READ-LEX-CLAENE
% mdparagraph:P724:begin
\textbf{clǣnum} — pure; clean (strong dative neuter singular). \textbf{clǣnum} — pure; clean (strong dative masculine singular).
% mdparagraph:P724:end

% teaching: NEW-U15-RECORD-EARNUNG-U15-SOURCE-DAT-PL
% teaching: NEW-U15-RECORD-EARNUNG-U15-SOURCE-SENSE
% teaching: U15-READ-LEX-EARNUNG
% support: U15-READ-LEX-EARNUNG
% mdparagraph:P725:begin
\textbf{earnungum} — merits; good services (dative feminine plural coordinated under for).
% mdparagraph:P725:end

% teaching: NEW-U15-RECORD-ENDE-U15-SOURCE-DAT-SG
% teaching: NEW-U15-RECORD-ENDE-U15-SOURCE-SENSE
% teaching: U15-READ-LEX-ENDE
% support: U15-READ-LEX-ENDE
% mdparagraph:P726:begin
\textbf{ęnde} — end (dative masculine singular after būtan).
% mdparagraph:P726:end

% teaching: NEW-U15-RECORD-FOR-U15-SOURCE-APPEAL
% teaching: U15-READ-LEX-FOR
% support: U15-READ-LEX-FOR
% mdparagraph:P727:begin
\textbf{for} — by; for the sake of, naming grounds of appeal in the prayer.
% mdparagraph:P727:end

% teaching: NEW-U15-RECORD-FUL-U15-SOURCE-ACC-F-SG-WEAK
% teaching: NEW-U15-RECORD-FUL-U15-SOURCE-SENSE
% teaching: U15-READ-LEX-FUL
% support: U15-READ-LEX-FUL
% mdparagraph:P728:begin
\textbf{fūlan} — foul (weak accusative feminine singular).
% mdparagraph:P728:end

% teaching: NEW-U15-RECORD-FULTUM-U15-SOURCE-NOM-SG
% teaching: NEW-U15-RECORD-FULTUM-U15-SOURCE-SENSE
% teaching: U15-READ-LEX-FULTUM
% support: U15-READ-LEX-FULTUM
% mdparagraph:P729:begin
\textbf{Fultum} — help; support (predicate nominative masculine singular, address to God).
% mdparagraph:P729:end

% teaching: NEW-U15-RECORD-GALNES-U15-SOURCE-ACC-SG
% teaching: NEW-U15-RECORD-GALNES-U15-SOURCE-SENSE
% teaching: U15-READ-LEX-GALNES
% support: U15-READ-LEX-GALNES
% mdparagraph:P730:begin
\textbf{gālnesse} — lust (accusative feminine singular).
% mdparagraph:P730:end

% teaching: NEW-U15-RECORD-GESCIELDAN-U15-SOURCE-IMP-2-SG
% teaching: NEW-U15-RECORD-GESCIELDAN-U15-SOURCE-SENSE
% teaching: U15-READ-LEX-GESCIELDAN
% support: U15-READ-LEX-GESCIELDAN
% mdparagraph:P731:begin
\textbf{gescield} — shield; protect (imperative second singular).
% mdparagraph:P731:end

% teaching: NEW-U15-RECORD-GESEWENLIC-U15-SOURCE-DAT-M-PL-STRONG
% teaching: NEW-U15-RECORD-GESEWENLIC-U15-SOURCE-SENSE
% teaching: U15-READ-LEX-GESEWENLIC
% support: U15-READ-LEX-GESEWENLIC
% mdparagraph:P732:begin
\textbf{gesewenlicum} — visible (strong dative masculine plural).
% mdparagraph:P732:end

% teaching: NEW-U15-RECORD-GESTATHELIAN-U15-SOURCE-IMP-2-SG
% teaching: NEW-U15-RECORD-GESTATHELIAN-U15-SOURCE-SENSE
% teaching: U15-READ-LEX-GESTATHELIAN
% support: U15-READ-LEX-GESTATHELIAN
% mdparagraph:P733:begin
\textbf{gestaðela} — establish; make steadfast (imperative second singular).
% mdparagraph:P733:end

% teaching: NEW-U15-RECORD-GESTRANGIAN-U15-SOURCE-IMP-2-SG
% teaching: NEW-U15-RECORD-GESTRANGIAN-U15-SOURCE-SENSE
% teaching: U15-READ-LEX-GESTRANGIAN
% support: U15-READ-LEX-GESTRANGIAN
% mdparagraph:P734:begin
\textbf{gestranga} — strengthen (imperative second singular).
% mdparagraph:P734:end

% teaching: NEW-U15-RECORD-GETHANC-U15-SOURCE-DAT-SG
% teaching: NEW-U15-RECORD-GETHANC-U15-SOURCE-SENSE
% teaching: U15-READ-LEX-GETHANC
% support: U15-READ-LEX-GETHANC
% mdparagraph:P735:begin
\textbf{geðance} — thought; mind (dative neuter singular after mid).
% mdparagraph:P735:end

% teaching: NEW-U15-RECORD-GEWISSIAN-U15-SOURCE-IMP-2-SG
% teaching: NEW-U15-RECORD-GEWISSIAN-U15-SOURCE-PRES-SUBJ-2-SG
% teaching: NEW-U15-RECORD-GEWISSIAN-U15-SOURCE-SENSE
% teaching: U15-READ-LEX-GEWISSIAN
% support: U15-READ-LEX-GEWISSIAN
% mdparagraph:P736:begin
\textbf{gewissie} — guide; direct (present subjunctive second singular in requested ðæt clause). \textbf{gewissa} — guide; direct (imperative second singular).
% mdparagraph:P736:end

% teaching: NEW-U15-RECORD-HALIG-U15-SOURCE-GEN-PL-SUBSTANTIVAL
% teaching: NEW-U15-RECORD-HALIG-U15-SOURCE-SENSE
% teaching: U15-READ-LEX-HALIG
% support: U15-READ-LEX-HALIG
% mdparagraph:P737:begin
\textbf{hālgena} — saints; holy ones (substantival genitive plural in -ena).
% mdparagraph:P737:end

% teaching: NEW-U15-RECORD-HE-U15-SOURCE-HIERA-GEN-PL
% teaching: NEW-U15-RECORD-HE-U15-SOURCE-SENSE
% teaching: U15-READ-LEX-HE
% support: U15-READ-LEX-HE
% mdparagraph:P738:begin
\textbf{hīera} — their (third-person genitive plural).
% mdparagraph:P738:end

% teaching: NEW-U15-RECORD-INWEARDLICE-U15-SOURCE-INVARIANT
% teaching: NEW-U15-RECORD-INWEARDLICE-U15-SOURCE-SENSE
% teaching: U15-READ-LEX-INWEARDLICE
% support: U15-READ-LEX-INWEARDLICE
% mdparagraph:P739:begin
\textbf{inweardlīce} — fervently; wholeheartedly.
% mdparagraph:P739:end

% teaching: NEW-U15-RECORD-LOF-U15-SOURCE-NOM-SG
% teaching: NEW-U15-RECORD-LOF-U15-SOURCE-SENSE
% teaching: U15-READ-LEX-LOF
% support: U15-READ-LEX-LOF
% mdparagraph:P740:begin
\textbf{lof} — praise (nominative masculine singular subject).
% mdparagraph:P740:end

% teaching: NEW-U15-RECORD-LUFE-U15-SOURCE-DAT-SG
% teaching: NEW-U15-RECORD-LUFE-U15-SOURCE-SENSE
% teaching: U15-READ-LEX-LUFE
% support: U15-READ-LEX-LUFE
% mdparagraph:P741:begin
\textbf{lufan} — love (weak feminine dative singular after for).
% mdparagraph:P741:end

% teaching: NEW-U15-RECORD-MAEGTHHAD-U15-SOURCE-DAT-SG
% teaching: NEW-U15-RECORD-MAEGTHHAD-U15-SOURCE-SENSE
% teaching: U15-READ-LEX-MAEGTHHAD
% support: U15-READ-LEX-MAEGTHHAD
% mdparagraph:P742:begin
\textbf{mægðhāde} — maidenhood; virginity (dative masculine singular after for).
% mdparagraph:P742:end

% teaching: NEW-U15-RECORD-MARIA-NAME-U15-SOURCE-GEN-CONTEXT
% teaching: NEW-U15-RECORD-MARIA-NAME-U15-SOURCE-SENSE
% teaching: U15-READ-LEX-MARIA-NAME
% support: U15-READ-LEX-MARIA-NAME
% mdparagraph:P743:begin
\textbf{Marian} — Mary (genitive form of a foreign name).
% mdparagraph:P743:end

% teaching: NEW-U15-RECORD-MICHAEL-NAME-U15-SOURCE-GEN-CONTEXT
% teaching: NEW-U15-RECORD-MICHAEL-NAME-U15-SOURCE-SENSE
% teaching: U15-READ-LEX-MICHAEL-NAME
% support: U15-READ-LEX-MICHAEL-NAME
% mdparagraph:P744:begin
\textbf{Michaeles} — Michael (genitive personal name, vernacular -es with Latin Sancti).
% mdparagraph:P744:end

% teaching: NEW-U13-VOC-MID
% teaching: U15-READ-LEX-MID
% support: U15-READ-LEX-MID
% mdparagraph:P745:begin
\textbf{mid} — with; in the manner of the accompanying quality or mental disposition.
% mdparagraph:P745:end

% teaching: NEW-U13-SOURCE-HELPS
% teaching: U15-READ-LEX-MOD
% support: U15-READ-LEX-MOD
% mdparagraph:P746:begin
\textbf{mōd} — mind; attention (accusative neuter singular).
% mdparagraph:P746:end

% teaching: NEW-U15-RECORD-SANCTUS-LATIN-U15-SOURCE-GEN-F-SG
% teaching: NEW-U15-RECORD-SANCTUS-LATIN-U15-SOURCE-GEN-M-SG
% teaching: NEW-U15-RECORD-SANCTUS-LATIN-U15-SOURCE-SENSE
% teaching: U15-READ-LEX-SANCTUS-LATIN
% support: U15-READ-LEX-SANCTUS-LATIN
% mdparagraph:P747:begin
\textbf{Sanctæ} — Saint; holy (Latin genitive feminine singular accompanying Marian). \textbf{Sancti} — Saint; holy (Latin genitive masculine singular).
% mdparagraph:P747:end

% teaching: NEW-U15-RECORD-SCIEPPEND-U15-SOURCE-NOM-SG
% teaching: NEW-U15-RECORD-SCIEPPEND-U15-SOURCE-SENSE
% teaching: U15-READ-LEX-SCIEPPEND
% support: U15-READ-LEX-SCIEPPEND
% mdparagraph:P748:begin
\textbf{Scieppend} — Creator (predicate nominative masculine singular title).
% mdparagraph:P748:end

% teaching: NEW-U15-RECORD-TAECAN-U15-SOURCE-IMP-2-SG
% teaching: NEW-U15-RECORD-TAECAN-U15-SOURCE-SENSE
% teaching: U15-READ-LEX-TAECAN
% support: U15-READ-LEX-TAECAN
% mdparagraph:P749:begin
\textbf{tǣc} — teach; show (imperative second singular).
% mdparagraph:P749:end

% teaching: NEW-U15-RECORD-THEARF-U15-SOURCE-DAT-SG
% teaching: NEW-U15-RECORD-THEARF-U15-SOURCE-SENSE
% teaching: U15-READ-LEX-THEARF
% support: U15-READ-LEX-THEARF
% mdparagraph:P750:begin
\textbf{ðearfe} — benefit; good; need (dative feminine singular after tō, with dependent sāwle).
% mdparagraph:P750:end

% teaching: NEW-U15-RECORD-TO-U15-SOURCE-BENEFIT
% teaching: NEW-U15-RECORD-TO-U15-SOURCE-CONFORMITY
% teaching: NEW-U15-RECORD-TO-U15-SOURCE-DIRECTED-CONDUCT
% teaching: NEW-U15-RECORD-TO-U15-SOURCE-ETERNITY
% teaching: U15-READ-LEX-TO
% support: U15-READ-LEX-TO
% mdparagraph:P751:begin
\textbf{tō} — towards; in relation to the person to whom conduct is directed. \textbf{tō} — to; in accordance with, a will as the goal of guidance. \textbf{tō} — for; to the benefit of. \textbf{tō} — for ever, in fixed tō worulde.
% mdparagraph:P751:end

% teaching: NEW-U15-RECORD-TOFORAN-U15-SOURCE-INVARIANT
% teaching: NEW-U15-RECORD-TOFORAN-U15-SOURCE-SENSE
% teaching: U15-READ-LEX-TOFORAN
% support: U15-READ-LEX-TOFORAN
% mdparagraph:P752:begin
\textbf{tōforan} — before; above all (preposition with dative of comparison or preference).
% mdparagraph:P752:end

% teaching: NEW-U15-RECORD-TOHOPA-U15-SOURCE-NOM-SG
% teaching: NEW-U15-RECORD-TOHOPA-U15-SOURCE-SENSE
% teaching: U15-READ-LEX-TOHOPA
% support: U15-READ-LEX-TOHOPA
% mdparagraph:P753:begin
\textbf{Tōhopa} — hope (predicate nominative masculine singular).
% mdparagraph:P753:end

% teaching: NEW-U15-RECORD-TREOWNES-U15-SOURCE-NOM-SG
% teaching: NEW-U15-RECORD-TREOWNES-U15-SOURCE-SENSE
% teaching: U15-READ-LEX-TREOWNES
% support: U15-READ-LEX-TREOWNES
% mdparagraph:P754:begin
\textbf{Trēownes} — trust; confidence (predicate nominative feminine singular).
% mdparagraph:P754:end

% teaching: NEW-U15-RECORD-UNGESEWENLIC-U15-SOURCE-DAT-M-PL-STRONG
% teaching: NEW-U15-RECORD-UNGESEWENLIC-U15-SOURCE-SENSE
% teaching: U15-READ-LEX-UNGESEWENLIC
% support: U15-READ-LEX-UNGESEWENLIC
% mdparagraph:P755:begin
\textbf{ungesewenlicum} — invisible (strong dative masculine plural).
% mdparagraph:P755:end

% teaching: NEW-U15-RECORD-WEALDEND-U15-SOURCE-SENSE
% teaching: NEW-U15-RECORD-WEALDEND-U15-SOURCE-VOC-SG
% teaching: U15-READ-LEX-WEALDEND
% support: U15-READ-LEX-WEALDEND
% mdparagraph:P756:begin
\textbf{Wealdend} — Ruler (nominative/\allowbreak{}vocative masculine singular title).
% mdparagraph:P756:end

% teaching: NEW-U15-RECORD-WITHERWINNA-U15-SOURCE-DAT-PL
% teaching: NEW-U15-RECORD-WITHERWINNA-U15-SOURCE-SENSE
% teaching: U15-READ-LEX-WITHERWINNA
% support: U15-READ-LEX-WITHERWINNA
% mdparagraph:P757:begin
\textbf{wiðerwinnum} — adversaries (dative masculine plural after wið).
% mdparagraph:P757:end

% teaching: NEW-U15-RECORD-WYRHTA-U15-SOURCE-MAKER
% teaching: U15-READ-LEX-WYRHTA
% support: U15-READ-LEX-WYRHTA
% mdparagraph:P758:begin
\textbf{wyrhta} — Maker; Creator (not the previously taught workman sense) (nominative singular masculine).
% mdparagraph:P758:end

\reserve{5\baselineskip}
% mdparagraph:P759:begin
\textbf{Modern construction help.}
% mdparagraph:P759:end

% teaching: NEW-U15-RECORD-U15-SOURCE-R-ADDRESS
% teaching: U15-READ-R-ADDRESS
% support: U15-READ-R-ADDRESS
% mdparagraph:P760:begin
Dryhten and God are direct address; ælmihtiga is weak singular. Wyrhta and Wealdend are “Maker and Ruler,” with genitive ealra gesceafta. This Maker sense is distinct from the previously taught workman sense.
% mdparagraph:P760:end

% teaching: NEW-U15-RECORD-U15-SOURCE-R-GROUNDS-OF-APPEAL
% teaching: U15-READ-R-GROUNDS-OF-APPEAL
% support: U15-READ-R-GROUNDS-OF-APPEAL
% mdparagraph:P761:begin
The for phrases name grounds of appeal. In ðǣre hālgan rōde tācne, the genitive phrase “of the holy cross” depends on dative tācne. Sanctæ Marian and Sancti Michaeles are supplied names in genitive function with mixed foreign-name inflection; do not derive a new productive Latin paradigm. Gehīersumnesse is obedience. Hālgena means saints; lufan is the weak noun’s dative; earnungum means merits.
% mdparagraph:P761:end

% teaching: NEW-U15-RECORD-U15-SOURCE-R-GUIDANCE
% teaching: U15-READ-R-GUIDANCE
% support: U15-READ-R-GUIDANCE
% mdparagraph:P762:begin
The requested ðæt clause has gewissie, “that you guide.” Bęt ðonne ic āworhte tō ðē means better than my conduct towards you has deserved. The following gewissa is a direct singular command. In mīnre sāwle ðearfe, mīnre sāwle is genitive, dependent on dative ðearfe: “for my soul’s benefit.” Cunne is first-singular subjunctive, “than I myself may know.”
% mdparagraph:P762:end

% teaching: NEW-U15-RECORD-U15-SOURCE-R-COMMANDS
% teaching: U15-READ-R-COMMANDS
% support: U15-READ-R-COMMANDS
% mdparagraph:P763:begin
Gestaðela, gestranga, āfierr and gescield are addressed singular commands. Establish my mind according to your will; strengthen me against the devil’s temptations; remove lust and unrighteousness; protect me against visible and invisible adversaries. The latter adjectives agree with dative plural wiðerwinnum.
% mdparagraph:P763:end

% teaching: NEW-U15-RECORD-U15-SOURCE-R-TEACHING-PURPOSE
% teaching: U15-READ-R-TEACHING-PURPOSE
% support: U15-READ-R-TEACHING-PURPOSE
% mdparagraph:P764:begin
Tǣc takes dative mē. Its following ðīnne willan is the accusative object of wyrceanne (“to do your will”), as Cook’s note 7 specifies. The ðæt clause gives purpose with mæge. Tōforan eallum ðingum means above all things; the two mid phrases give manner: with pure thought and a pure body.
% mdparagraph:P764:end

% teaching: NEW-U15-RECORD-U15-SOURCE-R-TITLES-AND-DOXOLOGY
% teaching: U15-READ-R-TITLES-AND-DOXOLOGY
% support: U15-READ-R-TITLES-AND-DOXOLOGY
% mdparagraph:P765:begin
The final causal clause names God as Creator, Redeemer, help, comfort, trust and hope. Sīe is an optative wish, “may there be”; ðē is dative. Ā is repeated three times, followed by tō worulde būtan ǣghwilcum ęnde, “for ever without any end.”
% mdparagraph:P765:end

\reserve{5\baselineskip}
\par\begin{minipage}{\linewidth}\raggedright
% mdparagraph:P766:begin
\textbf{Cook’s text}
% mdparagraph:P766:end

% mdparagraph:P767:begin
XI. A PRAYER OF KING ALFRED.
% mdparagraph:P767:end

% mdparagraph:P768:begin
(From the end of his translation of Boethius.)
% mdparagraph:P768:end

% teaching: NEW-U15-READ-FB-XI-ALFRED-PRAYER

% passage: FB-XI-ALFRED-PRAYER
% mdparagraph:P769:begin
\readingtext{Dryhten,\textsuperscript{R1} ælmihtiga God,\textsuperscript{R1} Wyrhta and Wealdend ealra gesceafta, ic bidde ðē for ðīnre miclan mildheortnesse, and for ðǣre hālgan rōde tācne,\textsuperscript{R2} and for Sanctæ Marian mægðhāde, and for Sancti Michaeles gehīersumnesse, and for ealra ðīnra hālgena\textsuperscript{R3} lufan and hīera earnungum, ðæt ðū mē gewissie\textsuperscript{R4} bęt ðonne ic āworhte tō ðē; and gewissa mē tō ðīnum willan, and tō mīnre sāwle ðearfe,\textsuperscript{R5} bęt ðonne ic self cunne\textsuperscript{R6}; and gestaðela mīn mōd tō ðīnum willan and tō mīnre sāwle ðearfe; and gestranga mē wið ðæs dēofles\textsuperscript{R3} costnungum; and āfierr fram mē ðā fūlan gālnesse and ǣlce unrihtwīsnesse; and gescield mē wið mīnum wiðerwinnum, gesewenlicum and ungesewenlicum; and tǣc mē ðīnne willan\textsuperscript{R7} tō wyrceanne; ðæt ic mæge\textsuperscript{R8} ðē inweardlīce lufian tōforan eallum ðingum, mid clǣnum geðance and mid clǣnum līchaman. For-ðon-ðe ðū eart mīn Scieppend,\textsuperscript{R9} and mīn Ālīesend, mīn Fultum, mīn Frōfor, mīn Trēownes, and mīn Tōhopa. Sīe ðē lof and wuldor nū and ā ā ā, tō worulde būtan ǣghwilcum ęnde. Amen.}
% mdparagraph:P769:end

\end{minipage}\par
\reserve{5\baselineskip}
% mdparagraph:P770:begin
\textbf{Original notes — Cook.}
% mdparagraph:P770:end

% mdparagraph:P771:begin
\textsuperscript{R1}: See 152.
% mdparagraph:P771:end

% mdparagraph:P772:begin
\textsuperscript{R2}: Governed by for.
% mdparagraph:P772:end

% mdparagraph:P773:begin
\textsuperscript{R3}: See 153. c.
% mdparagraph:P773:end

% mdparagraph:P774:begin
\textsuperscript{R4}: See 194. b.
% mdparagraph:P774:end

% mdparagraph:P775:begin
\textsuperscript{R5}: See 166.
% mdparagraph:P775:end

% mdparagraph:P776:begin
\textsuperscript{R6}: Optative (130).
% mdparagraph:P776:end

% mdparagraph:P777:begin
\textsuperscript{R7}: Object of wyrceanne.
% mdparagraph:P777:end

% mdparagraph:P778:begin
\textsuperscript{R8}: See 196. d.
% mdparagraph:P778:end

% mdparagraph:P779:begin
\textsuperscript{R9}: See 150.
% mdparagraph:P779:end

\par\begin{minipage}{\linewidth}\raggedright
% mdparagraph:P780:begin
\textbf{Tasks R01–R05 — new teaching material.} Keep the modern vocabulary, construction helps and original notes available for every task.
% mdparagraph:P780:end


% mdparagraph:P781:begin
\textbf{R01}
% mdparagraph:P781:end

% item: NEW-U15-R01
% mdparagraph:P782:begin
Identify the opening requested action after ic bidde and each of the six subsequent direct commands. For each give the actual finite form, its mood and its essential meaning.
% mdparagraph:P782:end

% output: NEW-U15-R01-O001
% output: NEW-U15-R01-O002
% output: NEW-U15-R01-O003
% output: NEW-U15-R01-O004
% output: NEW-U15-R01-O005
% output: NEW-U15-R01-O006
% output: NEW-U15-R01-O007
% output: NEW-U15-R01-O008
% output: NEW-U15-R01-O009
% output: NEW-U15-R01-O010
% output: NEW-U15-R01-O011
% output: NEW-U15-R01-O012
% output: NEW-U15-R01-O013
% output: NEW-U15-R01-O014
% output: NEW-U15-R01-O015
% output: NEW-U15-R01-O016
% output: NEW-U15-R01-O017
% output: NEW-U15-R01-O018
% output: NEW-U15-R01-O019
% output: NEW-U15-R01-O020
% output: NEW-U15-R01-O021
% mdtable:T83:begin
\par\begingroup
\fontsize{10.5}{13.2}\selectfont
\setlength{\LTleft}{0pt}\setlength{\LTright}{\fill}
\noindent
\begin{tabular}{@{}>{\raggedright\arraybackslash}p{39.78mm}>{\raggedright\arraybackslash}p{25.26mm}>{\raggedright\arraybackslash}p{25.26mm}>{\raggedright\arraybackslash}p{25.26mm}@{}}
\toprule
\textbf{Requested action} & \textbf{Finite form} & \textbf{Mood} & \textbf{Meaning} \\
\midrule
Opening requested action & (1) \_\_\_\_ & (2) \_\_\_\_ & (3) \_\_\_\_ \\
First direct command & (4) \_\_\_\_ & (5) \_\_\_\_ & (6) \_\_\_\_ \\
Second direct command & (7) \_\_\_\_ & (8) \_\_\_\_ & (9) \_\_\_\_ \\
Third direct command & (10) \_\_\_\_ & (11) \_\_\_\_ & (12) \_\_\_\_ \\
Fourth direct command & (13) \_\_\_\_ & (14) \_\_\_\_ & (15) \_\_\_\_ \\
Fifth direct command & (16) \_\_\_\_ & (17) \_\_\_\_ & (18) \_\_\_\_ \\
Sixth direct command & (19) \_\_\_\_ & (20) \_\_\_\_ & (21) \_\_\_\_ \\
\bottomrule
\end{tabular}
\par\endgroup

% mdtable:T83:end
\end{minipage}\par\addvspace{4pt}

\par\begin{minipage}{\linewidth}\raggedright
% mdparagraph:P783:begin
\textbf{R02}
% mdparagraph:P783:end

% item: NEW-U15-R02
% mdparagraph:P784:begin
Explain the grounds of entreaty introduced by the repeated for phrases. In for ðǣre hālgan rōde tācne, distinguish the noun governed by for from its dependent genitive.
% mdparagraph:P784:end

% output: NEW-U15-R02-O001
% output: NEW-U15-R02-O002
% mdparagraph:P785:begin
\textbf{Responses, in order:} (1) meaning of the for phrases; (2) case relations in ðǣre hālgan rōde tācne.
% mdparagraph:P785:end

\end{minipage}\par\addvspace{4pt}

\par\begin{minipage}{\linewidth}\raggedright
% mdparagraph:P786:begin
\textbf{R03}
% mdparagraph:P786:end

% item: NEW-U15-R03
% mdparagraph:P787:begin
Compare the two bęt ðonne clauses. State the standard in each and explain why the second one uses cunne in this petition.
% mdparagraph:P787:end

% output: NEW-U15-R03-O001
% output: NEW-U15-R03-O002
% output: NEW-U15-R03-O003
% mdparagraph:P788:begin
\textbf{Responses, in order:} (1) first comparison standard; (2) second comparison standard; (3) explanation of cunne.
% mdparagraph:P788:end

\end{minipage}\par\addvspace{4pt}

\par\begin{minipage}{\linewidth}\raggedright
% mdparagraph:P789:begin
\textbf{R04}
% mdparagraph:P789:end

% item: NEW-U15-R04
% mdparagraph:P790:begin
Distinguish the purpose clause ðæt ic mæge … lufian from the concluding Sīe ðē lof and wuldor. Identify the actor/\allowbreak{}addressee, the mood and the function of each.
% mdparagraph:P790:end

% output: NEW-U15-R04-O001
% output: NEW-U15-R04-O002
% output: NEW-U15-R04-O003
% output: NEW-U15-R04-O004
% output: NEW-U15-R04-O005
% output: NEW-U15-R04-O006
% mdparagraph:P791:begin
\textbf{Responses, in order:} (1) first clause: actor/\allowbreak{}addressee; (2) first clause: mood; (3) first clause: function; (4) second clause: actor/\allowbreak{}addressee; (5) second clause: mood; (6) second clause: function.
% mdparagraph:P791:end

\end{minipage}\par\addvspace{4pt}

\par\begin{minipage}{\linewidth}\raggedright
% mdparagraph:P792:begin
\textbf{R05}
% mdparagraph:P792:end

% item: NEW-U15-R05
% mdparagraph:P793:begin
Translate the complete prayer, including the grounds of entreaty, the commands, the purpose and the concluding praise.
% mdparagraph:P793:end

% output: NEW-U15-R05-O001
\end{minipage}\par\addvspace{4pt}

\reserve{8\baselineskip}
% mdparagraph:P794:begin
\section*{9. Sources and selected forms}
% mdparagraph:P794:end

\begingroup\fontsize{10}{12.8}\selectfont
\setlength{\parskip}{1.5pt}
% mdparagraph:P795:begin
\textsuperscript{1}: Cook, A First Book in Old English, 2nd ed., pp. 28–37, especially §§43–53. Tables give selected complete models; stated alternatives remain valid.
% mdparagraph:P795:end

% mdparagraph:P796:begin
\textsuperscript{2}: Wright, Old English Grammar, p. 184, §§382–383, supplies the complete strengu class and its variants. Bosworth–Toller, p. 500b, treats hǣlu as indeclinable. These facts support the selected singular tables without assigning plural count uses.
% mdparagraph:P796:end

% mdparagraph:P797:begin
\textsuperscript{3}: Cook, First Book, pp. 23, 74, 78–79 and 100. The selected negative indicative forms combine the positive systems with the stated contractions; the multiword entries remain ordinary ne plus a verb.
% mdparagraph:P797:end

% mdparagraph:P798:begin
\textsuperscript{4}: Cook, First Book, pp. 52–53; Wright, Old English Grammar, pp. 239–240. Further indefinite and free-choice families are acknowledged only within the limits stated above.
% mdparagraph:P798:end

% mdparagraph:P799:begin
\textsuperscript{5}: Cook, First Book, pp. 54–74 and the cited lexical entries. The class patterns determine endings and stem distribution; each newly required system is printed in full above.
% mdparagraph:P799:end

% mdparagraph:P800:begin
\textsuperscript{6}: Cook, First Book, pp. 39–41, §§57–59. The selected -ig syncopation and double-consonant spellings do not exclude the stated historical alternatives.
% mdparagraph:P800:end

% mdparagraph:P801:begin
\textsuperscript{7}: Cook, Exercises in Old English (1895), IX.28–30 and X, XI.31, printed pp. 15–18 (PDF pp. 23–26). The preserved source prompts are distinct from the modern directions and helps.
% mdparagraph:P801:end

\endgroup
\end{document}
